\documentclass[11pt,letterpaper]{article}
\usepackage[margin=1in]{geometry}
\usepackage{amsmath,amssymb,amsthm}
\usepackage{enumitem}
\usepackage{booktabs}
\usepackage{array}
\usepackage[table]{xcolor}
\usepackage[most]{tcolorbox}
\usepackage{hyperref}
\hypersetup{colorlinks=true,linkcolor=blue!60!black,urlcolor=blue!60!black}

% ---------- Study-primer environments (shared by every section) ----------
% Numbered worked example, exam style
\newcounter{example}[section]
\renewcommand{\theexample}{\thesection.\arabic{example}}
\newtcolorbox[use counter=example]{example}[1][]{%
  enhanced,breakable,colback=blue!3,colframe=blue!50!black,
  title={Example~\thesection.\arabic{example}\ifstrempty{#1}{}{: #1}},fonttitle=\bfseries}

% Solution to the immediately preceding example
\newtcolorbox{solution}{%
  enhanced,breakable,colback=white,colframe=gray!60,
  title={Solution},fonttitle=\bfseries,coltitle=black,colbacktitle=gray!15}

% Key fact / formula to memorise
\newtcolorbox{keyfact}[1][]{%
  enhanced,breakable,colback=green!4,colframe=green!40!black,
  title={Key fact\ifstrempty{#1}{}{: #1}},fonttitle=\bfseries}

% Common exam trap
\newtcolorbox{trap}[1][]{%
  enhanced,breakable,colback=red!4,colframe=red!50!black,
  title={Trap\ifstrempty{#1}{}{: #1}},fonttitle=\bfseries}

% Practice problems (answers collected at the end of each section)
\newcounter{practice}[section]
\renewcommand{\thepractice}{\thesection.\arabic{practice}}
\newtcolorbox[use counter=practice]{practice}{%
  enhanced,breakable,colback=yellow!5,colframe=orange!70!black,
  title={Practice~\thesection.\arabic{practice}},fonttitle=\bfseries}

\newcommand{\Prob}{\mathbb{P}}
\newcommand{\E}{\mathbb{E}}
\newcommand{\Var}{\operatorname{Var}}
\DeclareMathOperator{\comp}{c}
\newcommand{\cc}[1]{#1^{\mathrm{c}}}   % complement

\title{SOA Exam P Study Primer\\[4pt]\large Topic 1: General Probability (23--30\% of the exam)}
\author{}
\date{September 2026 syllabus}

\begin{document}
\maketitle

\section*{How to use this primer}
Topic 1 has seven learning outcomes, (a) through (g). Each section below covers one outcome:
what the outcome asks, the definitions and formulas needed, worked examples in the style of the
exam's five-choice multiple-choice questions, the traps that cost points, and practice problems
with answers. The exam is 30 questions in three hours, so a typical question must be solved in
about six minutes. Every section ends with a one-page summary of what to memorise.
\tableofcontents
\clearpage

\section{Outcome (a): Sets, sample spaces, events, and the axioms of probability}

\subsection{What the outcome asks}

The exam does not ask you to recite definitions. It tests whether you can convert
a word problem into set notation, fill a two-event or three-event Venn diagram from
partial information, and apply the axioms and their direct consequences (complement
rule, inclusion-exclusion, monotonicity) to produce a number. The most common
question gives four to seven proportions of policyholders holding various coverages
and asks for the proportion in one region of the diagram or in a union of regions
(``none'', ``exactly one'', ``at most one'', ``auto but not home''). A second type
gives two probabilities with no other information and asks for the largest or
smallest value a third probability can take; that is solved with the axioms alone.
The arithmetic is short; points are lost in the translation from English to sets.

\subsection{Definitions}

\paragraph{Sample space, outcomes, events.}
The \emph{sample space} $S$ of a random experiment (select a policyholder at
random, observe next year's claims) is the set of all possible results; each element of $S$ is
an \emph{outcome}. An \emph{event} is a subset of $S$. $A$ \emph{occurs} when the observed outcome lies in $A$. $S$ is the certain
event and $\emptyset$ the impossible event.

\paragraph{Set operations.}
For events $A, B \subseteq S$:
\begin{itemize}[nosep]
  \item Union $A \cup B$: outcomes in $A$ or $B$ or both. English: ``$A$ or $B$'',
        ``at least one of $A$, $B$''.
  \item Intersection $A \cap B$ (also written $AB$): outcomes in both. English:
        ``$A$ and $B$'', ``both''.
  \item Complement $\cc{A} = S \setminus A$: outcomes not in $A$. English: ``not $A$''.
  \item Difference $A \setminus B = A \cap \cc{B}$: outcomes in $A$ but not in $B$.
        English: ``$A$ but not $B$'', ``$A$ only'' (when only two events are in play).
  \item $A$ and $B$ are \emph{disjoint} (mutually exclusive) when $A \cap B = \emptyset$;
        $A \subseteq B$ when every outcome of $A$ is in $B$, so $A$ occurring forces $B$.
\end{itemize}
Union and intersection distribute over each other:
$A \cap (B \cup C) = (A\cap B) \cup (A \cap C)$. Any event $A$ splits as the
disjoint union $A = (A \cap B) \cup (A \cap \cc{B})$; this identity is used constantly.

\begin{keyfact}[De Morgan's laws]
\[
  \cc{(A \cup B)} = \cc{A} \cap \cc{B},
  \qquad
  \cc{(A \cap B)} = \cc{A} \cup \cc{B}.
\]
In words: ``neither $A$ nor $B$'' is the complement of ``at least one'';
``not both'' is the complement of ``both''. The same pattern holds for any number
of events: the complement of a union is the intersection of the complements, and
the complement of an intersection is the union of the complements.
\end{keyfact}

\paragraph{Set functions and the collection of events.}
A \emph{set function} is a function whose inputs are sets. The number of elements
$|A|$, the length of an interval, and probability are all set functions.
The domain of a probability is a collection $\mathcal{F}$ of subsets of $S$, the
\emph{events}. $\mathcal{F}$ must contain $S$, must contain $\cc{A}$ whenever it
contains $A$, and must contain $A_1 \cup A_2 \cup \cdots$ whenever it contains
each $A_i$; a collection with these closure properties is a \emph{sigma-field}
(or sigma-algebra). When $S$ is finite, $\mathcal{F}$ is every subset of $S$; on the exam, every set
you can describe is an event.

\paragraph{Probability as a set function.}
A \emph{probability} is a set function $\Prob : \mathcal{F} \to \mathbb{R}$
satisfying three axioms (Kolmogorov, 1933).

\begin{keyfact}[The three axioms of probability]
\begin{enumerate}[nosep,label=\textbf{A\arabic*.}]
  \item Nonnegativity: $\Prob(A) \ge 0$ for every event $A$.
  \item Normalisation: $\Prob(S) = 1$.
  \item Countable additivity: if $A_1, A_2, \ldots$ are pairwise disjoint
        ($A_i \cap A_j = \emptyset$ for $i \ne j$), then
        \[
          \Prob\bigl(A_1 \cup A_2 \cup \cdots\bigr) = \Prob(A_1) + \Prob(A_2) + \cdots .
        \]
\end{enumerate}
\end{keyfact}

Everything else in this outcome is a consequence of the three axioms.

\begin{keyfact}[Consequences of the axioms]
\begin{enumerate}[nosep,label=\textbf{C\arabic*.}]
  \item $\Prob(\emptyset) = 0$.
  \item Finite additivity: for pairwise disjoint $A_1, \ldots, A_n$,
        $\Prob(A_1 \cup \cdots \cup A_n) = \sum_{i=1}^{n} \Prob(A_i)$.
  \item Complement rule: $\Prob(\cc{A}) = 1 - \Prob(A)$.
  \item Monotonicity: if $A \subseteq B$ then $\Prob(A) \le \Prob(B)$ and
        $\Prob(B \setminus A) = \Prob(B) - \Prob(A)$.
  \item $0 \le \Prob(A) \le 1$.
  \item Difference rule: $\Prob(A \cap \cc{B}) = \Prob(A) - \Prob(A \cap B)$.
  \item Inclusion-exclusion, two events:
        \[
          \Prob(A \cup B) = \Prob(A) + \Prob(B) - \Prob(A \cap B).
        \]
  \item Inclusion-exclusion, three events:
        \begin{align*}
          \Prob(A \cup B \cup C) = {}& \Prob(A) + \Prob(B) + \Prob(C) \\
          & - \Prob(A\cap B) - \Prob(A \cap C) - \Prob(B \cap C) \\
          & + \Prob(A \cap B \cap C).
        \end{align*}
  \item Boole's inequality: $\Prob(A_1 \cup \cdots \cup A_n) \le \sum_{i=1}^{n}\Prob(A_i)$.
  \item Bonferroni's inequality: $\Prob(A \cap B) \ge \Prob(A) + \Prob(B) - 1$.
\end{enumerate}
\end{keyfact}

Each consequence comes from writing a set as a disjoint union and applying A3:
$S = A \cup \cc{A}$ gives C3; $B = A \cup (B \setminus A)$ gives C4 and C5;
$A = (A\cap B) \cup (A\cap\cc{B})$ gives C6; $A \cup B = A \cup (B \cap \cc{A})$
gives C7, and C7 applied twice gives C8. When the exam asks which statements
``must be true'', the answer is one of C1--C10 or an axiom.

\paragraph{Identities from adding two inclusion-exclusion equations.}
Some questions give probabilities of unions that involve complements, such as
$\Prob(A \cup B)$ and $\Prob(A \cup \cc{B})$, and ask for $\Prob(A)$. Write
inclusion-exclusion for each union and add. Since $\Prob(B) + \Prob(\cc{B}) = 1$ and
$\Prob(A\cap B) + \Prob(A \cap \cc{B}) = \Prob(A)$ (C6), the sum collapses.

\begin{keyfact}[Union-with-complement identities]
\begin{align*}
  \Prob(A \cup B) + \Prob(A \cup \cc{B}) &= 1 + \Prob(A), \\
  \Prob(A \cup B) + \Prob(\cc{A} \cup B) &= 1 + \Prob(B), \\
  \Prob(A \cup \cc{B}) &= 1 - \Prob(\cc{A} \cap B), \\
  \Prob(\cc{A} \cup \cc{B}) &= 1 - \Prob(A \cap B), \\
  \Prob(A \cup B) + \Prob(\cc{A} \cup \cc{B}) &= 1 + \Prob(\text{exactly one of } A, B), \\
  \Prob(A \cup \cc{B}) + \Prob(\cc{A} \cup B) &= 2 - \Prob(\text{exactly one of } A, B).
\end{align*}
\end{keyfact}

The third and fourth lines are De Morgan plus the complement rule: $A \cup \cc{B}$
is the complement of $\cc{A} \cap B$, so its probability is one minus the
``$B$ but not $A$'' region. Every identity in the box is also a statement about the
$2\times 2$ table: $A \cup \cc{B}$ is three of the four cells (all but
$\cc{A}\cap B$), so $\Prob(A \cup B) + \Prob(A \cup \cc{B})$ counts the two $A$ cells
twice and the two $\cc{A}$ cells once, which is $\Prob(A) + 1$. When an identity is not
remembered, fill the table with unknowns and add cells.

\subsection{Venn diagram method}

A Venn diagram partitions $S$ into disjoint regions, so the probability of any
event is the sum of the regions it contains (C2). Name the events, fill the regions
from the innermost intersection outward, then read off the requested set as a sum
of regions.

\paragraph{Two events (four regions).}
The four regions are $A\cap B$, $A \cap \cc{B}$, $\cc{A} \cap B$, and
$\cc{A} \cap \cc{B}$. A $2\times 2$ table is the fastest bookkeeping; row and column
totals are the marginal probabilities, and the grand total is $1$.
\begin{center}
\begin{tabular}{@{}lccc@{}}
\toprule
 & $B$ & $\cc{B}$ & Total \\
\midrule
$A$      & $\Prob(A \cap B)$        & $\Prob(A \cap \cc{B})$        & $\Prob(A)$ \\
$\cc{A}$ & $\Prob(\cc{A} \cap B)$   & $\Prob(\cc{A} \cap \cc{B})$   & $\Prob(\cc{A})$ \\
\midrule
Total    & $\Prob(B)$               & $\Prob(\cc{B})$               & $1$ \\
\bottomrule
\end{tabular}
\end{center}
Any three independent pieces of information determine the whole table.

\paragraph{Three events (eight regions).}
Fill in the order shown. Each region's probability is obtained by subtracting
regions already known from a given total.
\begin{center}
\begin{tabular}{@{}clll@{}}
\toprule
Order & Region & Description & How to compute \\
\midrule
1 & $A \cap B \cap C$                   & all three          & given, or found last from the total \\
2 & $A \cap B \cap \cc{C}$              & $A$ and $B$ only   & $\Prob(A\cap B) - \Prob(A\cap B\cap C)$ \\
3 & $A \cap \cc{B} \cap C$              & $A$ and $C$ only   & $\Prob(A\cap C) - \Prob(A\cap B\cap C)$ \\
4 & $\cc{A} \cap B \cap C$              & $B$ and $C$ only   & $\Prob(B\cap C) - \Prob(A\cap B\cap C)$ \\
5 & $A \cap \cc{B} \cap \cc{C}$         & $A$ only           & $\Prob(A) - (\text{regions } 1,2,3)$ \\
6 & $\cc{A} \cap B \cap \cc{C}$         & $B$ only           & $\Prob(B) - (\text{regions } 1,2,4)$ \\
7 & $\cc{A} \cap \cc{B} \cap C$         & $C$ only           & $\Prob(C) - (\text{regions } 1,3,4)$ \\
8 & $\cc{A} \cap \cc{B} \cap \cc{C}$    & none               & $1 - (\text{regions } 1\text{--}7)$ \\
\bottomrule
\end{tabular}
\end{center}
When the problem gives ``none'' instead of ``all three'', the same table is solved
in reverse: the sum of regions 1--7 equals $1 - \Prob(\text{none})$, which with
inclusion-exclusion determines $\Prob(A \cap B \cap C)$.

\paragraph{Binary classifications.}
A problem that classifies each policyholder as young or old, male or female, and
married or single is a three-event problem in disguise: pick one label from each
pair as the event ($Y$, $M$, $R$ for young, male, married) and the other label is its
complement. ``Young, female, single'' is the region $Y \cap \cc{M} \cap \cc{R}$,
which is ``$Y$ only'' in the eight-region table. Counts work exactly like
probabilities; the grand total replaces $1$.

\paragraph{A disjoint pair inside a three-event problem.}
When the problem states that two of the three events are mutually exclusive, say
$B \cap C = \emptyset$, two regions vanish: ``$B$ and $C$ only'' and ``all three''.
Six regions remain: $A$ only, $B$ only, $C$ only, $A \cap B$, $A \cap C$, and none.
Name them with letters, translate every sentence of the problem into a linear
equation in those letters, and solve. Useful translations:
\begin{itemize}[nosep]
  \item ``exactly two of the three'' is the sum $(A\cap B) + (A \cap C)$;
  \item ``exactly one'' is $(A \text{ only}) + (B \text{ only}) + (C \text{ only})$;
  \item ``$B$ but not $A$'' is $B$ only (because $B \cap C$ is empty);
  \item ``not $A$'' is $(B \text{ only}) + (C \text{ only}) + (\text{none})$;
  \item ``twice as many have $B$ as have $C$'' is $(B\text{ only}) + (A \cap B) = 2\bigl[(C \text{ only}) + (A \cap C)\bigr]$.
\end{itemize}
The total of the six regions is $1$ (or the population size). Five independent
facts plus the total determine all six regions; the exam usually gives exactly
enough, and the order of solving is: everything outside $A$ first, then the
``exactly two'' total, then the ratio.

\paragraph{A fourth set disjoint from the other three.}
A survey may add a fourth category $D$ with the statement that no member of $D$
belongs to $A$, $B$, or $C$. Then $D$ is disjoint from $A \cup B \cup C$, and
by finite additivity
\[
  \Prob(A \cup B \cup C \cup D) = \Prob(A \cup B \cup C) + \Prob(D),
\]
so ``none of the four'' is $1 - \Prob(A\cup B\cup C) - \Prob(D)$. Run the
three-event inclusion-exclusion for $A$, $B$, $C$, then add $D$ as a ninth region.
Since $D$ overlaps nothing, ``exactly one of the four'' is ``exactly one of the
three'' plus $\Prob(D)$, and ``exactly two of the four'' equals ``exactly two of the
three''.

\paragraph{From English to sets.}
The two tables below are the translation dictionary; every phrase in them appears
in exam questions. Write $p_A = \Prob(A)$, $p_B = \Prob(B)$, $p_{AB} = \Prob(A \cap B)$.
\begin{center}
\begin{tabular}{@{}lll@{}}
\toprule
English phrase (two events) & Set expression & Probability \\
\midrule
$A$ or $B$; at least one            & $A \cup B$                                 & $p_A + p_B - p_{AB}$ \\
$A$ and $B$; both                   & $A \cap B$                                 & $p_{AB}$ \\
not $A$                             & $\cc{A}$                                   & $1 - p_A$ \\
neither $A$ nor $B$                 & $\cc{A} \cap \cc{B} = \cc{(A\cup B)}$      & $1 - p_A - p_B + p_{AB}$ \\
not both; at most one               & $\cc{A} \cup \cc{B} = \cc{(A\cap B)}$      & $1 - p_{AB}$ \\
$A$ but not $B$; $A$ only           & $A \cap \cc{B}$                            & $p_A - p_{AB}$ \\
exactly one of $A$, $B$             & $(A\cap\cc{B}) \cup (\cc{A}\cap B)$        & $p_A + p_B - 2p_{AB}$ \\
\bottomrule
\end{tabular}
\end{center}
For three events write $\Sigma_1 = \Prob(A)+\Prob(B)+\Prob(C)$,
$\Sigma_2 = \Prob(A\cap B)+\Prob(A\cap C)+\Prob(B\cap C)$, and
$T = \Prob(A\cap B\cap C)$. Region numbers refer to the eight-region table.
\begin{center}
\begin{tabular}{@{}lll@{}}
\toprule
English phrase (three events) & Set expression & Probability \\
\midrule
at least one                        & $A \cup B \cup C$                          & $\Sigma_1 - \Sigma_2 + T$ \\
none                                & $\cc{A}\cap\cc{B}\cap\cc{C}$               & $1 - \Sigma_1 + \Sigma_2 - T$ \\
all three                           & $A \cap B \cap C$                          & $T$ \\
$A$ only                            & $A \cap \cc{B} \cap \cc{C}$                & $\Prob(A) - \Prob(AB) - \Prob(AC) + T$ \\
exactly one                         & regions 5, 6, 7                            & $\Sigma_1 - 2\Sigma_2 + 3T$ \\
exactly two                         & regions 2, 3, 4                            & $\Sigma_2 - 3T$ \\
at least two                        & regions 1, 2, 3, 4                         & $\Sigma_2 - 2T$ \\
at most one                         & complement of ``at least two''             & $1 - \Sigma_2 + 2T$ \\
\bottomrule
\end{tabular}
\end{center}
Every three-event formula is a sum of regions from the eight-region table; when in
doubt, fill the regions and add.

\subsection{Countable additivity and geometric series}

Axiom A3 covers infinite sequences of disjoint events, and the exam tests it with
regions whose areas or probabilities form a geometric progression: the $i$th region
has probability $r^i$ (or $k r^i$) for $i = 1, 2, \ldots$, with $0 < r < 1$. The regions
are disjoint, so the probability of hitting one of the first $N$ is the sum of the
first $N$ terms, and the probability of hitting any of infinitely many is the sum of
the series.

\begin{keyfact}[Geometric sums for disjoint regions]
For $0 < r < 1$,
\[
  \sum_{i=1}^{N} r^i = \frac{r\,(1 - r^N)}{1 - r},
  \qquad
  \sum_{i=1}^{\infty} r^i = \frac{r}{1 - r},
  \qquad
  \sum_{i=0}^{\infty} r^i = \frac{1}{1 - r}.
\]
The infinite sum must not exceed $1$ (A3 with A2), so a design with areas $r^i$,
$i \ge 1$, is admissible only when $r \le 1/2$. If the union of the regions is the
whole space, the constant $k$ in $\Prob(A_i) = k r^i$ is fixed by
$k\,r/(1-r) = 1$.
\end{keyfact}

The probability of missing all $N$ regions is $1 - r(1 - r^N)/(1-r)$, which decreases
to the floor $1 - r/(1-r) = (1 - 2r)/(1-r)$ as $N \to \infty$. A ``minimum $N$ such
that the miss probability is below $p$'' question is solved by setting the closed
form below $p$, isolating $r^N$, and taking logarithms:
\[
  r^N < c \iff N > \frac{\ln c}{\ln r}
  \qquad (\text{the inequality flips because } \ln r < 0),
\]
then rounding up to the next integer. If $p$ is at or below the floor, no finite
$N$ works, and no infinite design works either.

\subsection{Testing whether two described events are equal}

A question may define $F$ and $G$ in words, list five events in words, and ask how
many of the five equal $F \cap G$ (or $F \cup G$). Reading alone is unreliable; the
descriptions are chosen to differ by one boundary case. The systematic method:
\begin{enumerate}[nosep]
  \item Choose coordinates for an outcome. For claims in two years, an outcome is
        the pair $(N_1, N_2)$ of claim counts.
  \item Coarsen each coordinate to the fewest classes that every description in the
        list can still distinguish. If the descriptions mention ``none'', ``exactly
        one'', and ``two or more'', use the classes $0$, $1$, $2^{+}$.
  \item Draw the grid of classes and mark the cells that belong to the target event.
  \item For each listed description, mark its cells and compare. Equal events have
        identical cell sets; one differing cell settles the case.
\end{enumerate}
The listed descriptions are usually algebraic rewrites of the target, and each
rewrite is right or wrong by one of the set laws:
\begin{itemize}[nosep]
  \item De Morgan: ``not ($\cc{F}$ and $\cc{G}$)'' equals $F \cup G$; ``not
        ($\cc{F}$ or $\cc{G}$)'' equals $F \cap G$.
  \item Disjoint decomposition: $F \cup G = F \cup (\cc{F} \cap G)$.
  \item Absorption: $F \cup (F \cap H) = F$ and $F \cap (F \cup H) = F$; adding a
        subset of $F$ to a union with $F$ changes nothing.
  \item Distribution: $F \cap (G \cup H) = (F\cap G) \cup (F \cap H)$.
  \item Given one coordinate, a condition on a total becomes a condition on the
        other coordinate: with $N_1 = 1$ fixed, ``$N_1 + N_2 \ge 2$'' is
        ``$N_2 \ge 1$''.
\end{itemize}
A description that is a proper subset of the target (it adds a condition) or a
proper superset (it drops one) is not equal to it, even when it is ``almost'' the
same. Watch for ``more than one'' (at least two) versus ``one or more'' (at least one),
and for a condition on the two-year total that silently allows a year with zero.

\subsection{Worked examples}

\begin{example}[Three-event coverage survey, none of the three]
An insurance company surveys its policyholders about three types of coverage:
auto, homeowners, and life. The survey finds that 40\% have auto coverage, 35\%
have homeowners coverage, and 25\% have life coverage. It also finds that 15\%
have both auto and homeowners, 10\% have both auto and life, 8\% have both
homeowners and life, and 5\% have all three coverages.

Calculate the percentage of policyholders with none of the three coverages.

(A) 22\% \quad (B) 28\% \quad (C) 33\% \quad (D) 38\% \quad (E) 72\%
\end{example}
\begin{solution}
Let $A$, $H$, $L$ be the events that a randomly chosen policyholder has auto,
homeowners, and life coverage. ``None of the three'' is $\cc{(A\cup H\cup L)}$, so
by the complement rule and three-event inclusion-exclusion,
\[
  \Prob(A \cup H \cup L) = 0.40 + 0.35 + 0.25 - 0.15 - 0.10 - 0.08 + 0.05 = 0.72,
\]
and $\Prob(\text{none}) = 1 - 0.72 = 0.28$.

Answer: (B).
\end{solution}

\begin{example}[Exactly one of three]
An insurance agency's records show that, for a randomly selected customer,
\begin{itemize}[nosep]
  \item the probability of having auto coverage is $0.50$, homeowners coverage
        $0.40$, and life coverage $0.30$;
  \item the probability of having both auto and homeowners is $0.20$, both auto
        and life is $0.15$, and both homeowners and life is $0.10$;
  \item the probability of having all three is $0.05$.
\end{itemize}
Calculate the probability that a randomly selected customer has exactly one of the
three coverages.

(A) 0.35 \quad (B) 0.40 \quad (C) 0.45 \quad (D) 0.50 \quad (E) 0.55
\end{example}
\begin{solution}
Fill the eight regions in the standard order. All three: $0.05$.
Pairs only: auto--homeowners $0.20 - 0.05 = 0.15$; auto--life $0.15 - 0.05 = 0.10$;
homeowners--life $0.10 - 0.05 = 0.05$.
Singles only: auto $0.50 - 0.15 - 0.10 - 0.05 = 0.20$;
homeowners $0.40 - 0.15 - 0.05 - 0.05 = 0.15$; life $0.30 - 0.10 - 0.05 - 0.05 = 0.10$.
None: $1 - 0.80 = 0.20$.

Exactly one is the sum of the three ``only'' regions: $0.20 + 0.15 + 0.10 = 0.45$.

Formula check: $\Sigma_1 - 2\Sigma_2 + 3T = 1.20 - 2(0.45) + 3(0.05) = 1.20 - 0.90 + 0.15 = 0.45$.

Answer: (C).
\end{solution}

\begin{example}[Bounding a probability from the axioms alone]
For a randomly selected claim from an auto insurer's files, the probability that
the claim involves bodily injury is $0.60$ and the probability that it involves
property damage is $0.70$. No other information about the claims is available.

Determine the smallest possible value of the probability that a claim involves
both bodily injury and property damage.

(A) 0 \quad (B) 0.10 \quad (C) 0.30 \quad (D) 0.42 \quad (E) 0.60
\end{example}
\begin{solution}
Let $B$ be bodily injury and $D$ be property damage. Inclusion-exclusion gives
$\Prob(B \cap D) = \Prob(B) + \Prob(D) - \Prob(B \cup D) = 1.30 - \Prob(B \cup D)$.
Since $\Prob(B \cup D) \le 1$ (C5), $\Prob(B \cap D) \ge 1.30 - 1 = 0.30$. The
value $0.30$ is attained when $B \cup D = S$, so it is the minimum. (This is
Bonferroni's inequality, C10.)

The largest possible value is $\min(0.60, 0.70) = 0.60$, attained when
$B \subseteq D$ (C4). The choice $0.42 = (0.60)(0.70)$ assumes independence, which
is not given.

Answer: (C).
\end{solution}

\begin{example}[Complements and De Morgan]
A property insurer's records show that, for a randomly selected policyholder, the
probability of not having auto coverage is $0.65$, the probability of not having
homeowners coverage is $0.55$, and the probability of having neither coverage is
$0.35$.

Calculate the probability that a randomly selected policyholder has both auto and
homeowners coverage.

(A) 0.10 \quad (B) 0.15 \quad (C) 0.20 \quad (D) 0.35 \quad (E) 0.45
\end{example}
\begin{solution}
Let $A$ and $H$ be auto and homeowners coverage. Translate each statement.
$\Prob(\cc{A}) = 0.65$ so $\Prob(A) = 0.35$.
$\Prob(\cc{H}) = 0.55$ so $\Prob(H) = 0.45$.
``Neither'' is $\cc{A} \cap \cc{H} = \cc{(A \cup H)}$ by De Morgan, so
$\Prob(A \cup H) = 1 - 0.35 = 0.65$.
Then
\[
  \Prob(A \cap H) = \Prob(A) + \Prob(H) - \Prob(A \cup H) = 0.35 + 0.45 - 0.65 = 0.15 .
\]
Table check: $A \cap \cc{H} = 0.35 - 0.15 = 0.20$, $\cc{A} \cap H = 0.45 - 0.15 = 0.30$,
neither $= 0.35$; the four regions sum to $0.15 + 0.20 + 0.30 + 0.35 = 1.00$.

Answer: (B).
\end{solution}

\begin{example}[Translating ``only'', ``none'', ``no customer has both'', ``at most one'']
An insurance agency has 200 customers. Each customer has some combination of auto,
homeowners, and renters coverage, or none of them. The agency's records show:
\begin{itemize}[nosep]
  \item 120 customers have auto coverage;
  \item 60 customers have homeowners coverage and 50 have renters coverage; no
        customer has both homeowners and renters coverage;
  \item 65 customers have auto coverage only;
  \item 25 customers have none of the three coverages.
\end{itemize}
Calculate the number of customers who have at most one of the three coverages.

(A) 55 \quad (B) 65 \quad (C) 120 \quad (D) 145 \quad (E) 175
\end{example}
\begin{solution}
Let $A$, $H$, $R$ be the three coverages and work with counts out of 200.
``No customer has both homeowners and renters'' means $H \cap R = \emptyset$, so
$A \cap H \cap R$ is also empty and every two-coverage customer has auto plus one
other. ``Auto only'' is 65, so $120 - 65 = 55$ customers have auto plus another
coverage; these are exactly the two-coverage customers. Since $H$ and $R$ are
disjoint, $|H \cup R| = 60 + 50 = 110$, of whom 55 also have auto, leaving
$110 - 55 = 55$ with homeowners only or renters only. Exactly one coverage:
$65 + 55 = 120$. Check: $120 + 55 + 25 = 200$.

``At most one'' is none or exactly one: $25 + 120 = 145$.

Answer: (D).
\end{solution}

\begin{example}[Unions with complements]
For a randomly selected policy in a commercial lines portfolio, let $A$ be the event
that the policy carries general liability coverage and $B$ the event that it carries
property coverage. You are given
\[
  \Prob(A \cup B) = 0.78, \qquad \Prob(A \cup \cc{B}) = 0.67, \qquad \Prob(B) = 0.55 .
\]
Calculate the probability that the policy carries both coverages.

(A) 0.22 \qquad (B) 0.33 \qquad (C) 0.45 \qquad (D) 0.55 \qquad (E) 0.67
\end{example}
\begin{solution}
Write inclusion-exclusion for each union and add:
\begin{align*}
  \Prob(A \cup B) &= \Prob(A) + \Prob(B) - \Prob(A \cap B), \\
  \Prob(A \cup \cc{B}) &= \Prob(A) + \Prob(\cc{B}) - \Prob(A \cap \cc{B}), \\
  0.78 + 0.67 &= 2\Prob(A) + 1 - \bigl[\Prob(A\cap B) + \Prob(A \cap \cc{B})\bigr]
               = 2\Prob(A) + 1 - \Prob(A),
\end{align*}
so $\Prob(A) = 1.45 - 1 = 0.45$. Then
$\Prob(A \cap B) = \Prob(A) + \Prob(B) - \Prob(A \cup B) = 0.45 + 0.55 - 0.78 = 0.22$.

Shorter route: $A \cup \cc{B}$ is the complement of $\cc{A} \cap B$ (De Morgan), so
$\Prob(\cc{A} \cap B) = 1 - 0.67 = 0.33$, and
$\Prob(A \cap B) = \Prob(B) - \Prob(\cc{A}\cap B) = 0.55 - 0.33 = 0.22$.

Table check: $A\cap B = 0.22$, $A \cap \cc{B} = 0.45 - 0.22 = 0.23$,
$\cc{A} \cap B = 0.33$, neither $= 1 - 0.78 = 0.22$; the sum is $1.00$.

Answer: (A).
\end{solution}

\begin{example}[Disjoint pair inside a three-product profile]
An agent sells disability income insurance, term life insurance, and whole life
insurance. No client holds both term life and whole life. The agent's client
profile is:
\begin{enumerate}[nosep,label=(\roman*)]
  \item 20\% of clients hold none of the three products;
  \item 63\% of clients hold disability income insurance;
  \item three times as many clients hold term life as hold whole life;
  \item 23\% of clients hold exactly two of the three products;
  \item 12\% of clients hold term life but not disability income insurance.
\end{enumerate}
Calculate the percentage of clients who hold both disability income and whole life
insurance.

(A) 5\% \qquad (B) 8\% \qquad (C) 12\% \qquad (D) 15\% \qquad (E) 18\%
\end{example}
\begin{solution}
Let $D$, $T$, $W$ be the three products. Since $T \cap W = \emptyset$, the regions
$T \cap W$ only and $D \cap T \cap W$ are empty, leaving six regions: $d$ ($D$ only),
$t$ ($T$ only), $w$ ($W$ only), $dt$, $dw$, and $n$ (none). Translate:
\[
  n = 0.20, \quad d + dt + dw = 0.63, \quad t + dt = 3(w + dw), \quad dt + dw = 0.23,
  \quad t = 0.12 .
\]
Fact (v) gives $t$ directly because $T$ but not $D$ cannot include any $W$.
Outside $D$ lie $t$, $w$, and $n$, with total $1 - 0.63 = 0.37$, so
$w = 0.37 - 0.20 - 0.12 = 0.05$. Substitute into the ratio:
$0.12 + dt = 3(0.05 + dw) = 0.15 + 3\,dw$, so $dt = 0.03 + 3\,dw$. With
$dt + dw = 0.23$: $0.03 + 4\,dw = 0.23$, so $dw = 0.05$ and $dt = 0.18$.

Check: $d = 0.63 - 0.18 - 0.05 = 0.40$, and
$0.40 + 0.12 + 0.05 + 0.18 + 0.05 + 0.20 = 1.00$. Term life totals
$0.12 + 0.18 = 0.30$, whole life $0.05 + 0.05 = 0.10$, ratio $3$.

Answer: (A).
\end{solution}

\begin{example}[Three overlapping benefits and a fourth disjoint one]
A company offers its 250 employees four voluntary benefits: dental, vision, legal,
and a commuter plan. Enrollment records show:
\begin{enumerate}[nosep,label=(\roman*)]
  \item 96 are enrolled in dental, 58 in vision, and 40 in legal;
  \item 27 are enrolled in dental and vision, 18 in dental and legal, and 12 in
        vision and legal;
  \item 7 are enrolled in all three of dental, vision, and legal;
  \item 35 are enrolled in the commuter plan, and none of these is enrolled in
        dental, vision, or legal.
\end{enumerate}
Calculate the number of employees enrolled in exactly one of the four benefits.

(A) 71 \qquad (B) 101 \qquad (C) 114 \qquad (D) 136 \qquad (E) 179
\end{example}
\begin{solution}
Let $D$, $V$, $L$, $K$ be the four benefits. For the three overlapping ones,
$\Sigma_1 = 96 + 58 + 40 = 194$, $\Sigma_2 = 27 + 18 + 12 = 57$, $T = 7$.
Exactly one of $D$, $V$, $L$ is $\Sigma_1 - 2\Sigma_2 + 3T = 194 - 114 + 21 = 101$.
Region check: $D$ only $= 96 - 20 - 11 - 7 = 58$, $V$ only $= 58 - 20 - 5 - 7 = 26$,
$L$ only $= 40 - 11 - 5 - 7 = 17$, total $101$.

$K$ is disjoint from $D \cup V \cup L$, so every one of its 35 members has exactly one
benefit, and none of them changes the count for $D$, $V$, $L$. Exactly one of the
four: $101 + 35 = 136$.

For reference, $|D \cup V \cup L| = 194 - 57 + 7 = 144$, $|D\cup V\cup L\cup K| = 144 + 35 = 179$
by finite additivity, and none of the four $= 250 - 179 = 71$.

Answer: (D).
\end{solution}

\begin{example}[Countable additivity with geometric areas]
A carnival game uses a board of area 1 painted in red and blue regions. A
blindfolded player marks a point on the board and wins if the point lies in a blue
region. The board can be painted with any number of red regions, and the red regions
are designed so that the $i$th red region has area $(2/5)^i$, $i = 1, 2, \ldots$.

Calculate the minimum number of red regions that makes the probability of winning
less than 34\%.

(A) 2 \qquad (B) 3 \qquad (C) 4 \qquad (D) 5 \qquad (E) 6
\end{example}
\begin{solution}
With $N$ red regions, the red regions are disjoint, so by additivity the losing
probability is the sum of their areas and the winning probability is
\[
  w_N = 1 - \sum_{i=1}^{N} \Bigl(\frac{2}{5}\Bigr)^i
      = 1 - \frac{(2/5)\bigl(1 - (2/5)^N\bigr)}{1 - 2/5}
      = 1 - \frac{2}{3}\Bigl(1 - \Bigl(\frac{2}{5}\Bigr)^N\Bigr)
      = \frac{1}{3} + \frac{2}{3}\Bigl(\frac{2}{5}\Bigr)^N .
\]
Require $w_N < 0.34$:
\[
  \frac{2}{3}\Bigl(\frac{2}{5}\Bigr)^N < 0.34 - \frac{1}{3} = \frac{1}{150}
  \iff \Bigl(\frac{2}{5}\Bigr)^N < \frac{1}{100}
  \iff N > \frac{\ln(0.01)}{\ln(0.4)} = 5.03 .
\]
The smallest integer is $N = 6$. Check: $w_5 = 1/3 + (2/3)(0.01024) = 0.3402$, not
below $0.34$; $w_6 = 1/3 + (2/3)(0.004096) = 0.3361$.

The winning probability never drops below the floor $1/3$: the infinite series of red
areas sums to $(2/5)/(3/5) = 2/3$. A target of ``less than 33\%'' would have no
answer, with any finite or infinite number of red regions.

Answer: (E).
\end{solution}

\begin{example}[Which described events equal $F \cup G$]
A policyholder holds an auto policy for two years. Let $F$ be the event that the
policyholder has at most one accident in year one, and $G$ the event that the
policyholder has no accidents in year two. Consider the following events:
\begin{enumerate}[nosep,label=(\roman*)]
  \item The policyholder has at most one accident in the two-year period.
  \item It is not the case that the policyholder has two or more accidents in year
        one and at least one accident in year two.
  \item The policyholder has at most one accident in year one, or has two or more
        accidents in year one and no accidents in year two.
  \item The policyholder has at most one accident in year one, or has no accidents
        in either year.
  \item The policyholder has at most one accident in year one, or has more
        accidents in year one than in year two and no accidents in year two.
\end{enumerate}
Determine the number of events from the list that are the same as $F \cup G$.

(A) None \qquad (B) Exactly one \qquad (C) Exactly two \qquad (D) Exactly three \qquad (E) All
\end{example}
\begin{solution}
Let $N_1$ and $N_2$ be the accident counts in the two years. Every description
distinguishes only the classes $0$, $1$, and $2^{+}$ for each year, so the outcome
grid has nine cells. $F$ is the rows $N_1 = 0$ and $N_1 = 1$; $G$ is the column
$N_2 = 0$; $F \cup G$ is those two rows plus the cell $(2^{+}, 0)$, seven cells in all.
A bullet marks membership.
\begin{center}
\begin{tabular}{@{}lcccccc@{}}
\toprule
$(N_1, N_2)$ & $F \cup G$ & (i) & (ii) & (iii) & (iv) & (v) \\
\midrule
$(0,0)$          & $\bullet$ & $\bullet$ & $\bullet$ & $\bullet$ & $\bullet$ & $\bullet$ \\
$(0,1)$          & $\bullet$ & $\bullet$ & $\bullet$ & $\bullet$ & $\bullet$ & $\bullet$ \\
$(0,2^{+})$      & $\bullet$ &           & $\bullet$ & $\bullet$ & $\bullet$ & $\bullet$ \\
$(1,0)$          & $\bullet$ & $\bullet$ & $\bullet$ & $\bullet$ & $\bullet$ & $\bullet$ \\
$(1,1)$          & $\bullet$ &           & $\bullet$ & $\bullet$ & $\bullet$ & $\bullet$ \\
$(1,2^{+})$      & $\bullet$ &           & $\bullet$ & $\bullet$ & $\bullet$ & $\bullet$ \\
$(2^{+},0)$      & $\bullet$ &           & $\bullet$ & $\bullet$ &           & $\bullet$ \\
$(2^{+},1)$      &           &           &           &           &           &           \\
$(2^{+},2^{+})$  &           &           &           &           &           &           \\
\bottomrule
\end{tabular}
\end{center}
(i) is the event $N_1 + N_2 \le 1$, three cells; it omits $(0, 2^{+})$, which is in
$F$. Not equal.
(ii) is the complement of $\cc{F} \cap \cc{G}$, which by De Morgan is $F \cup G$. Equal.
(iii) is $F \cup (\cc{F} \cap G)$, the disjoint decomposition of $F \cup G$. Equal.
(iv) is $F \cup \{(0,0)\}$; the added cell is already in $F$, so this is just $F$,
which lacks $(2^{+}, 0)$. Not equal.
(v) is $F \cup \{N_1 > N_2,\ N_2 = 0\} = F \cup \{N_1 \ge 1,\ N_2 = 0\}$. The added
cells are $(1,0)$, already in $F$, and $(2^{+},0)$; the result is $F \cup G$. Equal.

Three of the five events equal $F \cup G$.

Answer: (D).
\end{solution}

\subsection{Traps}

\begin{trap}[``At least one'' versus ``exactly one'' versus ``at most one'']
``At least one of $A$, $B$'' is $A \cup B$ and includes the intersection. ``Exactly
one'' excludes the intersection: $\Prob(A) + \Prob(B) - 2\Prob(A\cap B)$. ``At most
one'' is the complement of ``both'': $1 - \Prob(A \cap B)$. For three events,
``at most one'' is not $1 - T$; it is $1 - (\Sigma_2 - 2T)$, the complement of ``at
least two''. Likewise ``neither $A$ nor $B$'' is $\cc{(A\cup B)}$ while ``not both''
is $\cc{(A \cap B)}$; De Morgan's laws tell you which is which. \end{trap}

\begin{trap}[$\Prob(A \cap \cc{B})$ is not $\Prob(A) - \Prob(B)$]
``$A$ but not $B$'' has probability $\Prob(A) - \Prob(A \cap B)$. Subtracting
$\Prob(B)$ instead removes outcomes that were never in $A$. The identity
$\Prob(B \setminus A) = \Prob(B) - \Prob(A)$ holds only when $A \subseteq B$.
\end{trap}

\begin{trap}[Assuming disjoint or independent when the problem does not say so]
If a question gives $\Prob(A)$ and $\Prob(B)$ and asks for $\Prob(A \cup B)$ with
no other information, the answer is not determined; it lies between
$\max(\Prob(A),\Prob(B))$ and $\min(1, \Prob(A)+\Prob(B))$. Adding the probabilities
assumes disjointness; multiplying them assumes independence. Use either only when
the problem states it or the structure forces it (a policyholder cannot be both
under 25 and over 65).
\end{trap}

\begin{trap}[``$A$ and $B$'' includes those who also have $C$]
In a three-event problem, ``15\% have both auto and homeowners'' is
$\Prob(A \cap H) = 0.15$, which includes the policyholders who also have life.
The region ``auto and homeowners only'' is $\Prob(A\cap H) - \Prob(A\cap H\cap L)$.
Conversely, when the problem says ``15\% have auto and homeowners but not life'',
it is giving the region directly. Confusing the two shifts the answer by a multiple of $T$.
\end{trap}

\begin{trap}[$\Prob(A \cup \cc{B})$ is not $1 - \Prob(A \cup B)$]
The complement of $A \cup B$ is $\cc{A} \cap \cc{B}$, not $A \cup \cc{B}$. The event
$A \cup \cc{B}$ is three of the four cells of the $2 \times 2$ table (everything except
$\cc{A} \cap B$), so its probability is $1 - \Prob(\cc{A} \cap B)$, and it is
typically large. When two unions involving complements are given, add the two
inclusion-exclusion equations rather than guessing a relation between them.
\end{trap}

\begin{trap}[An infinite union of disjoint regions has a floor, and a ceiling]
With region probabilities $r^i$, $i \ge 1$, the probability of missing every
region decreases toward $(1-2r)/(1-r)$, not toward $0$; a question that asks for a miss
probability below that floor has no answer. In the other direction, the series must
sum to at most $1$, so $r^i$ areas on a board of area $1$ require $r \le 1/2$. When
solving $r^N < c$, dividing by $\ln r$ (negative) reverses the inequality, and the
answer is the next integer above the bound, not the nearest integer.
\end{trap}

\begin{trap}[Coarsening the outcome grid too far]
The classes $0$, $1$, $2^{+}$ are enough when every description compares a count
to a fixed number. A description that compares the two years to each other (``more
accidents in year one than in year two'') can split the $2^{+}$ class: the outcomes
$(3,2)$ and $(2,3)$ both sit in the cell $(2^{+}, 2^{+})$ but fall on opposite sides of
the comparison. Either refine the classes or test the ambiguous cell with explicit
outcomes before declaring two events equal.
\end{trap}

\subsection{Practice problems}

\begin{practice}
A survey of homeowners insurance policyholders finds that 50\% have a security
system, 35\% have a fire sprinkler system, and 25\% have a backup generator.
Further, 15\% have both a security system and a sprinkler system, 10\% have both a
security system and a generator, 5\% have both a sprinkler system and a generator,
and 2\% have all three.

Calculate the percentage of policyholders with none of the three.

(A) 12\% \quad (B) 18\% \quad (C) 20\% \quad (D) 28\% \quad (E) 32\%
\end{practice}

\begin{practice}
For a randomly selected policyholder, let $A$ be the event that the policyholder
files a claim in the first policy year and $B$ the event that the policyholder
files a claim in the second policy year. You are given
$\Prob(\cc{A}) = 0.60$, $\Prob(B) = 0.35$, and $\Prob(A \cup B) = 0.60$.

Calculate the probability that the policyholder files a claim in exactly one of the
two years.

(A) 0.30 \quad (B) 0.45 \quad (C) 0.55 \quad (D) 0.60 \quad (E) 0.75
\end{practice}

\begin{practice}
For a randomly selected claim, the probability that the claim involves a rental
vehicle is $0.75$ and the probability that it involves a towing charge is $0.65$.
No other information is available.

Determine the largest possible value of the probability that the claim involves
neither a rental vehicle nor a towing charge.

(A) 0 \quad (B) 0.0875 \quad (C) 0.25 \quad (D) 0.35 \quad (E) 0.40
\end{practice}

\begin{practice}
An insurer classifies each policy by whether it has a young driver ($Y$), a
sports car ($S$), and an accident in the past three years ($C$). For a randomly
selected policy,
\begin{gather*}
  \Prob(Y) = 0.30,\quad \Prob(S) = 0.25,\quad \Prob(C) = 0.20,\\
  \Prob(Y\cap S) = 0.10,\quad \Prob(Y \cap C) = 0.08,\quad \Prob(S \cap C) = 0.06,\quad
  \Prob(Y \cap S \cap C) = 0.04 .
\end{gather*}
Calculate the probability that a randomly selected policy has at most one of the
three characteristics.

(A) 0.16 \quad (B) 0.24 \quad (C) 0.76 \quad (D) 0.84 \quad (E) 0.90
\end{practice}

\begin{practice}
The probability that a randomly selected auto policyholder lacks at least one of
collision coverage and comprehensive coverage is $0.90$. The probability of having
collision coverage is $0.45$ and the probability of having comprehensive coverage
is $0.30$.

Calculate the probability that the policyholder has at least one of the two
coverages.

(A) 0.10 \quad (B) 0.35 \quad (C) 0.55 \quad (D) 0.65 \quad (E) 0.75
\end{practice}

\begin{practice}
Of 500 auto policyholders, 280 have collision coverage, 340 have comprehensive
coverage, and 60 have neither.

Calculate the number of policyholders who have collision coverage only.

(A) 80 \quad (B) 100 \quad (C) 160 \quad (D) 180 \quad (E) 220
\end{practice}

\begin{practice}
For a randomly selected policyholder, let $A$ be the event that the policyholder
files a claim this year and $B$ the event that the policyholder renews at year end.
You are given
\[
  \Prob(A \cup B) = 0.85, \qquad \Prob(A \cup \cc{B}) = 0.75, \qquad \Prob(\cc{A} \cup B) = 0.90 .
\]
Calculate the probability that at most one of the events $A$ and $B$ occurs.

(A) 0.15 \qquad (B) 0.25 \qquad (C) 0.35 \qquad (D) 0.40 \qquad (E) 0.50
\end{practice}

\begin{practice}
An insurance agent sells auto, homeowners, and renters insurance. No client holds
both homeowners and renters insurance. The agent's client profile is:
\begin{enumerate}[nosep,label=(\roman*)]
  \item 12\% of clients hold none of the three products;
  \item 43\% of clients hold exactly one of the three products;
  \item 73\% of clients hold auto insurance;
  \item 5\% of clients hold renters insurance but not auto insurance;
  \item twice as many clients hold both auto and homeowners insurance as hold both
        auto and renters insurance.
\end{enumerate}
Calculate the percentage of clients who hold homeowners insurance.

(A) 20\% \qquad (B) 28\% \qquad (C) 30\% \qquad (D) 38\% \qquad (E) 40\%
\end{practice}

\begin{practice}
A bank surveys 400 customers about four products: a checking account, a savings
account, a brokerage account, and a certificate of deposit. The survey finds:
\begin{enumerate}[nosep,label=(\roman*)]
  \item 130 hold a checking account, 100 a savings account, and 70 a brokerage
        account;
  \item 40 hold checking and savings, 30 hold checking and brokerage, and 20 hold
        savings and brokerage;
  \item 60 hold a certificate of deposit, and none of these holds a checking,
        savings, or brokerage account;
  \item 115 hold none of the four products.
\end{enumerate}
Calculate the number of customers who hold all three of the checking, savings, and
brokerage accounts.

(A) 15 \qquad (B) 25 \qquad (C) 30 \qquad (D) 45 \qquad (E) 60
\end{practice}

\begin{practice}
Each incoming claim is assigned to exactly one adjuster from the sequence
$A_1, A_2, A_3, \ldots$, and the probability that a claim is assigned to adjuster
$A_i$ is $k\,(3/4)^i$ for $i = 1, 2, \ldots$, where $k$ is a constant.

Calculate the smallest $n$ such that the probability that a claim is assigned to one
of the first $n$ adjusters is at least $0.95$.

(A) 9 \qquad (B) 10 \qquad (C) 11 \qquad (D) 12 \qquad (E) 13
\end{practice}

\begin{practice}
A policyholder holds a policy for two years. Let $F$ be the event that the
policyholder has no claims in year one and $G$ the event that the policyholder has
at least two claims in year two. Consider the following events:
\begin{enumerate}[nosep,label=(\roman*)]
  \item The policyholder has at least one claim in year one and at most one claim in
        year two.
  \item It is not the case that the policyholder has no claims in year one or has at
        least two claims in year two.
  \item The policyholder has at least one claim in year one, and has either no claims
        or exactly one claim in year two.
  \item The policyholder has exactly one claim in year one and at most one claim in
        year two, or has two or more claims in year one and at most one claim in
        year two.
  \item The policyholder has at least one claim in year one and at most one claim in
        year two, or has exactly one claim in year one and no claims in year two.
\end{enumerate}
Determine the number of events from the list that are the same as $\cc{(F \cup G)}$.

(A) None \qquad (B) Exactly one \qquad (C) Exactly two \qquad (D) Exactly three \qquad (E) All
\end{practice}

\begin{practice}
An auto insurer has 5{,}000 policies. Each policy is classified as urban or rural,
as covering a sedan or a truck, and as commuting use or pleasure use. Of these
policies, 2{,}100 are urban, 2{,}800 cover sedans, and 3{,}200 are for commuting
use. Further, 1{,}150 are urban sedans, 1{,}900 are commuting sedans, and 1{,}350
are urban commuting policies. Finally, 600 policies are rural trucks for pleasure use.

Calculate the number of policies that are urban trucks for pleasure use.

(A) 300 \qquad (B) 450 \qquad (C) 650 \qquad (D) 700 \qquad (E) 950
\end{practice}

\subsection{Answers to practice problems}

\begin{enumerate}[nosep,label=\textbf{\arabic*.}]
  \item \textbf{(B)}. $\Prob(\text{at least one}) = 1.10 - 0.30 + 0.02 = 0.82$,
        so none is $0.18$.
  \item \textbf{(B)}. $\Prob(A) = 0.40$; $\Prob(A\cap B) = 0.40 + 0.35 - 0.60 = 0.15$;
        exactly one $= 0.40 + 0.35 - 2(0.15) = 0.45$ (equivalently $0.60 - 0.15$).
  \item \textbf{(C)}. Neither $= 1 - \Prob(R \cup T)$, and $\Prob(R \cup T) \ge \max(0.75, 0.65) = 0.75$
        by monotonicity, so neither $\le 0.25$, attained when $T \subseteq R$.
        (0.0875 assumes independence, which is not given.)
  \item \textbf{(D)}. At least two $= \Sigma_2 - 2T = 0.24 - 0.08 = 0.16$; at most one $= 0.84$.
  \item \textbf{(D)}. ``Lacks at least one'' is $\cc{(C \cap M)}$, so $\Prob(C \cap M) = 0.10$;
        $\Prob(C \cup M) = 0.45 + 0.30 - 0.10 = 0.65$.
  \item \textbf{(B)}. At least one $= 440$; both $= 280 + 340 - 440 = 180$;
        collision only $= 280 - 180 = 100$.
  \item \textbf{(E)}. $\Prob(A) = 0.85 + 0.75 - 1 = 0.60$; $\Prob(B) = 0.85 + 0.90 - 1 = 0.75$;
        $\Prob(A \cap B) = 0.60 + 0.75 - 0.85 = 0.50$; at most one $= 1 - 0.50 = 0.50$.
        (Neither $= 0.15$; exactly one $= 0.35$.)
  \item \textbf{(E)}. Exactly two $= 1 - 0.12 - 0.43 = 0.45$; auto only $= 0.73 - 0.45 = 0.28$;
        homeowners only $+$ renters only $= 0.43 - 0.28 = 0.15$, so homeowners only $= 0.10$;
        auto-renters $= 0.15$ and auto-homeowners $= 0.30$ from the ratio;
        homeowners $= 0.10 + 0.30 = 0.40$.
  \item \textbf{(A)}. The certificate holders are disjoint from the other three, so
        $|C \cup S \cup B| = 400 - 115 - 60 = 225 = 300 - 90 + T$, giving $T = 15$.
  \item \textbf{(C)}. The union of all $A_i$ is $S$, so $k(3/4)/(1/4) = 3k = 1$ and $k = 1/3$.
        $\Prob(A_1 \cup \cdots \cup A_n) = 1 - (3/4)^n \ge 0.95$ requires
        $n \ge \ln(0.05)/\ln(0.75) = 10.41$, so $n = 11$.
  \item \textbf{(E)}. The target is $\cc{F} \cap \cc{G} = \{N_1 \ge 1, N_2 \le 1\}$.
        (i) is the direct statement; (ii) is De Morgan; (iii) splits $N_2 \le 1$ into
        $N_2 = 0$ or $N_2 = 1$; (iv) splits $N_1 \ge 1$ into $N_1 = 1$ or $N_1 \ge 2$;
        (v) adds the cell $(1, 0)$, already in the target, so absorption leaves it unchanged.
  \item \textbf{(A)}. Union of urban, sedan, commuting $= 5000 - 600 = 4400 = 8100 - 4400 + T$,
        so $T = 700$; urban only $= 2100 - 1150 - 1350 + 700 = 300$.
\end{enumerate}

\subsection{Sample exam cross-reference}

After studying this section, test on the following questions from the sample exam.
\begin{itemize}[nosep]
  \item Bank 1, Q1: three-event inclusion-exclusion, none of the three (Example 1, Practice 1).
  \item Bank 1, Q3: add $\Prob(A \cup B)$ and $\Prob(A \cup \cc{B})$ to isolate $\Prob(A)$ (Example 6, Practice 7).
  \item Bank 1, Q5: three binary classifications as a three-event count; ``young, female, single'' is a single region (Example 2, Practice 12).
  \item Bank 1, Q87: disjoint pair inside three products, six-region linear system with an ``exactly two'' constraint and a ratio (Example 7, Practice 8).
  \item Bank 1, Q100: three overlapping sets plus a fourth disjoint from them; finite additivity for the union (Example 8, Practice 9).
  \item Bank 1, Q108: countable additivity over disjoint regions with geometric areas; minimum $N$ by logarithms (Example 9, Practice 10).
  \item Bank 1, Q153: event equality tested on an outcome grid; De Morgan, decomposition, absorption (Example 10, Practice 11).
\end{itemize}

\subsection{One-page summary}

\begin{itemize}[nosep]
  \item Sample space $S$: all outcomes. Event: a subset of $S$. Probability: a set
        function on the events with $\Prob(A) \ge 0$, $\Prob(S) = 1$, and countable
        additivity over pairwise disjoint events.
  \item Consequences: $\Prob(\emptyset) = 0$; $\Prob(\cc{A}) = 1 - \Prob(A)$;
        $A \subseteq B \Rightarrow \Prob(A) \le \Prob(B)$; $0 \le \Prob(A) \le 1$;
        $\Prob(A \cap \cc{B}) = \Prob(A) - \Prob(A \cap B)$.
  \item De Morgan: $\cc{(A \cup B)} = \cc{A} \cap \cc{B}$ (``neither'');
        $\cc{(A \cap B)} = \cc{A} \cup \cc{B}$ (``not both'').
  \item Inclusion-exclusion: $\Prob(A \cup B) = \Prob(A) + \Prob(B) - \Prob(A \cap B)$;
        $\Prob(A\cup B\cup C) = \Sigma_1 - \Sigma_2 + T$.
  \item Exactly one of two: $\Prob(A) + \Prob(B) - 2\Prob(A\cap B)$.
        Exactly one of three: $\Sigma_1 - 2\Sigma_2 + 3T$.
        Exactly two of three: $\Sigma_2 - 3T$. At least two of three: $\Sigma_2 - 2T$.
  \item Bounds with no other information: $\max(\Prob(A),\Prob(B)) \le \Prob(A \cup B) \le \min(1,\Prob(A)+\Prob(B))$
        and $\max(0, \Prob(A)+\Prob(B)-1) \le \Prob(A\cap B) \le \min(\Prob(A),\Prob(B))$.
        Boole: $\Prob(\bigcup A_i) \le \sum \Prob(A_i)$.
  \item Venn bookkeeping: fill regions from the triple intersection outward (all
        three; each pair only; each single only; none); every answer is a sum of regions.
  \item Translation: ``or''/``at least one'' is union; ``and''/``both'' is
        intersection; ``not'' is complement; ``but not'' and ``only'' subtract the
        intersection; ``every policyholder has at least one'' means the union is $S$;
        ``no one has both'' means disjoint.
  \item Unions with complements: $\Prob(A\cup B) + \Prob(A \cup \cc{B}) = 1 + \Prob(A)$;
        $\Prob(A \cup \cc{B}) = 1 - \Prob(\cc{A} \cap B)$; $\Prob(\cc{A}\cup\cc{B}) = 1 - \Prob(A \cap B)$.
        Derive any of these by adding two inclusion-exclusion equations.
  \item Disjoint pair inside three events: two regions vanish; name the six that
        remain, translate each fact into a linear equation, solve outside-in.
        A fourth set disjoint from the other three adds by finite additivity:
        $\Prob(A\cup B\cup C\cup D) = \Prob(A \cup B \cup C) + \Prob(D)$.
  \item Geometric regions: $\sum_{i=1}^{N} r^i = r(1-r^N)/(1-r)$,
        $\sum_{i=1}^{\infty} r^i = r/(1-r) \le 1$. Miss probability with $N$ regions
        is $1 - r(1-r^N)/(1-r)$, floor $(1-2r)/(1-r)$; solve $r^N < c$ as
        $N > \ln c / \ln r$ and round up.
  \item Event equality: draw the outcome grid (classes $0$, $1$, $2^{+}$ per
        coordinate), mark the target, mark each description; one differing cell
        breaks equality. Rewrites that preserve equality: De Morgan,
        $F \cup G = F \cup (\cc{F}\cap G)$, absorption $F \cup (F \cap H) = F$.
  \item Never assume disjointness or independence unless stated; check that the
        regions sum to $1$ (or to the stated count).
\end{itemize}

\clearpage
\section{Outcome (b): Combinatorics}

\subsection{What the outcome asks}

The exam tests classical probability: when every outcome of an experiment is equally likely,
\[
\Prob(A) = \frac{\text{number of outcomes in } A}{\text{total number of outcomes}},
\]
so the whole problem reduces to counting two sets correctly and consistently. The questions
ask you to decide whether order matters (permutation versus combination), whether selection is
with or without replacement, and how to split a group into several labelled subgroups
(multinomial partitions). The typical setting is selecting policies from a portfolio, forming
a committee from a group of policyholders, assigning claims to adjusters, or arranging files
in a queue. ``At least one'' questions are almost always fastest through the complement.
Exam arithmetic is done on a calculator, so answers are usually given to three decimals.

\subsection{Counting tools}

\paragraph{Multiplication principle.}
If a procedure consists of $k$ stages and stage $i$ can be done in $n_i$ ways regardless of
what happened at the earlier stages, the procedure can be done in $n_1 n_2 \cdots n_k$ ways.
Every formula below is a consequence of this principle.

\begin{keyfact}[Permutations: order matters, no replacement]
The number of ordered arrangements of all $n$ distinct objects is
\[
n! = n(n-1)(n-2)\cdots 2 \cdot 1, \qquad 0! = 1 .
\]
The number of ordered arrangements of $k$ objects chosen from $n$ distinct objects is
\[
{}_nP_k = \frac{n!}{(n-k)!} = n(n-1)\cdots(n-k+1) .
\]
\end{keyfact}

\begin{keyfact}[Combinations: order does not matter, no replacement]
The number of unordered subsets of size $k$ from $n$ distinct objects is
\[
\binom{n}{k} = \frac{n!}{k!\,(n-k)!} = \frac{{}_nP_k}{k!} .
\]
Useful identities:
\[
\binom{n}{k} = \binom{n}{n-k}, \qquad
\binom{n}{0} = \binom{n}{n} = 1, \qquad
\binom{n}{1} = n, \qquad
\binom{n}{2} = \frac{n(n-1)}{2},
\]
and Pascal's rule
\[
\binom{n}{k} = \binom{n-1}{k-1} + \binom{n-1}{k} .
\]
\end{keyfact}

Pascal's rule has a counting interpretation that recurs in exam problems: fix one particular
object; the $k$-subsets that contain it number $\binom{n-1}{k-1}$ and the $k$-subsets that
do not contain it number $\binom{n-1}{k}$. Dividing, the probability that a random $k$-subset
contains a specific object is $\binom{n-1}{k-1}/\binom{n}{k} = k/n$.

\begin{keyfact}[Binomial theorem]
\[
(x+y)^n = \sum_{k=0}^{n} \binom{n}{k} x^k y^{n-k}, \qquad
\sum_{k=0}^{n}\binom{n}{k} = 2^n, \qquad
\sum_{k=0}^{n}(-1)^k\binom{n}{k} = 0 .
\]
The identity $\sum_k \binom{n}{k}=2^n$ counts all subsets of an $n$-set. It is also why the
binomial probabilities in Topic 2 sum to one.
\end{keyfact}

\begin{keyfact}[Permutations with repeated objects (multinomial coefficient)]
The number of distinguishable arrangements of $n$ objects of which $n_1$ are identical of
type 1, $n_2$ identical of type 2, \dots, $n_r$ identical of type $r$
(with $n_1+\cdots+n_r=n$) is
\[
\binom{n}{n_1,\,n_2,\,\dots,\,n_r} = \frac{n!}{n_1!\,n_2!\cdots n_r!} .
\]
The same number counts the ways to partition $n$ distinct objects into $r$ labelled groups
of sizes $n_1,\dots,n_r$. It equals the product of successive combinations
$\binom{n}{n_1}\binom{n-n_1}{n_2}\cdots\binom{n_r}{n_r}$.
\end{keyfact}

For example, the word MISSISSIPPI has $11!/(1!\,4!\,4!\,2!) = 34650$ distinguishable
arrangements, and $12$ distinct claims can be divided among three adjusters, four claims each,
in $12!/(4!\,4!\,4!) = 34650$ ways.

\begin{keyfact}[With replacement]
Choosing $k$ times from $n$ objects with replacement, order recorded: $n^k$ outcomes.
Choosing $k$ times from $n$ objects with replacement, order ignored (a multiset): $\binom{n+k-1}{k}$ outcomes.
The second count is rarely needed on the exam; the first is needed constantly, in particular
as the denominator in ``at least one match'' problems.
\end{keyfact}

\begin{center}
\begin{tabular}{llp{3.4cm}p{5.2cm}}
\toprule
Order matters? & Replacement? & Count & Typical wording \\
\midrule
Yes & No  & ${}_nP_k = \dfrac{n!}{(n-k)!}$ & arrange, rank, sequence, fill distinct positions \\[6pt]
Yes & Yes & $n^k$ & each of $k$ items independently takes one of $n$ values \\[4pt]
No  & No  & $\dbinom{n}{k}$ & select, choose, committee, sample, hand \\[6pt]
No  & Yes & $\dbinom{n+k-1}{k}$ & multiset: only how many of each type is recorded \\
\bottomrule
\end{tabular}
\\[4pt]
Decision table for choosing $k$ from $n$ distinct objects.
\end{center}

\paragraph{Circular arrangements.}
$n$ distinct objects around a circle, where only relative position matters, can be arranged in
$(n-1)!$ ways: fix one object to remove the rotational symmetry and arrange the remaining $n-1$
in a row. If the circle can also be flipped over (a necklace), divide again by $2$.

\paragraph{Stars and bars.}
The number of ways to write $n$ as an ordered sum of $r$ nonnegative integers
$x_1 + x_2 + \cdots + x_r = n$, $x_i \ge 0$, is
\[
\binom{n+r-1}{r-1} .
\]
Picture $n$ identical stars and $r-1$ bars; each arrangement of the $n+r-1$ symbols is one
solution. If every $x_i \ge 1$ is required, give each variable one unit first and count
solutions of $y_1+\cdots+y_r = n-r$: $\binom{n-1}{r-1}$. For example, $10$ identical bonus
units can be distributed among $4$ agents in $\binom{13}{3}=286$ ways, and in
$\binom{9}{3}=84$ ways if every agent gets at least one. If all $286$ distributions are
equally likely, the probability that every agent gets at least one unit is $84/286 = 0.294$.

\paragraph{Restrictions.}
Two standard devices handle restrictions on arrangements.
\begin{itemize}[nosep]
\item \textbf{Objects that must be together}: glue them into one block, arrange the blocks,
then multiply by the number of internal orderings of the block. Three files that must be
adjacent among $7$ files in a row: $5!\cdot 3! = 720$ arrangements out of $7! = 5040$.
\item \textbf{Objects that must be separated}: arrange the unrestricted objects first, then
place the restricted objects in the gaps (including the two ends). No two of $3$ special
files adjacent among $7$: arrange the $4$ others ($4!$), choose $3$ of the $5$ gaps
($\binom{5}{3}$), order the special files in those gaps ($3!$): $24\cdot 10\cdot 6 = 1440$.
\end{itemize}

\subsection{Classical probability with counting}

With equally likely outcomes, $\Prob(A) = |A|/|S|$. The numerator and denominator must be
counted in the same way: both as ordered lists, or both as unordered sets. Either convention
gives the same probability if applied consistently.

\begin{keyfact}[The hypergeometric counting pattern]
A group contains $a$ objects of type I and $b$ objects of type II. A sample of $n$ objects is
drawn at random without replacement. The probability that the sample contains exactly $x$ of
type I is
\[
\Prob(X = x) = \frac{\dbinom{a}{x}\dbinom{b}{n-x}}{\dbinom{a+b}{n}} .
\]
With more than two types, the numerator becomes a product of one binomial coefficient per
type, and the denominator is $\binom{N}{n}$ with $N$ the total. Topic 2 names this the
hypergeometric distribution; on this outcome it is a counting pattern.
\end{keyfact}

\paragraph{Same problem, two methods.}
A portfolio holds $10$ policies: $4$ auto and $6$ homeowners. Three policies are selected at
random without replacement. Find the probability that exactly $2$ of the $3$ are auto.

\emph{Counting (unordered).} Total $\binom{10}{3} = 120$. Favourable
$\binom{4}{2}\binom{6}{1} = 6\cdot 6 = 36$. Probability $36/120 = 0.300$.

\emph{Sequential conditional probabilities (ordered).} The sequence auto, auto, home has
probability
\[
\frac{4}{10}\cdot\frac{3}{9}\cdot\frac{6}{8} = \frac{72}{720} = 0.100 .
\]
The sequences auto, home, auto and home, auto, auto each have the same probability (the
numerators $4,3,6$ appear in a different order, the denominators are $10,9,8$ regardless).
There are $\binom{3}{2}=3$ positions for the home policy, so the probability is
$3 \times 0.100 = 0.300$.

The two methods agree. The counting method is faster when the sample is unordered; the
sequential method is faster when the question involves only two or three draws, or asks
about a specific position (``the third policy selected is auto'').

\subsection{Worked examples}

\begin{example}[Committee with a constraint]
A group of $12$ policyholders consists of $7$ who hold auto coverage only and $5$ who hold
both auto and homeowners coverage. A committee of $5$ is chosen at random from the group.
Calculate the probability that the committee includes at least $2$ policyholders who hold
both coverages.
\par\smallskip\noindent (A) 0.44 \quad (B) 0.61 \quad (C) 0.69 \quad (D) 0.75 \quad (E) 0.81
\end{example}
\begin{solution}
Total committees: $\binom{12}{5} = 792$. Use the complement: fewer than $2$ from the
``both'' group means $0$ or $1$.
\[
\binom{5}{0}\binom{7}{5} + \binom{5}{1}\binom{7}{4} = 21 + 5\cdot 35 = 196 .
\]
\[
\Prob(\text{at least } 2) = 1 - \frac{196}{792} = \frac{596}{792} = 0.7525 .
\]
Answer: \textbf{(D)}.

The direct sum $\binom{5}{2}\binom{7}{3}+\binom{5}{3}\binom{7}{2}+\binom{5}{4}\binom{7}{1}+\binom{5}{5}\binom{7}{0}
= 350+210+35+1 = 596$ agrees. Choice (A) is the probability of exactly $2$, choice (C) is
the probability of at most $2$.
\end{solution}

\begin{example}[Sampling a portfolio with three policy types]
An insurer's portfolio consists of $20$ policies: $8$ auto, $7$ homeowners, and $5$ life.
Four policies are selected at random without replacement for audit. Calculate the probability
that the audit sample contains exactly $2$ auto, $1$ homeowners, and $1$ life policy.
\par\smallskip\noindent (A) 0.006 \quad (B) 0.058 \quad (C) 0.202 \quad (D) 0.263 \quad (E) 0.350
\end{example}
\begin{solution}
\[
\Prob = \frac{\binom{8}{2}\binom{7}{1}\binom{5}{1}}{\binom{20}{4}}
= \frac{28\cdot 7\cdot 5}{4845} = \frac{980}{4845} = 0.2023 .
\]
Answer: \textbf{(C)}.

Choice (B) results from writing $8\cdot 7\cdot 5$ in the numerator (forgetting that the
two auto policies are an unordered pair, $\binom{8}{2}=28$, not $8\cdot 7$ or $8$).
Choice (A) results from an ordered numerator $8\cdot7\cdot5\cdot3$ over the ordered
denominator ${}_{20}P_4$ without multiplying by the $4!/2! = 12$ arrangements of the types.
\end{solution}

\begin{example}[At least one match: birthday pattern]
A claims department has $8$ adjusters. Each of $4$ incoming claims is assigned at random to
one of the $8$ adjusters, independently of the other claims. Calculate the probability that at
least two of the four claims are assigned to the same adjuster.
\par\smallskip\noindent (A) 0.41 \quad (B) 0.50 \quad (C) 0.55 \quad (D) 0.59 \quad (E) 0.66
\end{example}
\begin{solution}
Assignments are with replacement and ordered (claim 1 to adjuster $x$, claim 2 to adjuster
$y$, \dots), so the sample space has $8^4 = 4096$ equally likely outcomes. The complement of
``at least two share'' is ``all four different'', with ${}_8P_4 = 8\cdot7\cdot6\cdot5 = 1680$
outcomes.
\[
\Prob(\text{at least one shared adjuster}) = 1 - \frac{1680}{4096} = 1 - 0.4102 = 0.5898 .
\]
Answer: \textbf{(D)}.

The same structure gives the birthday problem: with $23$ people and $365$ equally likely
birthdays, $1 - {}_{365}P_{23}/365^{23} = 0.507$.
\end{solution}

\begin{example}[Arrangement with a separation restriction]
Seven claim files are placed in a queue in random order. Three of the files are auto claims
and four are homeowners claims. Calculate the probability that no two auto claim files are
adjacent in the queue.
\par\smallskip\noindent (A) $1/7$ \quad (B) $2/7$ \quad (C) $3/7$ \quad (D) $4/7$ \quad (E) $5/7$
\end{example}
\begin{solution}
Treat the files as distinct. Total orderings: $7! = 5040$. Arrange the four homeowners files
first: $4! = 24$ ways. They create $5$ gaps (three interior, two ends). Choose $3$ of the $5$
gaps for the auto files, $\binom{5}{3} = 10$, and order the auto files within the chosen gaps,
$3! = 6$. Favourable orderings: $24\cdot 10\cdot 6 = 1440$.
\[
\Prob = \frac{1440}{5040} = \frac{2}{7} .
\]
Answer: \textbf{(B)}.

The same answer is reached treating the files as indistinguishable within type: the queue is
a pattern of $3$ A's and $4$ H's, $\binom{7}{3} = 35$ patterns, of which $\binom{5}{3} = 10$
have no two A's adjacent; $10/35 = 2/7$. Either model works if used for both numerator and
denominator. For contrast, the probability that the three auto files are all adjacent is
$5!\cdot 3!/7! = 720/5040 = 1/7$, choice (A).
\end{solution}

\begin{example}[Multinomial partition of claims among adjusters]
Twelve claims of different sizes are to be divided at random among three adjusters, with each
adjuster receiving exactly four claims. Calculate the probability that the three largest
claims are all assigned to the same adjuster.
\par\smallskip\noindent (A) 0.018 \quad (B) 0.027 \quad (C) 0.036 \quad (D) 0.045 \quad (E) 0.055
\end{example}
\begin{solution}
\emph{Counting.} Total partitions: $\dfrac{12!}{4!\,4!\,4!} = 34650$. Favourable: choose which
adjuster receives the three largest claims ($3$ ways), choose that adjuster's fourth claim from
the remaining $9$ ($\binom{9}{1}=9$), then split the remaining $8$ claims between the other two
adjusters ($\binom{8}{4} = 70$). Favourable count $3\cdot 9\cdot 70 = 1890$.
\[
\Prob = \frac{1890}{34650} = \frac{3}{55} = 0.0545 .
\]
\emph{Sequential.} Wherever the largest claim goes, that adjuster has $3$ open slots among
the $11$ remaining; the second largest lands there with probability $3/11$, and then the
third largest with probability $2/10$: $\dfrac{3}{11}\cdot\dfrac{2}{10} = \dfrac{6}{110} = \dfrac{3}{55}$.

Answer: \textbf{(E)}.

Choice (A) is $\binom{4}{3}/\binom{12}{3}$, which is the probability that a specified adjuster
gets all three; multiplying by $3$ adjusters gives the correct $0.055$.
\end{solution}

\begin{example}[Solving for the group size]
An insurer will audit $2$ of its $n$ agents, chosen at random. There are exactly $45$
different pairs of agents that could be audited. The insurer instead decides to audit $3$
agents chosen at random. Calculate the probability that a particular agent is among the three
audited.
\par\smallskip\noindent (A) 0.10 \quad (B) 0.20 \quad (C) 0.30 \quad (D) 0.33 \quad (E) 0.45
\end{example}
\begin{solution}
$\binom{n}{2} = \dfrac{n(n-1)}{2} = 45$, so $n(n-1) = 90 = 10\cdot 9$ and $n = 10$
(the other root of $n^2 - n - 90 = 0$ is $n=-9$, rejected).

With $n=10$, the number of $3$-subsets containing the particular agent is
$\binom{9}{2} = 36$ (the agent is in; choose $2$ companions from the other $9$), out of
$\binom{10}{3} = 120$. Probability $36/120 = 0.30$. Equivalently, $k/n = 3/10$.

Answer: \textbf{(C)}.
\end{solution}

\begin{example}[With replacement: every option chosen at least once]
A survey asks $5$ customers to select one of $3$ deductible options. Each customer selects an
option at random, independently, with each option equally likely. Calculate the probability
that every option is selected by at least one customer.
\par\smallskip\noindent (A) 0.383 \quad (B) 0.617 \quad (C) 0.667 \quad (D) 0.753 \quad (E) 0.802
\end{example}
\begin{solution}
Ordered outcomes with replacement: $3^5 = 243$. Let $A_i$ be the event that option $i$ is
selected by nobody. By inclusion--exclusion,
\[
\Prob(A_1\cup A_2\cup A_3) = 3\left(\frac{2}{3}\right)^5 - 3\left(\frac{1}{3}\right)^5 + 0
= \frac{3\cdot 32 - 3\cdot 1}{243} = \frac{93}{243} .
\]
\[
\Prob(\text{every option used}) = 1 - \frac{93}{243} = \frac{150}{243} = 0.6173 .
\]
Answer: \textbf{(B)}.

Direct count: the $5$ customers split over the $3$ options as $3{+}1{+}1$ or $2{+}2{+}1$.
Pattern $3{+}1{+}1$: choose the option used three times ($3$), choose those three customers
($\binom{5}{3} = 10$), assign the remaining two customers to the two remaining options
($2! = 2$): $60$. Pattern $2{+}2{+}1$: choose the option used once ($3$), choose that customer
($5$), split the remaining four customers into two labelled pairs ($\binom{4}{2} = 6$): $90$.
Total $150$, as before. Choice (A) is the probability of the complement.
\end{solution}

\subsection{Traps}

\begin{trap}[Ordered numerator over unordered denominator, or the reverse]
The most common combinatorics error is counting the favourable outcomes as ordered sequences
(products such as $8\cdot 7\cdot 5$) while the denominator is an unordered count
$\binom{N}{n}$, or the reverse. Decide at the start whether the sample space is ordered or
unordered and count both sets the same way. When in doubt, use unordered counts with
$\binom{\cdot}{\cdot}$ throughout.
\end{trap}

\begin{trap}[Double counting in ``at least $k$'' problems]
To count committees with at least $2$ from a group of $5$, it is wrong to ``choose $2$ from
the $5$, then choose the other $3$ from the remaining $10$'': $\binom{5}{2}\binom{10}{3} = 1200$
exceeds the total $\binom{12}{5}=792$, because a committee with three from the group is counted
three times. Either sum the exact cases ($\binom{5}{2}\binom{7}{3} + \binom{5}{3}\binom{7}{2} + \cdots$)
or use the complement.
\end{trap}

\begin{trap}[Forgetting to divide by $k!$]
${}_nP_k$ counts ordered selections; $\binom{n}{k} = {}_nP_k/k!$ counts unordered ones. A
question that asks for a committee, a sample, or a hand wants $\binom{n}{k}$. Writing
$8\cdot 7 = 56$ for the number of pairs of auto policies, when the number is
$\binom{8}{2} = 28$, doubles the numerator.
\end{trap}

\begin{trap}[Identical objects treated as distinct, or the reverse]
The word ``identical'' or ``indistinguishable'' in a question signals a multinomial
coefficient $n!/(n_1!\cdots n_r!)$, not $n!$. Conversely, when policies or claims are described
by type but are physically distinct items, they may be treated as distinct, provided the
denominator uses the same convention. The safe approach for probability questions is to treat
all objects as distinct in both numerator and denominator.
\end{trap}

\begin{trap}[``At least one'' by direct summation]
Summing the cases $1, 2, \dots, n$ for ``at least one'' wastes time and multiplies the chances
of an arithmetic slip. Compute $1 - \Prob(\text{none})$. For ``at least two'', compute
$1 - \Prob(0) - \Prob(1)$.
\end{trap}

\subsection{Practice problems}

\begin{practice}
A claims department has $8$ senior adjusters and $6$ junior adjusters. A review team of $4$ is
selected at random from the $14$. Calculate the probability that the team consists of exactly
$2$ senior and $2$ junior adjusters.
\par\smallskip\noindent (A) 0.21 \quad (B) 0.32 \quad (C) 0.42 \quad (D) 0.51 \quad (E) 0.60
\end{practice}

\begin{practice}
Of $20$ claims submitted in a month, $3$ are fraudulent. An auditor selects $5$ of the $20$
claims at random without replacement. Calculate the probability that at least one fraudulent
claim is selected.
\par\smallskip\noindent (A) 0.40 \quad (B) 0.52 \quad (C) 0.56 \quad (D) 0.60 \quad (E) 0.75
\end{practice}

\begin{practice}
Five auto claim files and three homeowners claim files are placed in a processing queue in
random order. Calculate the probability that the three homeowners files are processed
consecutively.
\par\smallskip\noindent (A) $1/56$ \quad (B) $3/28$ \quad (C) $1/8$ \quad (D) $3/14$ \quad (E) $3/8$
\end{practice}

\begin{practice}
The number of distinct pairs of policies that can be formed from a portfolio of $n$ policies
is $66$. Four policies are selected at random from the portfolio. Calculate the probability
that two specified policies are both among the four selected.
\par\smallskip\noindent (A) $1/66$ \quad (B) $1/33$ \quad (C) $1/22$ \quad (D) $1/12$ \quad (E) $1/11$
\end{practice}

\begin{practice}
Nine claims are divided at random among three adjusters so that each adjuster receives
exactly three claims. Calculate the probability that two particular claims are assigned to
the same adjuster.
\par\smallskip\noindent (A) $1/4$ \quad (B) $2/7$ \quad (C) $1/3$ \quad (D) $4/9$ \quad (E) $1/2$
\end{practice}

\begin{practice}
Four customers each independently select one of $5$ deductible options, with every option
equally likely. Calculate the probability that the four customers select four different
options.
\par\smallskip\noindent (A) 0.096 \quad (B) 0.192 \quad (C) 0.288 \quad (D) 0.384 \quad (E) 0.480
\end{practice}

\begin{practice}
A pricing committee of $4$ is selected at random from $5$ actuaries and $4$ underwriters.
Calculate the probability that the committee includes at least one actuary and at least one
underwriter.
\par\smallskip\noindent (A) 0.833 \quad (B) 0.881 \quad (C) 0.905 \quad (D) 0.952 \quad (E) 0.976
\end{practice}

\subsection{Answers to practice problems}

\begin{enumerate}[label=\arabic*.,nosep]
\item \textbf{(C)}. $\binom{8}{2}\binom{6}{2}/\binom{14}{4} = 28\cdot 15/1001 = 420/1001 = 0.4196$.
\item \textbf{(D)}. $1 - \binom{17}{5}/\binom{20}{5} = 1 - 6188/15504 = 1 - 0.3991 = 0.6009$.
Choice (C), $1-(17/20)^5 = 0.556$, wrongly assumes replacement.
\item \textbf{(B)}. Glue the three homeowners files into one block: $6!$ arrangements of six
blocks, times $3!$ internal orders, over $8!$: $6\cdot 720/40320 = 4320/40320 = 3/28$.
\item \textbf{(E)}. $\binom{n}{2}=66 \Rightarrow n(n-1)=132=12\cdot 11$, $n=12$. Both specified
policies in: $\binom{10}{2}/\binom{12}{4} = 45/495 = 1/11$. Sequentially:
$\frac{4}{12}\cdot\frac{3}{11} = 1/11$.
\item \textbf{(A)}. Wherever the first claim goes, that adjuster has $2$ open slots among the
$8$ remaining claims: $2/8 = 1/4$. Counting: $3\cdot\binom{7}{1}\binom{6}{3}/\dfrac{9!}{3!\,3!\,3!} = 420/1680 = 1/4$.
\item \textbf{(B)}. ${}_5P_4/5^4 = 120/625 = 0.192$.
\item \textbf{(D)}. Complement: all actuaries or all underwriters,
$\left[\binom{5}{4}+\binom{4}{4}\right]/\binom{9}{4} = 6/126$; probability $120/126 = 20/21 = 0.952$.
\end{enumerate}

\subsection{One-page summary}

\begin{itemize}[nosep]
\item Classical probability: $\Prob(A) = |A|/|S|$ for equally likely outcomes. Count numerator
and denominator by the same convention (both ordered or both unordered).
\item Multiplication principle: $n_1 n_2\cdots n_k$ for a $k$-stage procedure.
\item Ordered, no replacement: ${}_nP_k = n!/(n-k)!$; all $n$: $n!$.
\item Unordered, no replacement: $\binom{n}{k} = n!/[k!(n-k)!] = {}_nP_k/k!$.
\item Ordered, with replacement: $n^k$. Unordered, with replacement: $\binom{n+k-1}{k}$.
\item Identities: $\binom{n}{k}=\binom{n}{n-k}$; $\binom{n}{2}=n(n-1)/2$;
$\binom{n}{k}=\binom{n-1}{k-1}+\binom{n-1}{k}$; $\sum_k\binom{n}{k}=2^n$;
$(x+y)^n=\sum_k\binom{n}{k}x^ky^{n-k}$.
\item A random $k$-subset of $n$ contains a specified object with probability $k/n$.
\item Multinomial coefficient $n!/(n_1!\cdots n_r!)$: arrangements with identical objects, and
partitions of $n$ distinct objects into labelled groups of sizes $n_1,\dots,n_r$.
\item Hypergeometric pattern: $\binom{a}{x}\binom{b}{n-x}/\binom{a+b}{n}$; extend with one
factor per type. Equals the sequential product of conditional probabilities times the number
of orderings.
\item Circular arrangements: $(n-1)!$. Stars and bars: $x_1+\cdots+x_r=n$, $x_i\ge0$, has
$\binom{n+r-1}{r-1}$ solutions; $x_i\ge1$ has $\binom{n-1}{r-1}$.
\item Must be adjacent: glue into a block, arrange blocks, multiply by internal orderings.
Must be separated: arrange the others, then choose gaps.
\item ``At least one'': $1-\Prob(\text{none})$. Birthday pattern: $1-{}_nP_k/n^k$.
\item Solve $\binom{n}{2}=m$ by $n(n-1)=2m$ and factoring into consecutive integers.
\item Never ``fix $k$ then choose the rest'' for ``at least $k$''; sum exact cases or use
the complement.
\end{itemize}

\clearpage
\section{Outcome (c): Independence}

\subsection{What the outcome asks}

The exam tests whether you can recognise independence from the wording of a word problem and then use the multiplication rule $\Prob(A \cap B) = \Prob(A)\,\Prob(B)$ without being told to. The most common question shapes are: the probability that all of several independent risks occur (multiply), the probability that at least one of them occurs (one minus the product of the complements), and the probability that exactly one occurs (sum of the disjoint ``one yes, the rest no'' products). A second family gives you $\Prob(A)$, $\Prob(A \cup B)$ and the word ``independent'' and asks you to solve for $\Prob(B)$, or gives both risks the same unknown probability and makes you solve a quadratic. A third family gives a table of joint probabilities and asks whether the events are independent, or asks for the missing cell that makes them independent. Underlying all of these is the fact that if $A$ and $B$ are independent, so are $A$ and $\cc{B}$, which is what lets you convert ``fails'' probabilities into ``works'' probabilities freely.

\subsection{Definitions}

\begin{keyfact}[Independence of two events]
Events $A$ and $B$ are \emph{independent} if and only if
\[
\Prob(A \cap B) = \Prob(A)\,\Prob(B).
\]
If $\Prob(B) > 0$ this is equivalent to $\Prob(A \mid B) = \Prob(A)$, and if $\Prob(A) > 0$ it is equivalent to $\Prob(B \mid A) = \Prob(B)$: knowing that one event happened does not change the probability of the other.
\end{keyfact}

Use the product form when the question gives joint probabilities and the conditional form when it gives conditional probabilities.

\paragraph{Complements of independent events are independent.}
If $A$ and $B$ are independent, then each of the pairs $(A, \cc{B})$, $(\cc{A}, B)$ and $(\cc{A}, \cc{B})$ is also independent. Proof for $(A, \cc{B})$: since $A = (A \cap B) \cup (A \cap \cc{B})$ is a disjoint union,
\[
\Prob(A \cap \cc{B}) = \Prob(A) - \Prob(A \cap B) = \Prob(A) - \Prob(A)\Prob(B) = \Prob(A)\bigl(1 - \Prob(B)\bigr) = \Prob(A)\,\Prob(\cc{B}).
\]
Applying the same step again with the roles swapped gives $(\cc{A}, \cc{B})$, and by symmetry $(\cc{A}, B)$. The practical consequence: if two components fail independently, they also work independently, and one works independently of the other failing.

\begin{keyfact}[Working with complements]
For independent $A$ and $B$:
\begin{align*}
\Prob(\text{both occur}) &= \Prob(A)\Prob(B), \\
\Prob(\text{neither occurs}) &= \Prob(\cc{A} \cap \cc{B}) = \bigl(1 - \Prob(A)\bigr)\bigl(1 - \Prob(B)\bigr), \\
\Prob(\text{at least one occurs}) &= \Prob(A) + \Prob(B) - \Prob(A)\Prob(B) = 1 - \bigl(1-\Prob(A)\bigr)\bigl(1-\Prob(B)\bigr), \\
\Prob(\text{exactly one occurs}) &= \Prob(A)\bigl(1 - \Prob(B)\bigr) + \bigl(1 - \Prob(A)\bigr)\Prob(B).
\end{align*}
\end{keyfact}

\paragraph{Mutual independence of $n$ events.}
Events $A_1, A_2, \dots, A_n$ are \emph{mutually independent} if for \emph{every} subcollection $A_{i_1}, \dots, A_{i_k}$ ($2 \le k \le n$),
\[
\Prob(A_{i_1} \cap \cdots \cap A_{i_k}) = \Prob(A_{i_1}) \cdots \Prob(A_{i_k}).
\]
The events are \emph{pairwise independent} if the product rule holds for every pair only. Mutual independence implies pairwise independence; the converse fails.

Classic counterexample. Toss two fair coins. Let $A$ = first coin is heads, $B$ = second coin is heads, $C$ = exactly one head. Then $\Prob(A) = \Prob(B) = \Prob(C) = \tfrac12$, and
\[
\Prob(A \cap B) = \Prob(A \cap C) = \Prob(B \cap C) = \tfrac14 = \tfrac12 \cdot \tfrac12,
\]
so the three events are pairwise independent. But $A \cap B \cap C = \emptyset$ (two heads is not exactly one head), so $\Prob(A \cap B \cap C) = 0 \ne \tfrac18$. The three events are not mutually independent. When an exam question says ``independent'' about three or more events, it means mutually independent, and every product rule you need is available.

\begin{keyfact}[At least one of $n$ independent events]
If $A_1, \dots, A_n$ are mutually independent with $\Prob(A_i) = p_i$, then
\[
\Prob(\text{none occur}) = \prod_{i=1}^{n} (1 - p_i), \qquad
\Prob(\text{at least one occurs}) = 1 - \prod_{i=1}^{n} (1 - p_i).
\]
When all $p_i = p$: $\Prob(\text{at least one}) = 1 - (1-p)^n$. This is the single most used formula in this outcome.
\end{keyfact}

\paragraph{Independent versus mutually exclusive.}
Two events with positive probability cannot be both independent and mutually exclusive. Proof: if $A$ and $B$ are mutually exclusive then $\Prob(A \cap B) = 0$; if they are also independent then $0 = \Prob(A)\Prob(B)$, which forces $\Prob(A) = 0$ or $\Prob(B) = 0$, contradicting positive probability. Mutually exclusive events are in fact as dependent as possible: if one occurs, the other certainly does not.

\paragraph{Independent trials.}
An experiment is repeated so that the outcome of any trial does not affect any other trial (successive years of a policy, successive tosses, separate policyholders). Then any event defined by the outcomes of trial 1, any event defined by trial 2, and so on, are mutually independent. In particular, if each trial ``succeeds'' with probability $p$ and $q = 1 - p$,
\[
\Prob(\text{first success occurs on trial } k) = q^{k-1} p, \qquad
\Prob(\text{no success in the first } k \text{ trials}) = q^{k}.
\]
The first formula is just the multiplication rule applied to ``fail, fail, \dots, fail, succeed''. You do not need a named distribution to use it.

\subsection{How to recognise independence in a word problem}

\begin{center}
\begin{tabular}{@{}p{0.47\textwidth}p{0.47\textwidth}@{}}
\toprule
\textbf{Wording that signals independence} & \textbf{Wording that signals dependence} \\
\midrule
``independent'', ``independently'', ``mutually independent'' &
``given that'', ``of those who'', ``among policyholders with'' (conditional information) \\[3pt]
``the events are independent of each other'' &
Any conditional probability $\Prob(A \mid B)$ quoted with $\Prob(A \mid B) \ne \Prob(A)$ \\[3pt]
Separate policies, separate policyholders, separate risks ``each with probability'' &
``drawn without replacement'', ``selected from the remaining'' \\[3pt]
Separate years of a policy where each year's claim status is stated to be independent &
One population with overlapping groups (e.g.\ 40\% have auto, 30\% have home, 15\% have both) \\[3pt]
``components fail independently'', ``each component operates independently'' &
Components that share a power source, or a system where one failure causes another \\[3pt]
Repeated tosses, rolls, draws with replacement, independent trials &
Joint probabilities given in a table (independence must be checked, not assumed) \\
\bottomrule
\end{tabular}
\end{center}

Two operating rules. First, if the question does not say ``independent'' and does not describe independent trials, do not multiply marginal probabilities; look for the joint or conditional information that the question gives instead. Second, if the question says ``independent'' about several events, treat that as mutual independence and use every product you need.

\subsection{Worked examples}

\begin{example}[Solve for an unknown claim probability]
An insurer issues two policies to unrelated customers. The probability that a claim is filed on a given policy is the same for both policies, and claims on the two policies are independent. The probability that at least one of the two policies has a claim is $0.64$.

Calculate the probability that a claim is filed on a given policy.

\begin{center}
(A)~$0.20$ \quad (B)~$0.32$ \quad (C)~$0.36$ \quad (D)~$0.40$ \quad (E)~$0.60$
\end{center}
\end{example}
\begin{solution}
Let $p$ be the claim probability on each policy, and let $A$, $B$ be the claim events. Method 1, complement: ``at least one'' is the complement of ``neither'', and by independence
\[
\Prob(\text{neither}) = (1-p)^2 = 1 - 0.64 = 0.36 \quad\Longrightarrow\quad 1 - p = 0.6 \quad\Longrightarrow\quad p = 0.4.
\]
Method 2, inclusion--exclusion with the independence product:
\[
\Prob(A \cup B) = p + p - p^2 = 0.64 \quad\Longrightarrow\quad p^2 - 2p + 0.64 = 0 \quad\Longrightarrow\quad p = \frac{2 \pm \sqrt{4 - 2.56}}{2} = \frac{2 \pm 1.2}{2},
\]
so $p = 0.4$ or $p = 1.6$; discard $1.6$. Answer: \textbf{(D)}.
\end{solution}

\begin{example}[Series--parallel system]
A machine has three components, 1, 2 and 3, which operate independently. The machine works if component 1 works and at least one of components 2 and 3 works. The probabilities that components 1, 2 and 3 work are $0.9$, $0.8$ and $0.7$ respectively.

Calculate the probability that the machine works.

\begin{center}
(A)~$0.504$ \quad (B)~$0.720$ \quad (C)~$0.846$ \quad (D)~$0.940$ \quad (E)~$0.994$
\end{center}
\end{example}
\begin{solution}
Let $W_i$ be the event that component $i$ works. Components 2 and 3 are in parallel, so that block works unless both fail:
\[
\Prob(W_2 \cup W_3) = 1 - \Prob(\cc{W_2} \cap \cc{W_3}) = 1 - (0.2)(0.3) = 0.94.
\]
The complement step uses the fact that $\cc{W_2}$ and $\cc{W_3}$ are independent because $W_2$ and $W_3$ are. Component 1 is in series with that block, and $W_1$ is independent of any event built from $W_2$ and $W_3$, so
\[
\Prob(\text{machine works}) = \Prob(W_1)\,\Prob(W_2 \cup W_3) = (0.9)(0.94) = 0.846.
\]
Answer: \textbf{(C)}. Choice (A) is the probability that all three work, which is the wrong system logic.
\end{solution}

\begin{example}[At least one claim among many policyholders]
An insurer has $8$ policyholders. In a given year each policyholder files a claim with probability $0.05$, independently of the others.

Calculate the probability that at least one of the $8$ policyholders files a claim during the year.

\begin{center}
(A)~$0.264$ \quad (B)~$0.337$ \quad (C)~$0.400$ \quad (D)~$0.663$ \quad (E)~$0.736$
\end{center}
\end{example}
\begin{solution}
The complement of ``at least one claim'' is ``no claims''. By independence,
\[
\Prob(\text{no claims}) = (1 - 0.05)^8 = 0.95^8 = 0.6634, \qquad
\Prob(\text{at least one}) = 1 - 0.6634 = 0.3366.
\]
Answer: \textbf{(B)}. Choice (C) is $8 \times 0.05$, the result of adding the probabilities, which is only an upper bound. Choice (D) is the complement, left unfinished.
\end{solution}

\begin{example}[Independence from a joint-probability table]
For a randomly selected policyholder, let $S$ be the event that the policyholder is a smoker and $C$ the event that the policyholder filed a claim last year. An insurer's records give $\Prob(S \cap C) = 0.12$ and $\Prob(S \cap \cc{C}) = 0.28$. Assume $S$ and $C$ are independent.

Calculate the probability that a randomly selected policyholder is a nonsmoker who did not file a claim.

\begin{center}
(A)~$0.18$ \quad (B)~$0.30$ \quad (C)~$0.42$ \quad (D)~$0.48$ \quad (E)~$0.60$
\end{center}
\end{example}
\begin{solution}
The two given cells make up all of $S$: $\Prob(S) = 0.12 + 0.28 = 0.40$. Independence then gives $\Prob(C)$ from the joint cell:
\[
\Prob(S \cap C) = \Prob(S)\Prob(C) \quad\Longrightarrow\quad \Prob(C) = \frac{0.12}{0.40} = 0.30.
\]
Since $S$ and $C$ are independent, so are $\cc{S}$ and $\cc{C}$:
\[
\Prob(\cc{S} \cap \cc{C}) = \Prob(\cc{S})\Prob(\cc{C}) = (0.60)(0.70) = 0.42.
\]
Answer: \textbf{(C)}. The completed table is
\begin{center}
\begin{tabular}{@{}lccc@{}}
\toprule
 & $C$ & $\cc{C}$ & Total \\
\midrule
$S$        & $0.12$ & $0.28$ & $0.40$ \\
$\cc{S}$   & $0.18$ & $0.42$ & $0.60$ \\
\midrule
Total      & $0.30$ & $0.70$ & $1.00$ \\
\bottomrule
\end{tabular}
\end{center}
and every cell equals its row total times its column total, which is the test for independence when a full table is given. If the exam instead gives you the full table and asks whether the events are independent, check one cell: $0.12 = 0.40 \times 0.30$.
\end{solution}

\begin{example}[Exactly one of two independent claims]
A policyholder has an auto policy and a homeowners policy. In a given year the probability of an auto claim is $0.25$ and the probability of a homeowners claim is $0.10$. The two claim events are independent.

Calculate the probability that the policyholder files a claim on exactly one of the two policies during the year.

\begin{center}
(A)~$0.025$ \quad (B)~$0.300$ \quad (C)~$0.325$ \quad (D)~$0.350$ \quad (E)~$0.675$
\end{center}
\end{example}
\begin{solution}
Let $A$ = auto claim, $H$ = homeowners claim. ``Exactly one'' is the disjoint union $(A \cap \cc{H}) \cup (\cc{A} \cap H)$, and each piece is a product by independence (including complement independence):
\[
\Prob(\text{exactly one}) = (0.25)(0.90) + (0.75)(0.10) = 0.225 + 0.075 = 0.300.
\]
Answer: \textbf{(B)}. Check via the other route: $\Prob(A \cup B) = 0.25 + 0.10 - 0.025 = 0.325$ is ``at least one'' (choice (C)); subtracting ``both'' $= 0.025$ gives $0.300$.
\end{solution}

\begin{example}[First claim in a given year]
A policyholder files at most one claim per year. In each year the probability of filing a claim is $0.2$, independently of all other years.

Calculate the probability that the policyholder's first claim occurs in the fourth year.

\begin{center}
(A)~$0.0819$ \quad (B)~$0.1024$ \quad (C)~$0.4096$ \quad (D)~$0.5120$ \quad (E)~$0.5904$
\end{center}
\end{example}
\begin{solution}
``First claim in year 4'' means no claim in years 1, 2, 3 and a claim in year 4. The four yearly events are independent, so multiply:
\[
\Prob(\text{first claim in year 4}) = (0.8)(0.8)(0.8)(0.2) = 0.8^3 \times 0.2 = 0.512 \times 0.2 = 0.1024.
\]
Answer: \textbf{(B)}. Choice (A) is $0.8^4 \times 0.2$, an off-by-one error (that is the probability the first claim is in year 5). Choice (D) is the probability of no claim in the first three years, which ignores the claim in year 4. Choice (E) is the probability of at least one claim in the first four years.
\end{solution}

\begin{example}[First claim in an even-numbered year]
As in the previous example, a policyholder files a claim in each year with probability $0.2$, independently across years.

Calculate the probability that the first claim occurs in an even-numbered year.

\begin{center}
(A)~$0.160$ \quad (B)~$0.200$ \quad (C)~$0.444$ \quad (D)~$0.500$ \quad (E)~$0.556$
\end{center}
\end{example}
\begin{solution}
Let $q = 0.8$ and $p = 0.2$. The events ``first claim in year $2$'', ``first claim in year $4$'', \dots\ are mutually exclusive, and each is a product by independence:
\[
\Prob(\text{even}) = qp + q^3 p + q^5 p + \cdots = \frac{qp}{1 - q^2} = \frac{qp}{(1-q)(1+q)} = \frac{q}{1+q} = \frac{0.8}{1.8} = 0.4444.
\]
Answer: \textbf{(C)}. Sanity check: an odd year is more likely than an even year because year 1 comes first, and indeed $\Prob(\text{odd}) = 1/(1+q) = 0.5556 > 0.4444$.
\end{solution}

\subsection{Traps}

\begin{trap}[Adding when you should multiply]
``Each of $8$ policyholders files a claim with probability $0.05$, independently; find the probability at least one files'' is not $8 \times 0.05 = 0.40$. Adding gives $\Prob(A_1) + \cdots + \Prob(A_n)$, which counts overlaps multiple times and can exceed $1$. The correct route is $1 - 0.95^8 = 0.3366$. Adding is correct only for mutually exclusive events, and independent events with positive probability are never mutually exclusive.
\end{trap}

\begin{trap}[Treating independent as mutually exclusive]
Independent means $\Prob(A \cap B) = \Prob(A)\Prob(B)$, which is positive whenever both probabilities are. Mutually exclusive means $\Prob(A \cap B) = 0$. They are different, and for events with positive probability they are incompatible. If a question says two risks are independent and you write $\Prob(A \cup B) = \Prob(A) + \Prob(B)$, you have silently assumed they are mutually exclusive; the correct formula is $\Prob(A) + \Prob(B) - \Prob(A)\Prob(B)$.
\end{trap}

\begin{trap}[Assuming independence when it is not stated]
If a question says $40\%$ of policyholders have auto coverage and $30\%$ have homeowners coverage and asks for the probability a policyholder has both, $0.12$ is wrong unless the question says the coverages are independent. Without that word, the question gives $\Prob(A \cap B)$, $\Prob(A \cup B)$ or a conditional probability somewhere, and that is what to use. Overlapping populations, ``given that'', and sampling without replacement all signal dependence.
\end{trap}

\begin{trap}[Pairwise is not mutual]
If you are told only that $\Prob(A \cap B) = \Prob(A)\Prob(B)$, $\Prob(A \cap C) = \Prob(A)\Prob(C)$ and $\Prob(B \cap C) = \Prob(B)\Prob(C)$, you may not conclude $\Prob(A \cap B \cap C) = \Prob(A)\Prob(B)\Prob(C)$. The two-coin example (first head, second head, exactly one head) has all three pairwise products correct and the triple product wrong. Use the triple product only when the question says the events are (mutually) independent, or when the question describes independent trials.
\end{trap}

\begin{trap}[Forgetting that complements inherit independence]
Being told only that the ``works'' events $W_2$, $W_3$ are independent is enough: $\cc{W_2}$ and $\cc{W_3}$ are then independent too, so $\Prob(\cc{W_2} \cap \cc{W_3}) = (1 - 0.8)(1 - 0.7) = 0.06$ directly. The same fact is what makes ``exactly one'' equal to $p_1(1-p_2) + (1-p_1)p_2$ and ``none'' equal to $\prod (1 - p_i)$.
\end{trap}

\subsection{Practice problems}

\begin{practice}
Events $A$ and $B$ are independent, $\Prob(A) = 0.4$ and $\Prob(A \cup B) = 0.7$.

Calculate $\Prob(B)$.

\begin{center}
(A)~$0.30$ \quad (B)~$0.42$ \quad (C)~$0.50$ \quad (D)~$0.58$ \quad (E)~$0.75$
\end{center}
\end{practice}

\begin{practice}
A system has three components that operate independently, each of which works with probability $0.9$. The system works if at least two of the three components work.

Calculate the probability that the system works.

\begin{center}
(A)~$0.729$ \quad (B)~$0.810$ \quad (C)~$0.900$ \quad (D)~$0.972$ \quad (E)~$0.999$
\end{center}
\end{practice}

\begin{practice}
An insurer has $10$ policies in force. Each policy has a claim during the year with probability $0.03$, and the policies are independent.

Calculate the probability that the insurer receives at least one claim during the year.

\begin{center}
(A)~$0.030$ \quad (B)~$0.263$ \quad (C)~$0.300$ \quad (D)~$0.737$ \quad (E)~$0.970$
\end{center}
\end{practice}

\begin{practice}
A policyholder has three coverages. In a given year the probabilities of a claim on the three coverages are $0.1$, $0.2$ and $0.3$, and the three claim events are mutually independent.

Calculate the probability that the policyholder has a claim on exactly one of the three coverages.

\begin{center}
(A)~$0.006$ \quad (B)~$0.398$ \quad (C)~$0.496$ \quad (D)~$0.504$ \quad (E)~$0.600$
\end{center}
\end{practice}

\begin{practice}
Two independent risks each have the same probability $p$ of producing a loss in a year. The probability that at least one of the two risks produces a loss is $0.51$.

Calculate $p$.

\begin{center}
(A)~$0.255$ \quad (B)~$0.300$ \quad (C)~$0.490$ \quad (D)~$0.510$ \quad (E)~$0.700$
\end{center}
\end{practice}

\begin{practice}
A driver has an accident in any given year with probability $0.1$, independently of other years.

Calculate the probability that the driver's first accident occurs in the third year.

\begin{center}
(A)~$0.081$ \quad (B)~$0.090$ \quad (C)~$0.100$ \quad (D)~$0.271$ \quad (E)~$0.810$
\end{center}
\end{practice}

\begin{practice}
Events $A$, $B$ and $C$ each have probability $0.5$ and are pairwise independent, but $\Prob(A \cap B \cap C) = 0$.

Calculate $\Prob(A \cup B \cup C)$.

\begin{center}
(A)~$0.500$ \quad (B)~$0.625$ \quad (C)~$0.750$ \quad (D)~$0.875$ \quad (E)~$1.000$
\end{center}
\end{practice}

\subsection{Answers to practice problems}

\begin{enumerate}[label=\arabic*.,nosep]
\item \textbf{(C)}. $0.7 = 0.4 + b - 0.4b = 0.4 + 0.6b$, so $b = 0.3/0.6 = 0.5$. Equivalently $\Prob(\cc{A} \cap \cc{B}) = 0.3 = 0.6(1-b)$.
\item \textbf{(D)}. $\Prob(\text{exactly two}) + \Prob(\text{all three}) = 3(0.9)^2(0.1) + 0.9^3 = 0.243 + 0.729 = 0.972$. The three ``exactly two'' patterns are disjoint and each is a product by independence.
\item \textbf{(B)}. $1 - 0.97^{10} = 1 - 0.7374 = 0.2626$. Choice (C) is the sum $10 \times 0.03$.
\item \textbf{(B)}. $(0.1)(0.8)(0.7) + (0.9)(0.2)(0.7) + (0.9)(0.8)(0.3) = 0.056 + 0.126 + 0.216 = 0.398$. Choice (C) is ``at least one'', $1 - (0.9)(0.8)(0.7) = 0.496$.
\item \textbf{(B)}. $\Prob(\text{neither}) = (1-p)^2 = 1 - 0.51 = 0.49$, so $1 - p = 0.7$ and $p = 0.3$. Via the quadratic, $2p - p^2 = 0.51$ gives the same root.
\item \textbf{(A)}. No accident in years 1 and 2, accident in year 3: $(0.9)^2(0.1) = 0.081$. Choice (D) is ``first accident within three years'', $1 - 0.9^3$.
\item \textbf{(C)}. Inclusion--exclusion: $3(0.5) - 3(0.25) + 0 = 0.75$. The pairwise products supply the three pair terms; the triple term is given as $0$, not as $0.125$. This is the two-coin example (first head, second head, exactly one head), where $A \cup B \cup C$ is ``at least one head'' $= 3/4$.
\end{enumerate}

\subsection{One-page summary}

\begin{itemize}[nosep]
\item Definition: $A$ and $B$ independent $\iff \Prob(A \cap B) = \Prob(A)\Prob(B)$. Equivalent when $\Prob(B) > 0$: $\Prob(A \mid B) = \Prob(A)$.
\item If $A$, $B$ independent then so are $(A, \cc{B})$, $(\cc{A}, B)$, $(\cc{A}, \cc{B})$. Proof: $\Prob(A \cap \cc{B}) = \Prob(A) - \Prob(A)\Prob(B) = \Prob(A)\Prob(\cc{B})$.
\item Two events with positive probability cannot be both independent and mutually exclusive ($\Prob(A)\Prob(B) > 0 = \Prob(A \cap B)$ is impossible).
\item For independent $A$, $B$: $\Prob(A \cup B) = \Prob(A) + \Prob(B) - \Prob(A)\Prob(B) = 1 - (1-\Prob(A))(1-\Prob(B))$.
\item Exactly one of two: $p_1(1-p_2) + (1-p_1)p_2$. Neither: $(1-p_1)(1-p_2)$.
\item $n$ mutually independent events: $\Prob(\text{none}) = \prod(1-p_i)$; $\Prob(\text{at least one}) = 1 - \prod(1-p_i)$; equal $p$: $1 - (1-p)^n$.
\item Exactly $k$ of $n$ equal-probability independent events: $\binom{n}{k}p^k(1-p)^{n-k}$ (each pattern is a product; the patterns are disjoint).
\item Mutual independence requires the product rule for every subcollection; pairwise independence does not imply it. Counterexample: two coins, $A$ = first H, $B$ = second H, $C$ = exactly one H; $\Prob(A \cap B \cap C) = 0 \ne 1/8$.
\item ``Independent'' about three or more events on the exam means mutually independent.
\item Independent trials, success probability $p$, $q = 1-p$: $\Prob(\text{first success on trial } k) = q^{k-1}p$; $\Prob(\text{no success in } k \text{ trials}) = q^k$; $\Prob(\text{first success in an even trial}) = q/(1+q)$.
\item Series block: multiply the ``works'' probabilities. Parallel block: $1 - \prod(\text{fails})$. Reduce a mixed system block by block.
\item Solving for an unknown: write $\Prob(A \cup B) = a + b - ab$ (linear in one unknown) or $2p - p^2$ / $(1-p)^2$ (quadratic; take the root in $[0,1]$).
\item Table check: events are independent if and only if each cell equals its row total times its column total; checking one cell is enough for a $2 \times 2$ table.
\item Do not multiply marginals unless the problem says ``independent'' or describes independent trials; do not add unless the events are mutually exclusive.
\end{itemize}

\clearpage
\section{Outcome (d): Mutually exclusive events}

\subsection{What the outcome asks}

The exam tests whether you can recognise when two or more events cannot occur together
(mutually exclusive, or disjoint, events) and whether you then add their probabilities
correctly. It also tests the reverse direction: given $\Prob(A)$, $\Prob(B)$ and
$\Prob(A \cup B)$, decide whether the events are disjoint, or solve for whatever
intersection probability makes the numbers consistent. Underneath both directions sits one
technique that drives almost every Topic 1 problem: split the event you want into
non-overlapping pieces, find the probability of each piece, and add. Partitions of the
sample space (a family of disjoint events that together cover everything) are the standard
tool for that split. The phrases ``exactly one'', ``at least one'', ``at most one'' and
``none'' are all evaluated by writing the event as a union of disjoint pieces and adding.

\subsection{Definitions and rules}

Two events $A$ and $B$ are \emph{mutually exclusive} (equivalently \emph{disjoint}) if
$A \cap B = \varnothing$, so they cannot both occur on the same trial. A family
$A_1, A_2, \dots$ is \emph{pairwise disjoint} if $A_i \cap A_j = \varnothing$ for every
$i \neq j$. A family is \emph{exhaustive} if its union is the whole sample space $S$. A
family that is both pairwise disjoint and exhaustive is a \emph{partition} of $S$: every
outcome lies in exactly one of the events.

\begin{keyfact}[Addition rule for disjoint events]
If $A \cap B = \varnothing$ then
\[
\Prob(A \cup B) = \Prob(A) + \Prob(B).
\]
More generally, if $A_1, \dots, A_n$ are pairwise disjoint then
\[
\Prob\!\left(\bigcup_{i=1}^{n} A_i\right) = \sum_{i=1}^{n} \Prob(A_i),
\]
and the same holds for a countably infinite family $A_1, A_2, \dots$ (countable
additivity, one of the probability axioms). Additivity is the \emph{only} rule that lets you
add probabilities directly; it applies exactly when the events are disjoint.
\end{keyfact}

\begin{keyfact}[Complement as a partition]
For any event $A$, the pair $\{A, \cc{A}\}$ is a partition of $S$, so
\[
\Prob(A) + \Prob(\cc{A}) = 1, \qquad \Prob(\cc{A}) = 1 - \Prob(A).
\]
More generally, if $B_1, \dots, B_n$ partition $S$ then
$\Prob(B_1) + \dots + \Prob(B_n) = 1$. This is how the exam hides one probability: it gives
$n-1$ of them and expects you to recover the last one by subtraction.
\end{keyfact}

\begin{keyfact}[Disjoint decomposition identities]
For any events $A$ and $B$ (disjoint or not), the following unions are unions of disjoint
pieces, so their probabilities add:
\begin{align*}
A &= (A \cap B) \cup (A \cap \cc{B}),
 & \Prob(A) &= \Prob(A \cap B) + \Prob(A \cap \cc{B}), \\
A \cup B &= A \cup (B \cap \cc{A}),
 & \Prob(A \cup B) &= \Prob(A) + \Prob(B \cap \cc{A}), \\
A \cup B &= (A \cap \cc{B}) \cup (A \cap B) \cup (\cc{A} \cap B),
 & \Prob(A \cup B) &= \Prob(A \cap \cc{B}) + \Prob(A \cap B) + \Prob(\cc{A} \cap B).
\end{align*}
The first line is the law of total probability in its simplest form: condition on whether
$B$ happens or not. With a partition $B_1, \dots, B_n$ it reads
$\Prob(A) = \sum_{i} \Prob(A \cap B_i)$.
\end{keyfact}

\begin{keyfact}[Inclusion--exclusion collapses when events are disjoint]
For any two events, $\Prob(A \cup B) = \Prob(A) + \Prob(B) - \Prob(A \cap B)$. For three,
\[
\Prob(A \cup B \cup C) = \Prob(A) + \Prob(B) + \Prob(C)
 - \Prob(A \cap B) - \Prob(A \cap C) - \Prob(B \cap C) + \Prob(A \cap B \cap C).
\]
When the events are pairwise disjoint every intersection has probability $0$ and the
formula reduces to plain addition. Conversely, if the given numbers satisfy
$\Prob(A \cup B) = \Prob(A) + \Prob(B)$ then $\Prob(A \cap B) = 0$; if
$\Prob(A \cup B) < \Prob(A) + \Prob(B)$ the events overlap with positive probability.
\end{keyfact}

\begin{keyfact}[Disjoint is not independent]
If $A$ and $B$ are mutually exclusive and $\Prob(A) > 0$, $\Prob(B) > 0$, then $A$ and $B$
are \emph{dependent}.

\emph{Proof.} Disjointness gives $\Prob(A \cap B) = \Prob(\varnothing) = 0$, while
$\Prob(A)\Prob(B) > 0$. Hence $\Prob(A \cap B) \neq \Prob(A)\Prob(B)$, which is the
definition of dependence. Equivalently, $\Prob(A \mid B) = 0 \neq \Prob(A)$: knowing $B$
occurred rules $A$ out completely. \qed
\end{keyfact}

The three adjectives the exam uses for a family of events answer three different questions.

\begin{center}
\begin{tabular}{@{}lll@{}}
\toprule
Term & Definition & Consequence for probabilities \\
\midrule
Mutually exclusive & $A \cap B = \varnothing$ & $\Prob(A \cap B) = 0$; $\Prob(A \cup B) = \Prob(A) + \Prob(B)$ \\
Independent & $\Prob(A \cap B) = \Prob(A)\Prob(B)$ & $\Prob(A \mid B) = \Prob(A)$; can both occur \\
Exhaustive & $A \cup B = S$ & $\Prob(A \cup B) = 1$; at least one always occurs \\
\bottomrule
\end{tabular}
\end{center}

Mutually exclusive is a statement about the \emph{sets} (they do not overlap). Independent
is a statement about the \emph{numbers} (the product rule holds). Exhaustive is a statement
about \emph{coverage} (nothing is left out). Any combination is possible except mutually
exclusive together with independent when both probabilities are positive.

\subsection{Disjoint decomposition as a technique}

The standard move on a Topic 1 word problem is to rewrite the event asked for as a union of
pieces that cannot overlap, then add the probabilities of the pieces. Two things make the
move reliable. First, the pieces are usually intersections of the given events and their
complements, so each piece is a single ``cell'' of the Venn diagram. Second, every cell
probability can be recovered from the quantities the exam typically gives: the marginals
$\Prob(A)$, $\Prob(B)$, $\Prob(C)$, the pairwise intersections and the triple intersection.

\paragraph{Two events.} The four cells $A \cap B$, $A \cap \cc{B}$, $\cc{A} \cap B$ and
$\cc{A} \cap \cc{B}$ have probabilities $\Prob(AB)$, $\Prob(A) - \Prob(AB)$,
$\Prob(B) - \Prob(AB)$ and $1 - \Prob(A) - \Prob(B) + \Prob(AB)$, where $\Prob(AB)$
abbreviates $\Prob(A \cap B)$. ``Exactly one'' is $(A \cap \cc{B}) \cup (\cc{A} \cap B)$,
two disjoint cells, so its probability is $\Prob(A) + \Prob(B) - 2\Prob(AB)$. ``At least
one'' adds the cell $A \cap B$ back in; ``at most one'' removes it from the whole space;
``neither'' is the single cell $\cc{A} \cap \cc{B}$.

\paragraph{Three events.} Let $S_1 = \Prob(A) + \Prob(B) + \Prob(C)$,
$S_2 = \Prob(AB) + \Prob(AC) + \Prob(BC)$ and $S_3 = \Prob(ABC)$. The eight cells of the
three-set Venn diagram are disjoint. ``Exactly two'' is the union of the three cells
$A \cap B \cap \cc{C}$, $A \cap \cc{B} \cap C$ and $\cc{A} \cap B \cap C$; each pairwise
intersection contains the triple intersection, so each cell has probability (pairwise)
$- S_3$, and adding gives $S_2 - 3S_3$. Cell-by-cell reasoning of the same kind gives the
rest of the list below.

\begin{keyfact}[``Exactly'', ``at least'', ``at most'' formulas]
Two events, with $\Prob(A \cap B)$ written $\Prob(AB)$:
\begin{align*}
\Prob(\text{exactly one}) &= \Prob(A) + \Prob(B) - 2\Prob(AB), \\
\Prob(\text{at least one}) &= \Prob(A) + \Prob(B) - \Prob(AB), \\
\Prob(\text{at most one}) &= 1 - \Prob(AB), \\
\Prob(\text{none}) &= 1 - \Prob(A) - \Prob(B) + \Prob(AB).
\end{align*}
Three events, with $S_1 = \Prob(A) + \Prob(B) + \Prob(C)$,
$S_2 = \Prob(AB) + \Prob(AC) + \Prob(BC)$ and $S_3 = \Prob(ABC)$:
\begin{align*}
\Prob(\text{exactly three}) &= S_3, \\
\Prob(\text{exactly two}) &= S_2 - 3S_3, \\
\Prob(\text{exactly one}) &= S_1 - 2S_2 + 3S_3, \\
\Prob(\text{at least one}) &= S_1 - S_2 + S_3, \\
\Prob(\text{at least two}) &= S_2 - 2S_3, \\
\Prob(\text{at most one}) &= 1 - S_2 + 2S_3, \\
\Prob(\text{none}) &= 1 - S_1 + S_2 - S_3.
\end{align*}
Check on any problem: exactly one $+$ exactly two $+$ exactly three $=$ at least one.
\end{keyfact}

When the given events are themselves mutually exclusive, $S_2 = S_3 = 0$ and every formula
above collapses: ``at least one'' and ``exactly one'' coincide and equal $S_1$, ``at most
one'' has probability $1$, and ``none'' has probability $1 - S_1$.

\subsection{Worked examples}

\begin{example}[Claim size bands]
An insurer classifies each claim into exactly one of four size bands. The probability that a
claim is small (under 1{,}000) is $0.42$, medium (from 1{,}000 to under 5{,}000) is $0.28$,
and large (from 5{,}000 to under 25{,}000) is $0.18$. All other claims are catastrophic.
Calculate the probability that a randomly selected claim is either medium or catastrophic.

\smallskip\noindent (A)~$0.40$\hfill (B)~$0.46$\hfill (C)~$0.58$\hfill (D)~$0.34$\hfill (E)~$0.30$\hfill\mbox{}
\end{example}
\begin{solution}
The four bands partition the sample space, so their probabilities sum to $1$:
\[
\Prob(\text{catastrophic}) = 1 - 0.42 - 0.28 - 0.18 = 0.12.
\]
Medium and catastrophic are disjoint bands, so the addition rule applies directly:
\[
\Prob(\text{medium} \cup \text{catastrophic}) = 0.28 + 0.12 = 0.40.
\]
Answer: \textbf{(A)}.
\end{solution}

\begin{example}[Are the events disjoint?]
For a randomly chosen customer of an insurer, the probability of holding an auto policy is
$0.35$, the probability of holding a homeowners policy is $0.45$, and the probability of
holding at least one of the two is $0.70$. Calculate the probability that the customer holds
both policies.

\smallskip\noindent (A)~$0.00$\hfill (B)~$0.20$\hfill (C)~$0.15$\hfill (D)~$0.10$\hfill (E)~$0.05$\hfill\mbox{}
\end{example}
\begin{solution}
Let $A$ be auto and $B$ be homeowners. If $A$ and $B$ were disjoint, $\Prob(A \cup B)$ would
equal $0.35 + 0.45 = 0.80$. The stated value $0.70$ is smaller, so the events overlap and
the general addition rule is required:
\[
\Prob(A \cap B) = \Prob(A) + \Prob(B) - \Prob(A \cup B) = 0.35 + 0.45 - 0.70 = 0.10.
\]
Answer: \textbf{(D)}. The gap between the disjoint-case value $0.80$ and the actual value
$0.70$ is exactly the intersection probability.
\end{solution}

\begin{example}[Exactly one of two]
A policyholder has both an auto policy and a homeowners policy with the same company. In a
given year the probability of a claim on the auto policy is $0.40$, the probability of a claim
on the homeowners policy is $0.25$, and the probability of claims on both is $0.10$.
Calculate the probability that the policyholder files a claim on exactly one of the two
policies.

\smallskip\noindent (A)~$0.65$\hfill (B)~$0.45$\hfill (C)~$0.60$\hfill (D)~$0.35$\hfill (E)~$0.55$\hfill\mbox{}
\end{example}
\begin{solution}
Let $A$ be an auto claim and $H$ a homeowners claim. ``Exactly one'' is the disjoint union
$(A \cap \cc{H}) \cup (\cc{A} \cap H)$. Each piece is a marginal minus the intersection:
\[
\Prob(A \cap \cc{H}) = 0.40 - 0.10 = 0.30, \qquad
\Prob(\cc{A} \cap H) = 0.25 - 0.10 = 0.15.
\]
Adding the disjoint pieces gives $0.30 + 0.15 = 0.45$, which matches the formula
$\Prob(A) + \Prob(H) - 2\Prob(A \cap H) = 0.40 + 0.25 - 0.20 = 0.45$.
Answer: \textbf{(B)}. Choice (E) is $\Prob(A \cup H) = 0.55$, the ``at least one'' answer.
\end{solution}

\begin{example}[Exactly two of three]
A group health plan offers hospitalisation ($A$), dental ($B$) and vision ($C$) riders. For a
randomly selected member, $\Prob(A \cap B) = 0.12$, $\Prob(A \cap C) = 0.09$,
$\Prob(B \cap C) = 0.07$ and $\Prob(A \cap B \cap C) = 0.03$. Calculate the probability that
the member holds exactly two of the three riders.

\smallskip\noindent (A)~$0.25$\hfill (B)~$0.28$\hfill (C)~$0.22$\hfill (D)~$0.16$\hfill (E)~$0.19$\hfill\mbox{}
\end{example}
\begin{solution}
``Exactly two'' is the disjoint union of three cells, each a pairwise intersection with the
triple intersection removed:
\begin{align*}
\Prob(A \cap B \cap \cc{C}) &= 0.12 - 0.03 = 0.09, \\
\Prob(A \cap \cc{B} \cap C) &= 0.09 - 0.03 = 0.06, \\
\Prob(\cc{A} \cap B \cap C) &= 0.07 - 0.03 = 0.04.
\end{align*}
Sum: $0.09 + 0.06 + 0.04 = 0.19$. Equivalently $S_2 - 3S_3 = 0.28 - 0.09 = 0.19$.
Answer: \textbf{(E)}. Choice (B) is $S_2$ itself, which triple-counts the members holding
all three riders.
\end{solution}

\begin{example}[Mutually exclusive with an unknown]
An applicant for life insurance is classified as either preferred or standard, never both,
and some applicants are declined. Let $A$ be the event that an applicant is classified
preferred and $B$ the event that an applicant is classified standard. The probability that an
applicant is accepted (preferred or standard) is $0.60$, and an applicant is twice as likely
to be classified preferred as standard. Calculate $\Prob(B)$.

\smallskip\noindent (A)~$0.40$\hfill (B)~$0.30$\hfill (C)~$0.20$\hfill (D)~$0.15$\hfill (E)~$0.60$\hfill\mbox{}
\end{example}
\begin{solution}
The problem states that $A$ and $B$ cannot both occur, so they are mutually exclusive and
$\Prob(A \cup B) = \Prob(A) + \Prob(B)$. With $\Prob(A) = 2\Prob(B)$,
\[
0.60 = 2\Prob(B) + \Prob(B) = 3\Prob(B) \quad\Longrightarrow\quad \Prob(B) = 0.20,
\]
and $\Prob(A) = 0.40$. Answer: \textbf{(C)}. The probability of being declined is
$1 - 0.60 = 0.40$, which is choice (A); read the question stem carefully.
\end{solution}

\subsection{Traps}

\begin{trap}[Adding without disjointness]
$\Prob(A \cup B) = \Prob(A) + \Prob(B)$ is valid only when $A \cap B = \varnothing$. If the
problem does not state that the events cannot occur together (``exactly one of'', ``never
both'', ``classified into one of the following categories''), use the general rule
$\Prob(A) + \Prob(B) - \Prob(A \cap B)$. A quick sanity check: if the plain sum exceeds $1$,
the events cannot be disjoint. If the sum is at most $1$, the events might still overlap;
only the problem statement or a given intersection settles it.
\end{trap}

\begin{trap}[Disjoint is not independent]
Students frequently treat ``mutually exclusive'' as a synonym for ``independent''. It is
close to the opposite. Disjoint events with positive probability are always dependent,
because $\Prob(A \cap B) = 0$ while $\Prob(A)\Prob(B) > 0$. When a question says ``mutually
exclusive'', set every intersection to $0$; when it says ``independent'', set every
intersection to the product of the marginals. Never do both.
\end{trap}

\begin{trap}[Conditioning on a disjoint event]
If $A$ and $B$ are mutually exclusive and $\Prob(B) > 0$ then
$\Prob(A \mid B) = \Prob(A \cap B)/\Prob(B) = 0$. Likewise
$\Prob(A \cup \cc{B}) = \Prob(\cc{B})$ because $A \subseteq \cc{B}$. Answer choices in such
questions are built from the numbers $\Prob(A)$ and $\Prob(B)$ to catch anyone who computes
$\Prob(A)/\Prob(B)$ or $\Prob(A) + \Prob(\cc{B})$ mechanically.
\end{trap}

\begin{trap}[``Exactly'' versus ``at least'']
``Exactly one'' excludes the intersection; ``at least one'' includes it. For two events the
difference is exactly $\Prob(A \cap B)$, and the exam usually lists both values among the
choices. For three events, ``exactly two'' is $S_2 - 3S_3$, not $S_2$: every pairwise
intersection already contains the triple intersection, so $S_2$ counts the triple cell three
times. ``At least two'' is $S_2 - 2S_3$, which keeps the triple cell once.
\end{trap}

\subsection{Practice problems}

\begin{practice}
Claims on a commercial property policy fall into exactly one of four bands. The probability a
claim is under 500 is $0.35$, from 500 to under 2{,}000 is $0.40$, and 10{,}000 or more is
$0.05$. Calculate the probability that a claim is at least 2{,}000.

\smallskip\noindent (A)~$0.45$\hfill (B)~$0.20$\hfill (C)~$0.05$\hfill (D)~$0.25$\hfill (E)~$0.30$\hfill\mbox{}
\end{practice}

\begin{practice}
For a randomly selected policyholder, the probability of filing a medical claim in a year is
$0.30$, the probability of filing a prescription-drug claim is $0.50$, and the probability of
filing at least one of the two is $0.65$. Calculate the probability that the policyholder
files exactly one of the two types of claim.

\smallskip\noindent (A)~$0.50$\hfill (B)~$0.35$\hfill (C)~$0.65$\hfill (D)~$0.45$\hfill (E)~$0.55$\hfill\mbox{}
\end{practice}

\begin{practice}
Each policy in a portfolio is in exactly one of three risk classes: $A$, $B$ or $C$, or is
unclassified. Let $A$, $B$ and $C$ also denote the events that a randomly chosen policy is in
the respective class. You are given $\Prob(A) = 3\Prob(C)$, $\Prob(B) = 2\Prob(C)$ and
$\Prob(A \cup B \cup C) = 0.90$. Calculate $\Prob(A \cup C)$.

\smallskip\noindent (A)~$0.55$\hfill (B)~$0.45$\hfill (C)~$0.70$\hfill (D)~$0.75$\hfill (E)~$0.60$\hfill\mbox{}
\end{practice}

\begin{practice}
A survey of policyholders finds that $50\%$ carry collision coverage ($A$), $40\%$ carry
comprehensive coverage ($B$) and $30\%$ carry rental reimbursement ($C$). Also,
$20\%$ carry both $A$ and $B$, $15\%$ carry both $A$ and $C$, $10\%$ carry both $B$ and $C$,
and $5\%$ carry all three. Calculate the probability that a randomly selected policyholder
carries exactly two of the three coverages.

\smallskip\noindent (A)~$0.45$\hfill (B)~$0.30$\hfill (C)~$0.35$\hfill (D)~$0.25$\hfill (E)~$0.40$\hfill\mbox{}
\end{practice}

\begin{practice}
For a randomly selected insured, let $A$ be the event of a claim on a first-party coverage
and $B$ the event of a claim on a third-party coverage. You are given
$\Prob(A \cap B) = 0.12$, $\Prob(A \cap \cc{B}) = 0.28$ and $\Prob(\cc{A} \cap B) = 0.18$.
Calculate the probability that the insured has no claim on either coverage.

\smallskip\noindent (A)~$0.60$\hfill (B)~$0.30$\hfill (C)~$0.40$\hfill (D)~$0.42$\hfill (E)~$0.58$\hfill\mbox{}
\end{practice}

\begin{practice}
Events $A$ and $B$ are mutually exclusive with $\Prob(A) = 0.30$ and $\Prob(B) = 0.50$.
Calculate $\Prob(A \cup \cc{B})$.

\smallskip\noindent (A)~$0.50$\hfill (B)~$0.80$\hfill (C)~$0.30$\hfill (D)~$0.20$\hfill (E)~$0.65$\hfill\mbox{}
\end{practice}

\subsection{Answers to practice problems}

\begin{enumerate}[label=\arabic*.,nosep]
\item \textbf{(D)} The bands partition the sample space, so the missing band (2{,}000 to under
      10{,}000) has probability $1 - 0.35 - 0.40 - 0.05 = 0.20$. At least 2{,}000 is the union
      of the two disjoint top bands: $0.20 + 0.05 = 0.25$.
\item \textbf{(A)} $\Prob(A \cap B) = 0.30 + 0.50 - 0.65 = 0.15$. Exactly one is
      $\Prob(A \cup B) - \Prob(A \cap B) = 0.65 - 0.15 = 0.50$, or
      $0.30 + 0.50 - 2(0.15) = 0.50$.
\item \textbf{(E)} The classes are disjoint, so $3c + 2c + c = 6c = 0.90$ and $c = 0.15$.
      Then $\Prob(A \cup C) = 3c + c = 4c = 0.60$.
\item \textbf{(B)} $S_2 = 0.20 + 0.15 + 0.10 = 0.45$, $S_3 = 0.05$, and exactly two is
      $S_2 - 3S_3 = 0.45 - 0.15 = 0.30$. Check: exactly one is
      $1.20 - 0.90 + 0.15 = 0.45$, and $0.45 + 0.30 + 0.05 = 0.80 = S_1 - S_2 + S_3$.
\item \textbf{(D)} The three given cells are disjoint and make up $A \cup B$, so
      $\Prob(A \cup B) = 0.12 + 0.28 + 0.18 = 0.58$ and the answer is $1 - 0.58 = 0.42$.
      (Here $\Prob(A) = 0.40$, $\Prob(B) = 0.30$ and $\Prob(A)\Prob(B) = 0.12 = \Prob(A \cap B)$,
      so these events happen to be independent, and they are not disjoint.)
\item \textbf{(A)} Since $A \cap B = \varnothing$, $A \subseteq \cc{B}$, so
      $A \cup \cc{B} = \cc{B}$ and the probability is $1 - 0.50 = 0.50$. Choice (B) is
      the wrong result $\Prob(A) + \Prob(\cc{B}) = 0.80$, which double counts $A$.
\end{enumerate}

\subsection{One-page summary}

\begin{itemize}[nosep]
\item Mutually exclusive (disjoint): $A \cap B = \varnothing$, so $\Prob(A \cap B) = 0$ and
      $\Prob(A \mid B) = 0$. Pairwise disjoint: no two of the family overlap. Partition:
      pairwise disjoint and exhaustive; probabilities sum to $1$.
\item Addition rule: for disjoint events, $\Prob(A \cup B) = \Prob(A) + \Prob(B)$, and
      $\Prob(\bigcup A_i) = \sum \Prob(A_i)$ for any finite or countable pairwise disjoint
      family. This is the only situation in which probabilities add directly.
\item General rule: $\Prob(A \cup B) = \Prob(A) + \Prob(B) - \Prob(A \cap B)$. If the given
      union equals the sum of the marginals, the events are disjoint; if it is smaller, the
      shortfall is $\Prob(A \cap B)$.
\item Complement: $\Prob(\cc{A}) = 1 - \Prob(A)$. A partition with one probability missing:
      subtract the others from $1$.
\item Decomposition identities: $\Prob(A) = \Prob(A \cap B) + \Prob(A \cap \cc{B})$;
      $\Prob(A \cup B) = \Prob(A) + \Prob(\cc{A} \cap B)$;
      $\Prob(A \cup B) = \Prob(A \cap \cc{B}) + \Prob(A \cap B) + \Prob(\cc{A} \cap B)$.
\item Disjoint with positive probabilities implies dependent: $\Prob(A \cap B) = 0 \neq
      \Prob(A)\Prob(B)$. Mutually exclusive, independent and exhaustive are three unrelated
      properties.
\item Two events: exactly one $= \Prob(A) + \Prob(B) - 2\Prob(AB)$; at least one
      $= \Prob(A) + \Prob(B) - \Prob(AB)$; at most one $= 1 - \Prob(AB)$; none
      $= 1 - \Prob(A) - \Prob(B) + \Prob(AB)$.
\item Three events with $S_1, S_2, S_3$: exactly three $= S_3$; exactly two $= S_2 - 3S_3$;
      exactly one $= S_1 - 2S_2 + 3S_3$; at least one $= S_1 - S_2 + S_3$; at least two
      $= S_2 - 2S_3$; none $= 1 - S_1 + S_2 - S_3$.
\item Technique: rewrite the target event as a union of Venn-diagram cells, compute each cell
      from marginals and intersections, then add. Check that exactly one $+$ exactly two $+$
      exactly three $=$ at least one.
\item If $A$ and $B$ are disjoint then $A \subseteq \cc{B}$, so $\Prob(A \cup \cc{B}) =
      \Prob(\cc{B})$ and $\Prob(A \cap \cc{B}) = \Prob(A)$.
\end{itemize}

\clearpage
\section{Outcome (e): Addition and multiplication rules}

\subsection{What the outcome asks}

The exam tests two rules and the habit of choosing the right one. The addition rule
(inclusion-exclusion) gives the probability that at least one of several events occurs; the exam
states it for two or three events, usually with one quantity missing, and asks for that quantity or
for a derived one such as ``neither'' or ``exactly one''. The multiplication rule (chain rule)
gives the probability that several events occur in sequence, each conditioned on the ones before it;
the exam uses it for draws without replacement and for claim histories over successive years.
Many questions combine the two: split the outcome into disjoint cases with the addition rule, then
compute each case with the multiplication rule. A tree diagram is the bookkeeping device for the
multiplication rule, and its branches are the cases that the addition rule adds up.

\subsection{Addition rule}

\begin{keyfact}[General addition rule, two events]
For any events $A$ and $B$,
\[
\Prob(A \cup B) = \Prob(A) + \Prob(B) - \Prob(A \cap B).
\]
The subtraction removes the double count of outcomes that lie in both events. The rule holds
whether or not $A$ and $B$ are independent or disjoint.
\end{keyfact}

\begin{keyfact}[General addition rule, three events]
For any events $A$, $B$, $C$,
\[
\Prob(A \cup B \cup C) = \Prob(A) + \Prob(B) + \Prob(C)
 - \Prob(A \cap B) - \Prob(A \cap C) - \Prob(B \cap C)
 + \Prob(A \cap B \cap C).
\]
Sign pattern: singles added, pairs subtracted, triple added. Seven terms on the right; the exam
typically gives six of them (or the union) and asks for the seventh.
\end{keyfact}

Special cases reduce the number of terms.

\begin{itemize}[leftmargin=2em]
  \item \textbf{Disjoint (mutually exclusive).} If $A \cap B = \emptyset$ then
  $\Prob(A \cup B) = \Prob(A) + \Prob(B)$. For three pairwise disjoint events the union is the sum
  of the three probabilities.
  \item \textbf{Independent.} If $A$ and $B$ are independent then $\Prob(A \cap B) = \Prob(A)\Prob(B)$, so
  $\Prob(A \cup B) = \Prob(A) + \Prob(B) - \Prob(A)\Prob(B)$.
  \item \textbf{Complement shortcut.} By De Morgan, $\cc{(A \cup B)} = \cc{A} \cap \cc{B}$, so
  $\Prob(A \cup B) = 1 - \Prob(\cc{A} \cap \cc{B})$ and $\Prob(\text{neither}) = 1 - \Prob(A \cup B)$.
  For $n$ events, $\Prob(\text{at least one}) = 1 - \Prob(\text{none})$; when the events are
  independent this is $1 - \prod_{i=1}^{n}(1 - \Prob(A_i))$.
  \item \textbf{Bound.} Without any information about intersections,
  $\Prob(A_1 \cup \cdots \cup A_n) \le \sum_{i=1}^{n} \Prob(A_i)$ (Boole's inequality), with
  equality only when the events are pairwise disjoint.
\end{itemize}

\begin{keyfact}[Exactly one of two events]
\[
\Prob(\text{exactly one of } A, B) = \Prob(A) + \Prob(B) - 2\,\Prob(A \cap B)
 = \Prob(A \cup B) - \Prob(A \cap B).
\]
Given any two of $\Prob(A)+\Prob(B)$, $\Prob(A\cap B)$, $\Prob(A\cup B)$,
$\Prob(\text{exactly one})$, the remaining ones follow.
\end{keyfact}

\begin{keyfact}[Counting identity for three events]
Write $e_1$, $e_2$, $e_3$ for the probabilities that exactly one, exactly two, exactly three of
$A, B, C$ occur. Then
\[
\Prob(A \cup B \cup C) = e_1 + e_2 + e_3, \qquad
\Prob(A) + \Prob(B) + \Prob(C) = e_1 + 2e_2 + 3e_3.
\]
The second identity holds because an outcome in exactly $k$ of the events is counted $k$ times in
the sum of the marginals. The exam uses these two linear equations to recover $e_2$ or $e_3$ from
``at least one'' and ``exactly one'' data.
\end{keyfact}

\paragraph{Solving-for-unknown patterns.} The exam rarely asks for the union directly. The
standard forms are:

\begin{enumerate}[leftmargin=2em]
  \item Given $\Prob(A \cup B)$, $\Prob(A)$, $\Prob(B)$: the intersection is
  $\Prob(A \cap B) = \Prob(A) + \Prob(B) - \Prob(A \cup B)$.
  \item Given $\Prob(\cc{A} \cap \cc{B})$ (``neither''): first $\Prob(A \cup B) = 1 - \Prob(\cc{A} \cap \cc{B})$,
  then proceed as in pattern 1.
  \item Given $\Prob(\text{exactly one})$ and the two marginals: $\Prob(A \cap B) =
  \tfrac12\bigl[\Prob(A) + \Prob(B) - \Prob(\text{exactly one})\bigr]$.
  \item Three events with six of the seven inclusion-exclusion terms and either the union or
  ``none'' given: solve the seven-term equation for the missing term.
  \item Three events with ``at least one'' and ``exactly one'' given together with the sum of the
  marginals: solve the two linear equations in $e_2$ and $e_3$.
\end{enumerate}

\subsection{Multiplication rule}

\begin{keyfact}[General multiplication rule, two events]
For any events $A$ and $B$ with positive probability,
\[
\Prob(A \cap B) = \Prob(A)\,\Prob(B \mid A) = \Prob(B)\,\Prob(A \mid B).
\]
This is the definition of conditional probability rearranged. Use whichever ordering matches
the information given.
\end{keyfact}

\begin{keyfact}[Chain rule, three or more events]
\[
\Prob(A_1 \cap A_2 \cap A_3) = \Prob(A_1)\,\Prob(A_2 \mid A_1)\,\Prob(A_3 \mid A_1 \cap A_2),
\]
and in general
$\Prob(A_1 \cap \cdots \cap A_n) = \Prob(A_1) \prod_{k=2}^{n} \Prob(A_k \mid A_1 \cap \cdots \cap A_{k-1})$.
Each factor is conditioned on everything that happened before it.
\end{keyfact}

\begin{keyfact}[Independence collapses the chain]
If $A_1, \dots, A_n$ are mutually independent, every conditional factor equals the marginal and
$\Prob(A_1 \cap \cdots \cap A_n) = \Prob(A_1)\Prob(A_2)\cdots\Prob(A_n)$.
The product of marginals is a special case of the chain rule, not a separate rule. Multiply
unconditional probabilities only when the problem states independence or the setup implies it
(sampling with replacement, separate policyholders described as independent).
\end{keyfact}

\paragraph{Sequential sampling without replacement.} When items are drawn one at a time without
replacement, the composition of the pool changes after each draw, so the draws are dependent. The
chain rule handles this directly: the $k$th factor is the number of favourable items remaining
divided by the number of items remaining. From a pool of $N$ items of which $m$ are of type $X$,
\[
\Prob(\text{first two both type } X) = \frac{m}{N}\cdot\frac{m-1}{N-1}, \qquad
\Prob(\text{first is } X,\ \text{second is not}) = \frac{m}{N}\cdot\frac{N-m}{N-1}.
\]

\begin{keyfact}[Symmetry of draws without replacement]
The unconditional probability that the $k$th draw is of type $X$ equals the probability that the
first draw is of type $X$, namely $m/N$, for every $k \le N$. Every ordering of the $N$ items is
equally likely, so the item in position $k$ is a uniformly random item from the pool. Once the
results of earlier draws are known, the conditional probability changes.
\end{keyfact}

\paragraph{Tree diagrams.} A tree diagram lists the first event's outcomes as branches, then each
second event's outcomes as sub-branches, and so on. Each branch is labelled with a conditional
probability given the path that leads to it; the probability of a complete path is the product of
its branch labels, which is exactly the chain rule. Because the paths are disjoint, the probability
of any event is the sum of the probabilities of the paths in it, which is the addition rule for
disjoint cases. As a table, the tree for a two-year claim history with
$\Prob(C_1) = 0.20$, $\Prob(C_2 \mid C_1) = 0.35$, $\Prob(C_2 \mid \cc{C_1}) = 0.15$ has one row per path:

\begin{center}
\begin{tabular}{@{}llrr@{}}
\toprule
Year 1 (branch) & Year 2 (branch) & Product & Path probability \\
\midrule
Claim ($0.20$)    & Claim ($0.35$)    & $0.20 \times 0.35$ & $0.070$ \\
Claim ($0.20$)    & No claim ($0.65$) & $0.20 \times 0.65$ & $0.130$ \\
No claim ($0.80$) & Claim ($0.15$)    & $0.80 \times 0.15$ & $0.120$ \\
No claim ($0.80$) & No claim ($0.85$) & $0.80 \times 0.85$ & $0.680$ \\
\midrule
 & & Total & $1.000$ \\
\bottomrule
\end{tabular}
\end{center}

The path probabilities sum to $1.000$, which is the standard check. Any event is a set of rows and
its probability is the sum of those rows.

\subsection{Combining the rules}

The most common exam structure is a two-stage experiment: a classification or a first draw, then a
second outcome whose probability depends on the first. The method is fixed: list the disjoint cases
for the first stage (the branches of the tree), which must cover every possibility without
overlapping; for each case compute $\Prob(\text{case}) \times \Prob(\text{target} \mid \text{case})$
with the multiplication rule; add the results with the addition rule for disjoint events.
Written out for cases $B_1, \dots, B_n$ that partition the sample space and a target event $A$,
\[
\Prob(A) = \sum_{i=1}^{n} \Prob(B_i)\,\Prob(A \mid B_i).
\]
This is the law of total probability. Outcome (g) develops it fully and adds Bayes' theorem for
reversing the conditioning; for this outcome it is enough to recognise that ``split into cases,
multiply along each path, add the paths'' is the same computation. Example: 70\% of policyholders
are low risk with claim probability $0.10$ and 30\% are high risk with claim probability $0.30$;
the probability that a randomly selected policyholder has a claim is
$0.70 \times 0.10 + 0.30 \times 0.30 = 0.07 + 0.09 = 0.16$.

A question that gives conditional probabilities along a chain and asks for ``at least one'' over
the chain is usually fastest through the ``none'' path: multiply along it and subtract from 1.

\subsection{Worked examples}

\begin{example}[Three coverages, solve for the triple intersection]
An insurer's policyholders may hold auto, homeowners, and life coverage. Among them, 60\% hold
auto, 45\% hold homeowners, 30\% hold life, 25\% hold both auto and homeowners, 15\% hold both auto
and life, 10\% hold both homeowners and life, and 5\% hold none of the three.
Calculate the percentage of policyholders who hold all three coverages.
\par\smallskip\noindent (A)~10\% \quad (B)~5\% \quad (C)~8\% \quad (D)~12\% \quad (E)~15\%
\end{example}
\begin{solution}
Let $A$, $H$, $L$ be the events that a policyholder holds auto, homeowners, life. The
complement of the union is ``none'', so $\Prob(A \cup H \cup L) = 1 - 0.05 = 0.95$. Three-event
inclusion-exclusion with $x = \Prob(A \cap H \cap L)$:
\[
0.95 = 0.60 + 0.45 + 0.30 - 0.25 - 0.15 - 0.10 + x = 0.85 + x,
\]
so $x = 0.10$. Answer \textbf{(A)}.
\end{solution}

\begin{example}[Claims in two successive years]
For a policyholder, the probability of at least one claim in year 1 is $0.20$. If the
policyholder has a claim in year 1, the probability of a claim in year 2 is $0.35$; if there is
no claim in year 1, the probability of a claim in year 2 is $0.15$.
Calculate the probability that the policyholder has claims in both years or in neither year.
\par\smallskip\noindent (A)~0.68 \quad (B)~0.70 \quad (C)~0.80 \quad (D)~0.75 \quad (E)~0.93
\end{example}
\begin{solution}
Let $C_1$, $C_2$ be the claim events. By the multiplication rule,
\begin{align*}
\Prob(C_1 \cap C_2) &= \Prob(C_1)\Prob(C_2 \mid C_1) = 0.20 \times 0.35 = 0.07,\\
\Prob(\cc{C_1} \cap \cc{C_2}) &= \Prob(\cc{C_1})\Prob(\cc{C_2} \mid \cc{C_1}) = 0.80 \times 0.85 = 0.68.
\end{align*}
``Both'' and ``neither'' are disjoint, so the addition rule gives $0.07 + 0.68 = 0.75$.
Answer \textbf{(D)}. As a check, ``exactly one year'' is $0.20 \times 0.65 + 0.80 \times 0.15 = 0.25$
and $0.75 + 0.25 = 1$.
\end{solution}

\begin{example}[Two policies of the same type]
A portfolio contains 20 policies: 10 auto, 6 homeowners, and 4 life. An auditor selects two
policies at random without replacement.
Calculate the probability that the two selected policies are of the same type.
\par\smallskip\noindent (A)~0.30 \quad (B)~0.35 \quad (C)~0.33 \quad (D)~0.38 \quad (E)~0.40
\end{example}
\begin{solution}
``Same type'' splits into three disjoint cases: both auto, both homeowners, both life. Each case
is a chain-rule product, and the cases are added:
\[
\frac{10}{20}\cdot\frac{9}{19} + \frac{6}{20}\cdot\frac{5}{19} + \frac{4}{20}\cdot\frac{3}{19}
 = \frac{90 + 30 + 12}{380} = \frac{132}{380} = \frac{33}{95} \approx 0.3474.
\]
Answer \textbf{(B)}. Counting unordered pairs gives the same value, $(45 + 15 + 6)/190$.
\end{solution}

\begin{example}[Type of the second policy drawn]
From the portfolio of the previous example (10 auto, 6 homeowners, 4 life), the auditor selects two
policies at random without replacement.
Calculate the probability that the second policy selected is a life policy.
\par\smallskip\noindent (A)~0.15 \quad (B)~0.16 \quad (C)~0.19 \quad (D)~0.21 \quad (E)~0.20
\end{example}
\begin{solution}
Condition on the first draw. Either the first policy is life (probability $4/20$, leaving 3 life
among 19) or it is not (probability $16/20$, leaving 4 life among 19):
\[
\Prob(\text{second is life}) = \frac{4}{20}\cdot\frac{3}{19} + \frac{16}{20}\cdot\frac{4}{19}
 = \frac{12 + 64}{380} = \frac{76}{380} = \frac{4}{20} = 0.20.
\]
Answer \textbf{(E)}. The result equals the first-draw probability $4/20$. This is not a coincidence:
if nothing is known about the first draw, the second policy is a uniformly random policy from the
20, so its type distribution is the same as the pool's. Algebraically,
$\frac{m}{N}\cdot\frac{m-1}{N-1} + \frac{N-m}{N}\cdot\frac{m}{N-1} = \frac{m(N-1)}{N(N-1)} = \frac{m}{N}$.
Recognising the symmetry saves the computation; the computation serves as a check.
\end{solution}

\begin{example}[Three coverages from ``at least one'' and ``exactly one'']
A survey of an insurer's policyholders finds that 60\% hold auto coverage, 35\% hold homeowners
coverage, and 20\% hold life coverage. Also, 80\% hold at least one of the three coverages and
50\% hold exactly one.
Calculate the percentage of policyholders who hold all three coverages.
\par\smallskip\noindent (A)~0\% \quad (B)~10\% \quad (C)~5\% \quad (D)~15\% \quad (E)~25\%
\end{example}
\begin{solution}
Let $e_1, e_2, e_3$ be the proportions holding exactly one, two, three coverages. ``At least one''
gives $e_1 + e_2 + e_3 = 0.80$ and $e_1 = 0.50$, so $e_2 + e_3 = 0.30$. The sum of the marginals
counts each policyholder once per coverage held:
\[
e_1 + 2e_2 + 3e_3 = 0.60 + 0.35 + 0.20 = 1.15
 \quad\Longrightarrow\quad 2e_2 + 3e_3 = 0.65.
\]
Subtract twice the first equation from the second: $e_3 = 0.65 - 2(0.30) = 0.05$, and then
$e_2 = 0.25$. Answer \textbf{(C)}. The individual pairwise intersections are neither needed nor
determined by the data; only their sum $e_2 + 3e_3 = 0.40$ is.
\end{solution}

\begin{example}[Complement of a union from count data]
Of 200 liability claims closed last year, 90 involved property damage, 60 involved bodily injury,
and 120 involved property damage or bodily injury or both. A claim is selected at random.
Calculate the probability that the selected claim involved neither property damage nor bodily
injury.
\par\smallskip\noindent (A)~0.40 \quad (B)~0.15 \quad (C)~0.25 \quad (D)~0.45 \quad (E)~0.60
\end{example}
\begin{solution}
Work in counts, then convert. The union has 120 claims, so its complement has $200 - 120 = 80$
claims and the probability is $80/200 = 0.40$. Answer \textbf{(A)}. The intersection is
$90 + 60 - 120 = 30$ claims, and ``exactly one'' is $120 - 30 = 90$ claims. The 90 and 60 are
counts, not percentages; treating them as proportions produces a nonsense answer.
\end{solution}

\begin{example}[Chain rule over three years]
A policyholder's claim history is modelled year by year. The probability of no claim in year 1 is
$0.80$. Given no claim in year 1, the probability of no claim in year 2 is $0.85$. Given no claim
in years 1 and 2, the probability of no claim in year 3 is $0.90$.
Calculate the probability that the policyholder has at least one claim during the three years.
\par\smallskip\noindent (A)~0.29 \quad (B)~0.35 \quad (C)~0.45 \quad (D)~0.39 \quad (E)~0.55
\end{example}
\begin{solution}
Let $N_k$ be the event of no claim in year $k$. By the chain rule,
\[
\Prob(N_1 \cap N_2 \cap N_3) = \Prob(N_1)\Prob(N_2 \mid N_1)\Prob(N_3 \mid N_1 \cap N_2)
 = 0.80 \times 0.85 \times 0.90 = 0.612.
\]
At least one claim is the complement: $1 - 0.612 = 0.388$. Answer \textbf{(D)}. The problem gives no
conditional probabilities after a claim, so the paths containing a claim cannot be computed and
the complement route is the only one available.
\end{solution}

\subsection{Traps}

\begin{trap}[Adding without subtracting the intersection]
$\Prob(A \cup B) = \Prob(A) + \Prob(B)$ only when $A$ and $B$ are disjoint. Policyholders can hold
two coverages and a claim can involve two causes. Unless the
question says ``mutually exclusive'' or the events cannot occur together by construction, subtract
$\Prob(A \cap B)$. A quick check: if $\Prob(A) + \Prob(B) > 1$, the events cannot be disjoint.
\end{trap}

\begin{trap}[Multiplying marginals without independence]
$\Prob(A \cap B) = \Prob(A)\Prob(B)$ requires independence. Draws without replacement, successive
years of the same policyholder's claims, and risk class followed by claim status are all dependent.
The correct factor for the second event is the conditional probability given the first. If the
question supplies $\Prob(B \mid A)$, it is telling you the events are not independent.
\end{trap}

\begin{trap}[Confusing the joint with the conditional]
$\Prob(A \cap B)$ is the probability that both happen; $\Prob(B \mid A)$ is the probability of $B$
in the world where $A$ has already happened. ``The probability that a policyholder with a claim
in year 1 has a claim in year 2'' is $\Prob(C_2 \mid C_1)$. ``The probability that a policyholder
has claims in both years'' is $\Prob(C_1 \cap C_2) = \Prob(C_1)\Prob(C_2 \mid C_1)$. The
conditional is never smaller than the joint.
\end{trap}

\begin{trap}[Sign pattern in three-event inclusion-exclusion]
Singles are added, pairwise intersections are subtracted, and the triple intersection is added
back. The common errors are subtracting the triple intersection or omitting it. Write the full
seven-term equation before substituting. Check the result: each pairwise intersection must be at
least the triple intersection, and each marginal at least any pairwise intersection containing it.
\end{trap}

\begin{trap}[Counts, percentages, and the base]
Some questions give raw counts (90 of 200 claims), some give percentages of the whole group, and
some give percentages of a subgroup (``of those with auto coverage, 40\% also hold homeowners'').
The third kind is a conditional probability, not a joint probability, and must be multiplied by
the subgroup's proportion before it enters the addition rule. With counts, work in counts and
divide by the total once at the end.
\end{trap}

\subsection{Practice problems}

\begin{practice}
An insurer's records show that 55\% of policyholders are male, 40\% are over age 50, and 20\% are
males over age 50. A policyholder is selected at random.
Calculate the probability that the selected policyholder is a female who is age 50 or younger.
\par\smallskip\noindent (A)~0.05 \quad (B)~0.25 \quad (C)~0.20 \quad (D)~0.30 \quad (E)~0.45
\end{practice}

\begin{practice}
Among an insurer's policyholders, 50\% hold auto coverage, 40\% hold homeowners coverage, 30\%
hold life coverage, 20\% hold auto and homeowners, 15\% hold auto and life, 10\% hold homeowners
and life, and 5\% hold all three.
Calculate the percentage of policyholders who hold none of the three coverages.
\par\smallskip\noindent (A)~10\% \quad (B)~15\% \quad (C)~25\% \quad (D)~30\% \quad (E)~20\%
\end{practice}

\begin{practice}
A claims adjuster has 12 open files: 5 auto, 4 homeowners, and 3 life. The adjuster selects three
files at random without replacement to review first.
Calculate the probability that none of the three selected files is a life file.
\par\smallskip\noindent (A)~0.30 \quad (B)~0.35 \quad (C)~0.38 \quad (D)~0.42 \quad (E)~0.45
\end{practice}

\begin{practice}
A batch of 25 newly issued policies contains 7 with a data-entry error. A quality reviewer examines
the policies one at a time in random order without replacement.
Calculate the probability that the third policy examined contains an error.
\par\smallskip\noindent (A)~0.28 \quad (B)~0.24 \quad (C)~0.26 \quad (D)~0.30 \quad (E)~0.32
\end{practice}

\begin{practice}
For a policyholder, the probability of a claim in year 1 is $0.15$. In any year, if the
policyholder had a claim in the preceding year, the probability of a claim is $0.40$; if the
policyholder had no claim in the preceding year, the probability of a claim is $0.10$.
Calculate the probability that the policyholder has a claim in exactly one of years 1 and 2.
\par\smallskip\noindent (A)~0.135 \quad (B)~0.150 \quad (C)~0.200 \quad (D)~0.175 \quad (E)~0.235
\end{practice}

\begin{practice}
Among an auto insurer's policyholders, 45\% carry collision coverage, 35\% carry comprehensive
coverage, and 50\% carry exactly one of the two.
Calculate the percentage of policyholders who carry both collision and comprehensive coverage.
\par\smallskip\noindent (A)~5\% \quad (B)~15\% \quad (C)~10\% \quad (D)~20\% \quad (E)~30\%
\end{practice}

\begin{practice}
A life insurer classifies applicants as standard, preferred, or substandard; 60\% of applicants are
standard, 30\% are preferred, and 10\% are substandard. The probability of a claim within the
first policy year is $0.08$ for a standard applicant, $0.03$ for a preferred applicant, and $0.20$
for a substandard applicant.
Calculate the probability that a randomly selected applicant has a claim within the first policy
year.
\par\smallskip\noindent (A)~0.058 \quad (B)~0.065 \quad (C)~0.085 \quad (D)~0.103 \quad (E)~0.077
\end{practice}

\subsection{Answers to practice problems}

\begin{enumerate}[leftmargin=2em]
  \item \textbf{(B)} $0.25$. $\Prob(M \cup O) = 0.55 + 0.40 - 0.20 = 0.75$; female and 50 or younger
  is the complement of the union, $1 - 0.75 = 0.25$.
  \item \textbf{(E)} $20\%$. Union $= 0.50 + 0.40 + 0.30 - 0.20 - 0.15 - 0.10 + 0.05 = 0.80$; none
  $= 1 - 0.80 = 0.20$.
  \item \textbf{(C)} $0.38$. Nine non-life files among 12: $\frac{9}{12}\cdot\frac{8}{11}\cdot\frac{7}{10}
  = \frac{504}{1320} = \frac{21}{55} \approx 0.3818$.
  \item \textbf{(A)} $0.28$. By symmetry of draws without replacement, the third policy is a uniformly
  random policy from the batch, so the probability is $7/25 = 0.28$.
  \item \textbf{(D)} $0.175$. Claim then no claim: $0.15 \times 0.60 = 0.090$; no claim then claim:
  $0.85 \times 0.10 = 0.085$; disjoint, so add: $0.175$.
  \item \textbf{(B)} $15\%$. Exactly one $= \Prob(A) + \Prob(B) - 2\Prob(A \cap B)$, so
  $0.50 = 0.80 - 2x$ and $x = 0.15$.
  \item \textbf{(E)} $0.077$. Split by class and multiply along each path:
  $0.60(0.08) + 0.30(0.03) + 0.10(0.20) = 0.048 + 0.009 + 0.020 = 0.077$.
\end{enumerate}

\subsection{One-page summary}

\begin{itemize}[leftmargin=2em]
  \item Two events: $\Prob(A \cup B) = \Prob(A) + \Prob(B) - \Prob(A \cap B)$. Always subtract
  the intersection unless the events are disjoint.
  \item Three events: add the three singles, subtract the three pairs, add the triple. Seven terms;
  the exam hides one.
  \item Complement of the union is ``neither'' or ``none'':
  $\Prob(\cc{A} \cap \cc{B}) = 1 - \Prob(A \cup B)$. ``At least one'' $= 1 - \Prob(\text{none})$.
  \item Exactly one of two events: $\Prob(A) + \Prob(B) - 2\Prob(A \cap B) = \Prob(A \cup B) - \Prob(A \cap B)$.
  \item Three events with $e_1, e_2, e_3$ (exactly one, two, three): union $= e_1 + e_2 + e_3$ and
  sum of marginals $= e_1 + 2e_2 + 3e_3$.
  \item Multiplication rule: $\Prob(A \cap B) = \Prob(A)\Prob(B \mid A) = \Prob(B)\Prob(A \mid B)$.
  Chain rule: $\Prob(A_1 \cap A_2 \cap A_3) = \Prob(A_1)\Prob(A_2 \mid A_1)\Prob(A_3 \mid A_1 \cap A_2)$;
  each factor conditions on all earlier events.
  \item Under independence every conditional factor becomes the marginal and the chain is a plain
  product of marginals. Do not use the plain product otherwise.
  \item Without replacement: $k$th factor is (favourable items remaining)/(items remaining).
  Unconditionally, the $k$th draw has the same type distribution as the first: $m/N$.
  \item Tree diagram: branch labels are conditional probabilities; a path's probability is the
  product of its labels; an event's probability is the sum of its paths. Paths sum to 1.
  \item Two-stage problems: split into disjoint cases, multiply along each case, add. This is the
  law of total probability, $\Prob(A) = \sum_i \Prob(B_i)\Prob(A \mid B_i)$ (outcome (g)).
  \item Joint versus conditional: $\Prob(A \cap B) = \Prob(A)\Prob(B \mid A) \le \Prob(B \mid A)$.
  A percentage ``of those with $A$'' is a conditional probability. With counts, divide by the
  total once at the end.
\end{itemize}

\clearpage
\section{Outcome (f): Conditional probability}

\subsection{What the outcome asks}

The exam tests whether you can read the phrase ``given that'' (or one of its
disguises) and respond by restricting the sample space to the conditioning event
before you count or divide. Questions supply the information in one of a few
forms: joint and marginal probabilities stated in words, a two-way table of
counts or percentages, a Venn setup with two or three coverages, or a small
population from which items are drawn. You compute $\Prob(A \mid B)$ from that
information, and you apply the complement and addition rules \emph{inside} the
conditional world when the question asks for ``not'' or ``or'' after ``given''.
Questions that give $\Prob(A\mid B)$ and ask for $\Prob(B\mid A)$ across several
classes are outcome (g) (Bayes and total probability); the definition that makes
those work is the subject of this section, and the two-class version is solved
here directly from the definition.

\subsection{Definition and properties}

\begin{keyfact}[Definition of conditional probability]
For events $A$ and $B$ with $\Prob(B) > 0$,
\[
\Prob(A \mid B) \;=\; \frac{\Prob(A \cap B)}{\Prob(B)} .
\]
Read: the probability of $A$ given that $B$ has occurred. The denominator is the
probability of the \emph{conditioning} event, the one named after ``given''.
\end{keyfact}

Interpretation. Conditioning on $B$ throws away every outcome outside $B$ and
treats $B$ as the new sample space. Probabilities of outcomes inside $B$ are then
rescaled by dividing by $\Prob(B)$ so that they again sum to $1$. Only the part
of $A$ that lies inside $B$, namely $A \cap B$, survives the restriction.

Because $\Prob(\,\cdot \mid B)$ is a genuine probability measure on the reduced
sample space $B$, it satisfies all three axioms: $\Prob(A\mid B) \ge 0$,
$\Prob(B \mid B) = 1$, and it is additive over disjoint events. Every rule
derived from the axioms therefore has a conditional version, obtained by
writing ``$\mid B$'' after every event.

\begin{keyfact}[Conditional complement and addition rules]
Holding the conditioning event $B$ fixed:
\begin{align*}
\Prob(\cc{A} \mid B) &= 1 - \Prob(A \mid B), \\
\Prob(A \cup C \mid B) &= \Prob(A \mid B) + \Prob(C \mid B) - \Prob(A \cap C \mid B), \\
\Prob(A \cap \cc{C} \mid B) &= \Prob(A \mid B) - \Prob(A \cap C \mid B).
\end{align*}
These hold because the event \emph{after} the bar is the same throughout.
\end{keyfact}

\begin{keyfact}[What is NOT true]
Changing the event \emph{after} the bar is not a complement operation:
\[
\Prob(A \mid \cc{B}) \;\ne\; 1 - \Prob(A \mid B) \quad \text{in general.}
\]
$\Prob(A\mid B)$ and $\Prob(A\mid \cc{B})$ live on two different reduced sample
spaces and need not be related at all. To get $\Prob(A \mid \cc{B})$ you need
$\Prob(A \cap \cc{B})$ and $\Prob(\cc{B})$, typically via
$\Prob(A \cap \cc{B}) = \Prob(A) - \Prob(A \cap B)$.
\end{keyfact}

\begin{keyfact}[Multiplication rule and chain rule]
Rearranging the definition gives
\[
\Prob(A \cap B) = \Prob(B)\,\Prob(A \mid B) = \Prob(A)\,\Prob(B \mid A),
\]
and for three events
\[
\Prob(A \cap B \cap C) = \Prob(A)\,\Prob(B \mid A)\,\Prob(C \mid A \cap B).
\]
Conversely, conditioning on an intersection is the same definition with a
compound conditioning event:
\[
\Prob(C \mid A \cap B) = \frac{\Prob(A \cap B \cap C)}{\Prob(A \cap B)} .
\]
\end{keyfact}

\begin{keyfact}[Link to independence]
If $\Prob(B) > 0$, then $A$ and $B$ are independent exactly when
$\Prob(A \mid B) = \Prob(A)$, equivalently $\Prob(A \mid B) = \Prob(A \mid \cc{B})$.
Knowing $B$ happened does not change the probability of $A$. If
$\Prob(A \mid B) \ne \Prob(A)$ the events are dependent, and the multiplication
rule must use the conditional probability, not the marginal.
\end{keyfact}

Two consequences. $\Prob(A \mid B)$ and $\Prob(B \mid A)$ share the numerator
$\Prob(A \cap B)$ but have different denominators, so
$\Prob(B \mid A) = \Prob(A \mid B)\,\Prob(B)/\Prob(A)$, which reverses a single
condition whenever $\Prob(A)$ is given. If $A \subseteq B$ then
$\Prob(A \mid B) = \Prob(A)/\Prob(B)$; if $A \cap B = \emptyset$ then
$\Prob(A \mid B) = 0$.

\subsection{Three computational settings}

\subsubsection*{(i) Two-way (contingency) table}

The conditioning event names a row or a column. Restrict to that row or column,
then divide the cell you want by the row or column total. Nothing outside that
row or column is used.

A table of counts. An insurer classifies $200$ life policyholders by sex and
smoking status.

\begin{center}
\begin{tabular}{@{}lrrr@{}}
\toprule
 & Smoker & Nonsmoker & Total \\
\midrule
Male   & 30 & 90  & 120 \\
Female & 24 & 56  & 80  \\
\midrule
Total  & 54 & 146 & 200 \\
\bottomrule
\end{tabular}
\end{center}

\begin{itemize}[nosep]
\item $\Prob(\text{Smoker} \mid \text{Male}) = 30/120 = 0.25$: restrict to the Male row.
\item $\Prob(\text{Smoker} \mid \text{Female}) = 24/80 = 0.30$: restrict to the Female row.
\item $\Prob(\text{Male} \mid \text{Smoker}) = 30/54 = 0.5556$: restrict to the Smoker column.
\item $\Prob(\text{Male} \cap \text{Smoker}) = 30/200 = 0.15$: no restriction; divide by the grand total.
\end{itemize}
Since $\Prob(\text{Smoker} \mid \text{Male}) \ne \Prob(\text{Smoker} \mid \text{Female})$,
sex and smoking status are dependent in this population.

A table of proportions. The same procedure works when the cells are joint
probabilities that sum to $1$; the row and column totals are then the marginal
probabilities.

\begin{center}
\begin{tabular}{@{}lrrr@{}}
\toprule
 & Claim & No claim & Total \\
\midrule
Young & 0.12 & 0.28 & 0.40 \\
Old   & 0.09 & 0.51 & 0.60 \\
\midrule
Total & 0.21 & 0.79 & 1.00 \\
\bottomrule
\end{tabular}
\end{center}

\begin{itemize}[nosep]
\item $\Prob(\text{Claim} \mid \text{Young}) = 0.12/0.40 = 0.30$.
\item $\Prob(\text{Claim} \mid \text{Old}) = 0.09/0.60 = 0.15$.
\item $\Prob(\text{Young} \mid \text{Claim}) = 0.12/0.21 = 0.5714$.
\item $\Prob(\text{No claim} \mid \text{Young}) = 1 - 0.30 = 0.70 = 0.28/0.40$: conditional complement, same row.
\item If the table had several age rows, conditioning on ``age at least 30''
would mean adding the relevant rows first: the denominator is the sum of those
row totals and the numerator is the sum of their Claim cells.
\end{itemize}

Exam problems often give the table indirectly: a marginal (``40\% of
policyholders are young''), a conditional (``30\% of young policyholders filed
a claim''), and an overall rate (``21\% filed a claim''). Rebuild the table:
Young--Claim $= 0.40 \times 0.30 = 0.12$, Old--Claim $= 0.21 - 0.12 = 0.09$,
the rest by subtraction. Every conditional is then a ratio of two entries.

\subsubsection*{(ii) Venn diagram / region bookkeeping}

With two or three overlapping coverages, label each disjoint region with its
probability (or count), then
\[
\Prob(A \mid B) = \frac{\text{total of the regions inside } A \cap B}{\text{total of the regions inside } B}.
\]
For three sets, the ``only'' regions come from inclusion--exclusion:
$\Prob(A \text{ only}) = \Prob(A) - \Prob(A\cap B) - \Prob(A \cap C) + \Prob(A \cap B \cap C)$.
``Of those with auto, what fraction also have home'' means the denominator is
every region inside the auto circle and the numerator is the lens $A \cap H$
(both of its regions, with and without life).

\subsubsection*{(iii) Sequential / tree}

When the experiment happens in stages (select a class, then observe a claim;
draw one file, then another), the branch labels on a probability tree are
conditional probabilities given the path so far. Multiplying along a path gives
the joint probability of the path. ``Given the first was $X$, what is the
probability the second is $Y$'' is read directly off a branch label; ``given the
second was $Y$'' reverses the tree, which is outcome (g).

\subsubsection*{Phrases that signal conditioning}

\begin{center}
\begin{tabular}{@{}p{0.42\textwidth}p{0.50\textwidth}@{}}
\toprule
Phrase in the problem & Meaning \\
\midrule
``given that $B$'' / ``if it is known that $B$'' & Denominator is $\Prob(B)$. \\
``of those who $B$'' / ``among $B$ policyholders'' & The $B$ group is the whole population; take the fraction inside it. \\
``restricted to $B$'' / ``within the $B$ class'' & Same. \\
``$x$\% of $B$ policyholders are $A$'' & $\Prob(A \mid B) = x/100$, not $\Prob(A \cap B)$. \\
``$x$\% of policyholders are both $A$ and $B$'' & Joint $\Prob(A \cap B)$; no conditioning. \\
``the probability that a $B$ policyholder is $A$'' & $\Prob(A \mid B)$. \\
\bottomrule
\end{tabular}
\end{center}

\subsection{Worked examples}

\begin{example}[Two-way table of policyholders]
An auto insurer has $500$ policyholders, classified by age and sex as follows.

\begin{center}
\begin{tabular}{@{}lrrr@{}}
\toprule
 & Male & Female & Total \\
\midrule
Young & 140 & 85  & 225 \\
Old   & 160 & 115 & 275 \\
\midrule
Total & 300 & 200 & 500 \\
\bottomrule
\end{tabular}
\end{center}

A policyholder is selected at random. Calculate the probability that the
policyholder is female, given that the policyholder is young.

\medskip
(A) 0.17 \quad (B) 0.38 \quad (C) 0.40 \quad (D) 0.43 \quad (E) 0.45
\end{example}

\begin{solution}
The conditioning event is ``young'', so restrict to the Young row, which has
$225$ policyholders. Of these, $85$ are female.
\[
\Prob(\text{Female} \mid \text{Young}) = \frac{85}{225} = 0.3778 .
\]
The distractors are the other ratios one could form from the same cell: $85/500 = 0.17$ is
$\Prob(\text{Female} \cap \text{Young})$; $200/500 = 0.40$ is the marginal $\Prob(\text{Female})$;
$85/200 = 0.425$ is the reversed conditional $\Prob(\text{Young} \mid \text{Female})$.
\textbf{Answer: (B).}
\end{solution}

\begin{example}[Both, given at least one]
For a randomly selected policyholder holding both auto and homeowners coverage
with an insurer, the probability of an auto claim in a given year is $0.30$, the
probability of a homeowners claim is $0.20$, and the probability of both is
$0.08$. Given that the policyholder files at least one claim in the year,
calculate the probability that the policyholder files both.

\medskip
(A) 0.08 \quad (B) 0.16 \quad (C) 0.19 \quad (D) 0.27 \quad (E) 0.42
\end{example}

\begin{solution}
Let $A$ = auto claim and $H$ = homeowners claim. ``At least one'' is $A \cup H$,
and ``both'' is $A \cap H$, which is a subset of $A \cup H$, so the intersection
of the two events is just $A \cap H$.
\[
\Prob(A \cup H) = 0.30 + 0.20 - 0.08 = 0.42, \qquad
\Prob(A \cap H \mid A \cup H) = \frac{0.08}{0.42} = 0.1905 .
\]
\textbf{Answer: (C).}
\end{solution}

\begin{example}[Venn with three coverages]
A survey of an insurer's customers finds that $55\%$ have auto coverage, $35\%$
have homeowners coverage, and $25\%$ have life coverage. Also, $15\%$ have both
auto and homeowners, $12\%$ have both auto and life, $8\%$ have both homeowners
and life, and $5\%$ have all three. Of those customers who have auto coverage,
calculate the proportion who have neither homeowners nor life coverage.

\medskip
(A) 0.33 \quad (B) 0.40 \quad (C) 0.45 \quad (D) 0.55 \quad (E) 0.60
\end{example}

\begin{solution}
Let $A$, $H$, $L$ be the three coverages. The denominator is $\Prob(A) = 0.55$.
The numerator is the ``auto only'' region:
\[
\Prob(A \cap \cc{H} \cap \cc{L}) = \Prob(A) - \Prob(A\cap H) - \Prob(A \cap L) + \Prob(A\cap H \cap L)
= 0.55 - 0.15 - 0.12 + 0.05 = 0.33 .
\]
Then
\[
\Prob(\cc{H} \cap \cc{L} \mid A) = \frac{0.33}{0.55} = 0.60 .
\]
Choice (A) is the joint probability $0.33$, not the conditional.
\textbf{Answer: (E).}
\end{solution}

\begin{example}[Reversing one condition from the definition]
In a block of life policies, $30\%$ of policyholders are smokers. Among
smokers, $20\%$ filed a claim last year. Overall, $12\%$ of policyholders filed
a claim last year. Calculate the probability that a policyholder who filed a
claim last year is a smoker.

\medskip
(A) 0.06 \quad (B) 0.20 \quad (C) 0.30 \quad (D) 0.50 \quad (E) 0.60
\end{example}

\begin{solution}
Let $S$ = smoker and $C$ = filed a claim. Given: $\Prob(S) = 0.30$,
$\Prob(C \mid S) = 0.20$, $\Prob(C) = 0.12$. The joint probability is
\[
\Prob(S \cap C) = \Prob(S)\,\Prob(C \mid S) = 0.30 \times 0.20 = 0.06 ,
\]
and the requested conditional is
\[
\Prob(S \mid C) = \frac{\Prob(S \cap C)}{\Prob(C)} = \frac{0.06}{0.12} = 0.50 .
\]
No Bayes machinery is needed because $\Prob(C)$ is supplied directly; the
definition does the whole job.
\textbf{Answer: (D).}
\end{solution}

\begin{example}[Conditioning on the complement]
A life insurer classifies $60\%$ of its policyholders as standard and the rest
as preferred. The probability that a standard policyholder lapses within one
year is $0.15$. The probability that a randomly selected policyholder lapses
within one year is $0.11$. Calculate the probability that a preferred
policyholder lapses within one year.

\medskip
(A) 0.02 \quad (B) 0.05 \quad (C) 0.085 \quad (D) 0.15 \quad (E) 0.85
\end{example}

\begin{solution}
Let $S$ = standard and $L$ = lapse. Given: $\Prob(S) = 0.60$, so
$\Prob(\cc{S}) = 0.40$; $\Prob(L \mid S) = 0.15$; $\Prob(L) = 0.11$.
Build the pieces of $L$:
\[
\Prob(L \cap S) = 0.60 \times 0.15 = 0.09, \qquad
\Prob(L \cap \cc{S}) = \Prob(L) - \Prob(L \cap S) = 0.11 - 0.09 = 0.02 .
\]
Then
\[
\Prob(L \mid \cc{S}) = \frac{0.02}{0.40} = 0.05 .
\]
Choice (E), $0.85 = 1 - 0.15$, is $\Prob(\cc{L} \mid S)$, the probability that
a \emph{standard} policyholder does \emph{not} lapse; it says nothing about
preferred policyholders. Choice (A) is the joint probability, not the
conditional.
\textbf{Answer: (B).}
\end{solution}

\begin{example}[Conditioning on an intersection]
An insurer's records show that $50\%$ of its customers have auto coverage,
$40\%$ have homeowners coverage, and $20\%$ have both. Also, $12\%$ of all
customers filed a claim last year, $8\%$ have auto coverage and filed a claim
last year, and $5\%$ have both coverages and filed a claim last year. Calculate
the probability that a customer who has both auto and homeowners coverage filed
a claim last year.

\medskip
(A) 0.05 \quad (B) 0.10 \quad (C) 0.16 \quad (D) 0.25 \quad (E) 0.40
\end{example}

\begin{solution}
Let $A$ = auto, $H$ = homeowners, $C$ = claim. The conditioning event is
$A \cap H$ with $\Prob(A \cap H) = 0.20$; the numerator is $\Prob(A \cap H \cap C) = 0.05$.
\[
\Prob(C \mid A \cap H) = \frac{0.05}{0.20} = 0.25 .
\]
The other figures are distractors: $\Prob(C \mid A) = 0.08/0.50 = 0.16$ conditions
on auto alone; $\Prob(C) = 0.12$ and $\Prob(H) = 0.40$ are not needed.
\textbf{Answer: (D).}
\end{solution}

\subsection{Traps}

\begin{trap}[Dividing by the wrong marginal]
The denominator is always the probability of the event named after ``given''
(or ``of those who'', ``among''). In a table, that means the total of the row
or column you restricted to, never the grand total and never the other
margin. Before dividing, say out loud which group is the new population.
\end{trap}

\begin{trap}[$\Prob(A \mid B)$ versus $\Prob(B \mid A)$]
These share the numerator $\Prob(A \cap B)$ but have different denominators
and are generally different numbers. ``$20\%$ of smokers filed a claim'' is
$\Prob(C \mid S)$; ``$50\%$ of claimants are smokers'' is $\Prob(S \mid C)$.
Exam answer choices routinely include the reversed conditional as a
distractor. Write the conditional with the bar before computing anything.
\end{trap}

\begin{trap}[$\Prob(A \mid \cc{B})$ is not $1 - \Prob(A \mid B)$]
The complement rule flips the event \emph{before} the bar:
$\Prob(\cc{A} \mid B) = 1 - \Prob(A \mid B)$. Flipping the event \emph{after} the
bar changes the population, and the two conditionals need not add to $1$.
To find $\Prob(A \mid \cc{B})$, compute $\Prob(A \cap \cc{B}) = \Prob(A) - \Prob(A \cap B)$
and divide by $\Prob(\cc{B}) = 1 - \Prob(B)$.
\end{trap}

\begin{trap}[Joint versus conditional in the wording]
``$x$\% of policyholders are young and filed a claim'' is a joint probability
$\Prob(Y \cap C)$. ``$x$\% of young policyholders filed a claim'' is a
conditional $\Prob(C \mid Y)$. ``$x$\% of policyholders are young'' is a
marginal $\Prob(Y)$. Misreading one of these as another produces an answer
that is present among the choices. When the problem mixes the three, rebuild the
two-way table from them before answering.
\end{trap}

\begin{trap}[Mixing counts and percentages]
A conditional probability is a ratio, so it can be computed from counts or from
proportions, but not from a mixture. If a row total is a count and a cell is a
percentage of the whole population, convert both to the same units first. In
a proportion table, the grand total is $1$, and the marginals are row and
column sums; do not divide a proportion by a count.
\end{trap}

\subsection{Practice problems}

\begin{practice}
A life insurer classifies $400$ policyholders as preferred or standard, and as
smokers or nonsmokers. There are $20$ preferred smokers, $140$ preferred
nonsmokers, $90$ standard smokers, and $150$ standard nonsmokers. A
policyholder is selected at random. Calculate the probability that the
policyholder is preferred, given that the policyholder is a nonsmoker.

\medskip
(A) 0.35 \quad (B) 0.40 \quad (C) 0.48 \quad (D) 0.72 \quad (E) 0.88
\end{practice}

\begin{practice}
For a randomly selected employee with a group benefits plan, the probability of
a dental claim in a year is $0.45$, the probability of a vision claim is
$0.35$, and the probability of both is $0.15$. Given that the employee has at
least one of the two types of claim in the year, calculate the probability that
the employee has both.

\medskip
(A) 0.15 \quad (B) 0.23 \quad (C) 0.33 \quad (D) 0.43 \quad (E) 0.65
\end{practice}

\begin{practice}
In a survey of an insurer's customers, $40\%$ have auto coverage, $30\%$ have
homeowners coverage, and $25\%$ have life coverage; $12\%$ have auto and
homeowners, $10\%$ have auto and life, $8\%$ have homeowners and life, and
$4\%$ have all three. Of those customers with homeowners coverage, calculate the
proportion who have neither auto nor life coverage.

\medskip
(A) 0.14 \quad (B) 0.40 \quad (C) 0.47 \quad (D) 0.53 \quad (E) 0.60
\end{practice}

\begin{practice}
An auto insurer classifies $20\%$ of its policies as high risk. The probability
that a high-risk policy has a claim in a year is $0.35$. The probability that a
randomly selected policy has a claim in a year is $0.16$. Calculate the
probability that a policy with a claim in the year is a high-risk policy.

\medskip
(A) 0.07 \quad (B) 0.20 \quad (C) 0.35 \quad (D) 0.44 \quad (E) 0.56
\end{practice}

\begin{practice}
Of an insurer's policyholders, $55\%$ are female. The probability that a
female policyholder files a claim in a year is $0.10$. The probability that a
randomly selected policyholder files a claim in a year is $0.145$. Calculate
the probability that a male policyholder files a claim in the year.

\medskip
(A) 0.09 \quad (B) 0.145 \quad (C) 0.20 \quad (D) 0.30 \quad (E) 0.90
\end{practice}

\begin{practice}
An insurer's records show that $60\%$ of customers have auto coverage. Among
customers with auto coverage, $50\%$ carry collision coverage. Also, $12\%$ of
all customers have auto coverage with collision and filed a claim last year.
Calculate the probability that a customer with auto coverage and collision
coverage filed a claim last year.

\medskip
(A) 0.12 \quad (B) 0.20 \quad (C) 0.30 \quad (D) 0.40 \quad (E) 0.72
\end{practice}

\begin{practice}
A file cabinet contains $12$ claim files, of which $5$ are fraudulent. An
auditor selects $2$ files at random without replacement. Given that at least
one of the selected files is fraudulent, calculate the probability that both
are fraudulent.

\medskip
(A) 0.15 \quad (B) 0.22 \quad (C) 0.36 \quad (D) 0.41 \quad (E) 0.68
\end{practice}

\subsection{Answers to practice problems}

\begin{enumerate}[label=\arabic*., nosep]
\item \textbf{(C).} Nonsmokers total $140 + 150 = 290$; of these $140$ are preferred.
$140/290 = 0.4828$. ($140/400 = 0.35$ is the joint; $160/400 = 0.40$ is the marginal;
$140/160 = 0.875$ is the reversed conditional.)
\item \textbf{(B).} $\Prob(D \cup V) = 0.45 + 0.35 - 0.15 = 0.65$;
$\Prob(D \cap V \mid D \cup V) = 0.15/0.65 = 0.2308$.
\item \textbf{(C).} Homeowners only $= 0.30 - 0.12 - 0.08 + 0.04 = 0.14$;
divide by $\Prob(H) = 0.30$: $0.14/0.30 = 0.4667$. ($0.12/0.30 = 0.40$ is the
fraction of homeowners customers with auto.)
\item \textbf{(D).} $\Prob(R \cap C) = 0.20 \times 0.35 = 0.07$;
$\Prob(R \mid C) = 0.07/0.16 = 0.4375$.
\item \textbf{(C).} $\Prob(F \cap C) = 0.55 \times 0.10 = 0.055$;
$\Prob(M \cap C) = 0.145 - 0.055 = 0.09$; $\Prob(C \mid M) = 0.09/0.45 = 0.20$.
(Choice (E), $0.90 = 1 - 0.10$, is $\Prob(\cc{C} \mid F)$, not a statement about males.)
\item \textbf{(D).} $\Prob(A \cap K) = 0.60 \times 0.50 = 0.30$;
$\Prob(C \mid A \cap K) = 0.12/0.30 = 0.40$.
\item \textbf{(B).} $\Prob(\text{both}) = \binom{5}{2}/\binom{12}{2} = 10/66$;
$\Prob(\text{none}) = \binom{7}{2}/\binom{12}{2} = 21/66$;
$\Prob(\text{at least one}) = 45/66$;
ratio $= 10/45 = 2/9 = 0.2222$. ($4/11 = 0.36$ is $\Prob(\text{second fraudulent} \mid \text{first fraudulent})$, a different condition.)
\end{enumerate}

\subsection{One-page summary}

\begin{itemize}[nosep]
\item Definition: $\Prob(A \mid B) = \Prob(A \cap B)/\Prob(B)$, $\Prob(B) > 0$. The
event after the bar is the new sample space; divide by its probability.
\item Signal words for conditioning: given, of those who, among, if it is known
that, restricted to, within the class, ``$x$\% of $B$ policyholders are $A$''.
\item ``$x$\% of policyholders are $A$ and $B$'' is joint; ``$x$\% of $B$
policyholders are $A$'' is conditional; ``$x$\% of policyholders are $B$'' is marginal.
\item Two-way table: restrict to the row or column named after ``given'', divide
the cell by that row or column total. Works for counts or proportions, not a mix.
\item Rebuild the table from words: cell $= \text{marginal} \times \text{conditional}$;
remaining cells by subtraction from marginals.
\item Venn: $\Prob(A \mid B) = (\text{regions in } A \cap B)/(\text{regions in } B)$;
``only'' regions by inclusion--exclusion.
\item Tree: branch labels are conditionals; multiply along a path for the joint.
\item Conditional complement: $\Prob(\cc{A} \mid B) = 1 - \Prob(A \mid B)$.
\item Conditional addition: $\Prob(A \cup C \mid B) = \Prob(A \mid B) + \Prob(C \mid B) - \Prob(A \cap C \mid B)$.
\item Not a rule: $\Prob(A \mid \cc{B}) \ne 1 - \Prob(A \mid B)$. Get it from
$\Prob(A \cap \cc{B}) = \Prob(A) - \Prob(A \cap B)$ divided by $1 - \Prob(B)$.
\item Multiplication: $\Prob(A \cap B) = \Prob(B)\Prob(A \mid B) = \Prob(A)\Prob(B \mid A)$;
chain: $\Prob(A \cap B \cap C) = \Prob(A)\Prob(B \mid A)\Prob(C \mid A \cap B)$.
\item Conditioning on an intersection: $\Prob(C \mid A \cap B) = \Prob(A \cap B \cap C)/\Prob(A \cap B)$.
\item Reversing one condition: $\Prob(B \mid A) = \Prob(A \mid B)\Prob(B)/\Prob(A)$ when $\Prob(A)$ is given.
\item Both given at least one: $\Prob(A \cap B)/\Prob(A \cup B)$ with
$\Prob(A \cup B) = \Prob(A) + \Prob(B) - \Prob(A \cap B)$.
\item Independence: $\Prob(A \mid B) = \Prob(A) = \Prob(A \mid \cc{B})$.
\item Subset: $A \subseteq B \Rightarrow \Prob(A \mid B) = \Prob(A)/\Prob(B)$;
disjoint: $\Prob(A \mid B) = 0$.
\end{itemize}

\clearpage
\section{Outcome (g): Law of total probability and Bayes' theorem}

\subsection{What the outcome asks}

The exam partitions a population into classes (risk classes of drivers, smokers and
non-smokers, suppliers of a part, diseased and healthy patients) and gives the proportion of
the population in each class. These proportions are the \emph{prior} probabilities. Each
class has its own conditional rate of some outcome (an accident, a death, a defect, a positive
test), and the question runs in one of two directions. The forward direction is the law of
total probability: combine the class proportions and the class rates to find the overall rate
of the outcome. The reverse direction is Bayes' theorem: given that the outcome occurred,
find the probability that the individual came from a particular class. Almost every question
on this outcome adds one twist: two years of experience that are independent only within a
class, a class proportion left as an unknown to be solved from a given overall rate, or
conditioning on the outcome \emph{not} occurring.

\subsection{Law of total probability}

Events $B_1, B_2, \dots, B_n$ form a \emph{partition} of the sample space if they are
mutually exclusive ($B_i \cap B_j = \emptyset$ for $i \neq j$) and exhaustive
($B_1 \cup \dots \cup B_n = S$). Every outcome lies in exactly one $B_i$. The simplest
partition is $\{B, \cc{B}\}$.

\begin{keyfact}[Law of total probability]
If $B_1, \dots, B_n$ partition the sample space and $\Prob(B_i) > 0$ for every $i$, then for
any event $A$,
\[
\Prob(A) \;=\; \sum_{i=1}^{n} \Prob(A \mid B_i)\,\Prob(B_i)
\;=\; \Prob(A \mid B_1)\Prob(B_1) + \dots + \Prob(A \mid B_n)\Prob(B_n).
\]
Two-event case:
\[
\Prob(A) \;=\; \Prob(A \mid B)\,\Prob(B) + \Prob(A \mid \cc{B})\,\Prob(\cc{B}).
\]
In words: the overall rate of $A$ is the weighted average of the class rates of $A$, the
weights being the class proportions.
\end{keyfact}

\textbf{Proof.} Because the $B_i$ partition $S$, the event $A$ decomposes into the disjoint
pieces $A \cap B_1, \dots, A \cap B_n$, so $\Prob(A) = \sum_i \Prob(A \cap B_i)$ by
finite additivity. By the multiplication rule, $\Prob(A \cap B_i) = \Prob(A \mid B_i)\Prob(B_i)$.
Substituting gives the formula. \qed

\textbf{Tree view.} The formula is a two-level tree. The first level branches on the class,
with branch probabilities $\Prob(B_i)$; the second level branches on whether $A$ occurs,
with branch probabilities $\Prob(A \mid B_i)$ and $\Prob(\cc{A} \mid B_i) = 1 - \Prob(A \mid B_i)$.
Multiply along each path to get a joint probability, then add the paths that end in $A$.
\begin{itemize}[nosep,leftmargin=2em]
  \item $B_1$ with probability $\Prob(B_1)$
  \begin{itemize}[nosep]
    \item $A$: path probability $\Prob(B_1)\Prob(A \mid B_1)$
    \item $\cc{A}$: path probability $\Prob(B_1)\bigl(1 - \Prob(A \mid B_1)\bigr)$
  \end{itemize}
  \item $B_2$ with probability $\Prob(B_2)$
  \begin{itemize}[nosep]
    \item $A$: path probability $\Prob(B_2)\Prob(A \mid B_2)$
    \item $\cc{A}$: path probability $\Prob(B_2)\bigl(1 - \Prob(A \mid B_2)\bigr)$
  \end{itemize}
  \item $\vdots$
  \item $B_n$ with probability $\Prob(B_n)$
  \begin{itemize}[nosep]
    \item $A$: path probability $\Prob(B_n)\Prob(A \mid B_n)$
    \item $\cc{A}$: path probability $\Prob(B_n)\bigl(1 - \Prob(A \mid B_n)\bigr)$
  \end{itemize}
\end{itemize}
The $2n$ path probabilities sum to 1. $\Prob(A)$ is the sum of the $n$ paths ending in $A$;
$\Prob(\cc{A})$ is the sum of the $n$ paths ending in $\cc{A}$.

\subsection{Bayes' theorem}

Bayes' theorem reverses the direction of conditioning. The data give
$\Prob(A \mid B_i)$ (the rate of the outcome within a class); the question asks for
$\Prob(B_i \mid A)$ (the probability of the class given the outcome).

\begin{keyfact}[Bayes' theorem]
Two events, with $\Prob(A) > 0$ and $0 < \Prob(B) < 1$:
\[
\Prob(B \mid A) \;=\; \frac{\Prob(A \mid B)\,\Prob(B)}{\Prob(A)}
\;=\; \frac{\Prob(A \mid B)\,\Prob(B)}{\Prob(A \mid B)\,\Prob(B) + \Prob(A \mid \cc{B})\,\Prob(\cc{B})}.
\]
Partition $B_1, \dots, B_n$:
\[
\Prob(B_k \mid A) \;=\; \frac{\Prob(A \mid B_k)\,\Prob(B_k)}{\displaystyle\sum_{i=1}^{n} \Prob(A \mid B_i)\,\Prob(B_i)}.
\]
In words:
\[
\text{posterior} \;=\; \frac{\text{prior} \times \text{likelihood}}{\text{sum over all classes of (prior} \times \text{likelihood)}}.
\]
Here the prior is $\Prob(B_k)$, the likelihood is $\Prob(A \mid B_k)$, and the denominator is
$\Prob(A)$ from the law of total probability.
\end{keyfact}

\textbf{Derivation.} By the definition of conditional probability,
$\Prob(B_k \mid A) = \Prob(A \cap B_k)/\Prob(A)$. The numerator is
$\Prob(A \mid B_k)\Prob(B_k)$ by the multiplication rule, and the denominator is
$\sum_i \Prob(A \mid B_i)\Prob(B_i)$ by the law of total probability. \qed

Two consequences are worth keeping in view. First, the posterior probabilities
$\Prob(B_i \mid A)$ for $i = 1, \dots, n$ sum to 1, because the numerators sum to the
denominator. Second, the ratio of two posteriors is
\[
\frac{\Prob(B_j \mid A)}{\Prob(B_k \mid A)} = \frac{\Prob(B_j)\,\Prob(A \mid B_j)}{\Prob(B_k)\,\Prob(A \mid B_k)},
\]
because the denominator cancels. The observation $A$ multiplies each class's prior odds by
that class's likelihood.

\subsection{The table method}

Every question on this outcome can be organised in one table with one row per class. The
steps:
\begin{enumerate}[nosep,leftmargin=2em]
  \item List the classes $B_i$ (they must cover everyone and not overlap).
  \item Write the prior $\Prob(B_i)$ in column 2. The column must sum to 1.
  \item Write the likelihood $\Prob(A \mid B_i)$ in column 3, where $A$ is the event
        actually observed. If the observed event is ``no accident'', column 3 holds
        $1 - (\text{accident rate})$.
  \item Multiply across: joint $\Prob(A \cap B_i) = \Prob(B_i)\,\Prob(A \mid B_i)$ in column 4.
  \item Sum column 4. The sum is $\Prob(A)$ (law of total probability).
  \item Divide each joint by the column sum: posterior $\Prob(B_i \mid A)$ in column 5
        (Bayes' theorem). Column 5 must sum to 1.
\end{enumerate}

\textbf{Worked table.} A health insurer classifies its policyholders as smokers (20\%),
former smokers (30\%), and never-smokers (50\%). The probability of a claim in a year is 0.30
for a smoker, 0.20 for a former smoker, and 0.10 for a never-smoker. A randomly chosen
policyholder files a claim. Find the probability that the policyholder is a smoker.

\begin{center}
\begin{tabular}{@{}lcccc@{}}
\toprule
Class $B_i$ & Prior $\Prob(B_i)$ & Likelihood $\Prob(A \mid B_i)$ & Joint $\Prob(A \cap B_i)$ & Posterior $\Prob(B_i \mid A)$ \\
\midrule
Smoker        & 0.20 & 0.30 & $0.20 \times 0.30 = 0.06$ & $0.06/0.17 = 0.3529$ \\
Former smoker & 0.30 & 0.20 & $0.30 \times 0.20 = 0.06$ & $0.06/0.17 = 0.3529$ \\
Never-smoker  & 0.50 & 0.10 & $0.50 \times 0.10 = 0.05$ & $0.05/0.17 = 0.2941$ \\
\midrule
Total         & 1.00 &      & $\Prob(A) = 0.17$          & 1.0000 \\
\bottomrule
\end{tabular}
\end{center}

The overall claim rate is $\Prob(A) = 0.17$, and $\Prob(\text{smoker} \mid \text{claim}) = 0.06/0.17 = 0.3529$.
Note that the prior for smokers was 0.20; observing a claim raised it to 0.35.

\textbf{Natural frequencies.} When the priors are proportions of a population, the same
computation can be done by counting in a hypothetical population. Take 1000 policyholders.
\begin{itemize}[nosep,leftmargin=2em]
  \item Smokers: $0.20 \times 1000 = 200$ people, of whom $0.30 \times 200 = 60$ file a claim.
  \item Former smokers: $0.30 \times 1000 = 300$ people, of whom $0.20 \times 300 = 60$ file a claim.
  \item Never-smokers: $0.50 \times 1000 = 500$ people, of whom $0.10 \times 500 = 50$ file a claim.
\end{itemize}
Total claimants: $60 + 60 + 50 = 170$. Of the 170 claimants, 60 are smokers, so
$\Prob(\text{smoker} \mid \text{claim}) = 60/170 = 0.3529$. The counts are the joint column
multiplied by 1000; the arithmetic is the same, but the counts are harder to get wrong under
time pressure. Use 10\,000 or 100\,000 when the rates are small (a prevalence of 0.001, say)
so that every count is a whole number.

\subsection{Worked examples}

\begin{example}[Three risk classes of drivers]
An auto insurer classifies its drivers as standard, preferred, or ultra-preferred. Of the
insurer's drivers, 60\% are standard, 30\% are preferred, and 10\% are ultra-preferred. In a
given year, the probability of an accident is 0.10 for a standard driver, 0.05 for a
preferred driver, and 0.01 for an ultra-preferred driver. A driver chosen at random has an
accident during the year. Calculate the probability that the driver is ultra-preferred.
\par\smallskip\noindent (A) 0.001\qquad (B) 0.010\qquad (C) 0.013\qquad (D) 0.076\qquad (E) 0.132
\end{example}
\begin{solution}
Let $A$ be the event of an accident. Build the table.
\begin{center}
\begin{tabular}{@{}lccc@{}}
\toprule
Class & Prior & Likelihood & Joint \\
\midrule
Standard        & 0.60 & 0.10 & 0.060 \\
Preferred       & 0.30 & 0.05 & 0.015 \\
Ultra-preferred & 0.10 & 0.01 & 0.001 \\
\midrule
Total           & 1.00 &      & $\Prob(A) = 0.076$ \\
\bottomrule
\end{tabular}
\end{center}
\[
\Prob(\text{ultra} \mid A) = \frac{0.001}{0.076} = 0.0132.
\]
Answer (C). Choice (A) is the joint probability, not the posterior; choice (D) is $\Prob(A)$.
Note the base-rate effect: ultra-preferred drivers are 10\% of the population but only 1.3\%
of the drivers who have accidents.
\end{solution}

\begin{example}[Medical test with low prevalence]
A blood test for a disease has sensitivity 0.95 (the probability of a positive result given
that the disease is present) and specificity 0.90 (the probability of a negative result given
that the disease is absent). The disease is present in 1\% of the population. A person chosen
at random tests positive. Calculate the probability that the person has the disease.
\par\smallskip\noindent (A) 0.009\qquad (B) 0.088\qquad (C) 0.095\qquad (D) 0.475\qquad (E) 0.950
\end{example}
\begin{solution}
Let $D$ be disease present and $T$ a positive test. The likelihood of a positive test in the
healthy class is the false-positive rate $1 - \text{specificity} = 0.10$.
\begin{center}
\begin{tabular}{@{}lccc@{}}
\toprule
Class & Prior & $\Prob(T \mid \cdot)$ & Joint \\
\midrule
Disease    & 0.01 & 0.95 & 0.0095 \\
No disease & 0.99 & 0.10 & 0.0990 \\
\midrule
Total      & 1.00 &      & $\Prob(T) = 0.1085$ \\
\bottomrule
\end{tabular}
\end{center}
\[
\Prob(D \mid T) = \frac{0.0095}{0.1085} = 0.0876.
\]
Answer (B). In natural frequencies with 10\,000 people: 100 have the disease and 95 of them
test positive; 9\,900 do not and 990 of them test positive; $95/(95 + 990) = 0.0876$. A test
that is 95\% sensitive and 90\% specific still leaves a positive result with under a 9\%
chance of disease, because the healthy class is 99 times larger than the diseased class and
its false positives dominate. Choice (E) confuses $\Prob(T \mid D)$ with $\Prob(D \mid T)$.
\end{solution}

\begin{example}[Claim in year two given no claim in year one]
An insurer has two classes of policyholders. Low-risk policyholders make up 70\% of the
insured population and have probability 0.10 of a claim in any given year. High-risk
policyholders make up the remaining 30\% and have probability 0.40 of a claim in any given
year. For each policyholder, claims in different years are independent given the
policyholder's class. A policyholder chosen at random has no claim in the first year.
Calculate the probability that the policyholder has a claim in the second year.
\par\smallskip\noindent (A) 0.135\qquad (B) 0.167\qquad (C) 0.190\qquad (D) 0.210\qquad (E) 0.250
\end{example}
\begin{solution}
Let $N_1$ be no claim in year 1 and $C_2$ a claim in year 2. Then
$\Prob(C_2 \mid N_1) = \Prob(N_1 \cap C_2)/\Prob(N_1)$, and both numerator and denominator
are computed by total probability over the classes. Within a class the years are independent,
so $\Prob(N_1 \cap C_2 \mid \text{class}) = \Prob(N_1 \mid \text{class})\,\Prob(C_2 \mid \text{class})$.
\begin{center}
\begin{tabular}{@{}lcccc@{}}
\toprule
Class & Prior & $\Prob(N_1 \mid \cdot)$ & Joint $\Prob(N_1 \cap \cdot)$ & $\Prob(N_1 \cap C_2 \cap \cdot)$ \\
\midrule
Low  & 0.70 & 0.90 & 0.630 & $0.630 \times 0.10 = 0.063$ \\
High & 0.30 & 0.60 & 0.180 & $0.180 \times 0.40 = 0.072$ \\
\midrule
Total & 1.00 & & $\Prob(N_1) = 0.810$ & $\Prob(N_1 \cap C_2) = 0.135$ \\
\bottomrule
\end{tabular}
\end{center}
\[
\Prob(C_2 \mid N_1) = \frac{0.135}{0.810} = 0.1667.
\]
Answer (B). An equivalent route is to update the class probabilities first:
$\Prob(\text{low} \mid N_1) = 0.630/0.810 = 7/9$ and $\Prob(\text{high} \mid N_1) = 2/9$; then
$\Prob(C_2 \mid N_1) = (7/9)(0.10) + (2/9)(0.40) = 0.1667$. Choice (C) is the unconditional
one-year claim rate $0.7(0.10) + 0.3(0.40) = 0.19$; the clean first year shifts weight toward
the low-risk class and lowers the second-year rate. The years are \emph{not} unconditionally
independent, so $\Prob(C_2 \mid N_1) \neq \Prob(C_2)$.
\end{solution}

\begin{example}[Unknown class proportion]
An insurer issues policies to two types of policyholders. Type A policyholders have
probability 0.05 of a claim in a year; type B policyholders have probability 0.20. The
insurer's overall claim rate is 0.11 per policyholder per year. A policyholder chosen at
random files a claim. Calculate the probability that the policyholder is type B.
\par\smallskip\noindent (A) 0.40\qquad (B) 0.60\qquad (C) 0.67\qquad (D) 0.73\qquad (E) 0.80
\end{example}
\begin{solution}
Let $x = \Prob(\text{A})$, so $\Prob(\text{B}) = 1 - x$. The law of total probability gives
\[
0.11 = 0.05x + 0.20(1 - x) = 0.20 - 0.15x
\quad\Longrightarrow\quad
x = \frac{0.09}{0.15} = 0.60.
\]
So 60\% are type A and 40\% are type B. Now apply Bayes' theorem; the denominator is the
given overall rate 0.11, so there is nothing further to add up:
\[
\Prob(\text{B} \mid \text{claim}) = \frac{0.40 \times 0.20}{0.11} = \frac{0.08}{0.11} = 0.7273.
\]
Answer (D). Choice (A) is the prior for type B; choice (B) is the proportion of type A.
\end{solution}

\begin{example}[Classes given as a ratio]
A company has twice as many male employees as female employees. In a given year, 6\% of male
employees and 9\% of female employees file a disability claim. An employee chosen at random
files a disability claim. Calculate the probability that the employee is female.
\par\smallskip\noindent (A) 0.333\qquad (B) 0.400\qquad (C) 0.429\qquad (D) 0.500\qquad (E) 0.600
\end{example}
\begin{solution}
``Twice as many male as female'' means the ratio male : female is $2 : 1$, so the priors are
$\Prob(M) = 2/3$ and $\Prob(F) = 1/3$. (The same information stated as odds ``$2$ to $1$
in favour of male'' converts the same way: odds $a : b$ means probabilities $a/(a+b)$ and
$b/(a+b)$.)
\begin{center}
\begin{tabular}{@{}lccc@{}}
\toprule
Class & Prior & Likelihood & Joint \\
\midrule
Male   & $2/3$ & 0.06 & 0.04 \\
Female & $1/3$ & 0.09 & 0.03 \\
\midrule
Total  & 1 & & 0.07 \\
\bottomrule
\end{tabular}
\end{center}
\[
\Prob(F \mid \text{claim}) = \frac{0.03}{0.07} = 0.4286.
\]
Answer (C). Using natural frequencies with 300 employees: 200 males with 12 claims, 100
females with 9 claims; $9/21 = 0.4286$. Choice (A) is the prior. Choice (E) comes from
using the wrong ratio (treating female as the larger class).
\end{solution}

\begin{example}[Which supplier]
A manufacturer buys a component from three suppliers. Supplier A provides 50\% of the
components, supplier B provides 30\%, and supplier C provides 20\%. The proportion of
defective components is 2\% for supplier A, 3\% for supplier B, and 5\% for supplier C. A
component chosen at random from inventory is found to be defective. Calculate the probability
that it came from supplier C.
\par\smallskip\noindent (A) 0.200\qquad (B) 0.310\qquad (C) 0.345\qquad (D) 0.400\qquad (E) 0.500
\end{example}
\begin{solution}
\begin{center}
\begin{tabular}{@{}lcccc@{}}
\toprule
Supplier & Prior & Likelihood & Joint & Posterior \\
\midrule
A & 0.50 & 0.02 & 0.010 & $0.010/0.029 = 0.345$ \\
B & 0.30 & 0.03 & 0.009 & $0.009/0.029 = 0.310$ \\
C & 0.20 & 0.05 & 0.010 & $0.010/0.029 = 0.345$ \\
\midrule
Total & 1.00 & & 0.029 & 1.000 \\
\bottomrule
\end{tabular}
\end{center}
$\Prob(\text{C} \mid \text{defective}) = 0.010/0.029 = 0.3448$. Answer (C). Choice (B) is the
posterior for supplier B, a distractor for a candidate who reads the wrong row. In natural
frequencies with 1000 components: A supplies 500 with 10 defects, B supplies 300 with 9,
C supplies 200 with 10; $10/29$.
\end{solution}

\begin{example}[Conditioning on the complement]
An insurer's drivers fall into three age groups: 30\% are young, 50\% are middle-aged, and
20\% are senior. The probability of an accident in a year is 0.20 for a young driver, 0.08
for a middle-aged driver, and 0.12 for a senior driver. A driver chosen at random has no
accident during the year. Calculate the probability that the driver is young.
\par\smallskip\noindent (A) 0.240\qquad (B) 0.274\qquad (C) 0.300\qquad (D) 0.484\qquad (E) 0.516
\end{example}
\begin{solution}
The observed event is $\cc{A}$, no accident. The likelihood in each row is
$\Prob(\cc{A} \mid \text{class}) = 1 - \Prob(A \mid \text{class})$, in every row.
\begin{center}
\begin{tabular}{@{}lccc@{}}
\toprule
Class & Prior & $\Prob(\cc{A} \mid \cdot)$ & Joint \\
\midrule
Young  & 0.30 & 0.80 & 0.240 \\
Middle & 0.50 & 0.92 & 0.460 \\
Senior & 0.20 & 0.88 & 0.176 \\
\midrule
Total  & 1.00 & & $\Prob(\cc{A}) = 0.876$ \\
\bottomrule
\end{tabular}
\end{center}
\[
\Prob(\text{young} \mid \cc{A}) = \frac{0.240}{0.876} = 0.2740.
\]
Answer (B). Choice (D), 0.484, is $\Prob(\text{young} \mid A) = 0.060/0.124$, the answer to
the opposite question; choice (E) is $1 - 0.484$, which is $\Prob(\text{not young} \mid A)$
and not $\Prob(\text{young} \mid \cc{A})$. A clean year lowers the probability of the young
class from 0.30 to 0.27, a smaller shift than an accident produces, because a clean year is
common in every class.
\end{solution}

\begin{example}[Observation spanning two years]
An insurer classifies 80\% of its policyholders as good risks and 20\% as bad risks. In each
year a good risk has probability 0.05 of at least one claim and a bad risk has probability
0.25. Given the class, claim activity in different years is independent. A policyholder
chosen at random has at least one claim during a two-year period. Calculate the probability
that the policyholder is a bad risk.
\par\smallskip\noindent (A) 0.200\qquad (B) 0.400\qquad (C) 0.529\qquad (D) 0.556\qquad (E) 0.714
\end{example}
\begin{solution}
Let $E$ be the event of at least one claim in two years. Within a class,
$\Prob(E \mid \text{class}) = 1 - (1 - p)^2$ where $p$ is the class's one-year rate, using
conditional independence.
\begin{center}
\begin{tabular}{@{}lccc@{}}
\toprule
Class & Prior & $\Prob(E \mid \cdot)$ & Joint \\
\midrule
Good & 0.80 & $1 - 0.95^2 = 0.0975$ & 0.0780 \\
Bad  & 0.20 & $1 - 0.75^2 = 0.4375$ & 0.0875 \\
\midrule
Total & 1.00 & & $\Prob(E) = 0.1655$ \\
\bottomrule
\end{tabular}
\end{center}
\[
\Prob(\text{bad} \mid E) = \frac{0.0875}{0.1655} = 0.5287.
\]
Answer (C). Choice (D), 0.556, is $\Prob(\text{bad} \mid \text{claim in a single year}) = 0.05/0.09$;
the likelihood must match the event actually observed, which here covers two years.
\end{solution}

\subsection{Traps}

\begin{trap}[Answering $\Prob(A \mid B)$ when $\Prob(B \mid A)$ is asked]
The data always give the rate of the outcome within each class, $\Prob(A \mid B_i)$. The
question almost always asks for the class given the outcome, $\Prob(B_i \mid A)$. These are
different numbers: in the medical-test example, $\Prob(T \mid D) = 0.95$ while
$\Prob(D \mid T) = 0.088$. Before computing, write down in symbols which event is observed
(goes to the right of the bar) and which is asked for (goes to the left). The answer choices
routinely include the likelihood as a distractor.
\end{trap}

\begin{trap}[Forgetting a class in the denominator]
The denominator of Bayes' theorem is $\Prob(A)$, the sum of prior times likelihood over
\emph{every} class. Omitting a class (typically the third class in a three-class problem, or
the ``no disease'' class in a test problem) makes the denominator too small and the posterior
too large. Check that the prior column sums to 1 before summing the joint column. If the
priors given in the problem do not sum to 1, the remaining mass belongs to a class the
problem describes in words (``the remaining policyholders'') and that class must be in the
table.
\end{trap}

\begin{trap}[Using the overall rate as a likelihood]
The likelihood column holds class-specific rates $\Prob(A \mid B_i)$. The overall rate
$\Prob(A)$ belongs in the total row, not in any class row. When a problem gives the overall
rate together with all but one of the class rates or priors, the overall rate is there so
that the missing quantity can be solved from the law of total probability, as in the
unknown-proportion example. Once solved, the given overall rate is the Bayes denominator; do
not recompute it from the table and get a different number through rounding.
\end{trap}

\begin{trap}[Conditional independence is given the class, not unconditional]
In multi-year problems, ``claims in different years are independent'' means independent
\emph{within each class}: $\Prob(C_1 \cap C_2 \mid B_i) = \Prob(C_1 \mid B_i)\Prob(C_2 \mid B_i)$.
It does not mean $\Prob(C_1 \cap C_2) = \Prob(C_1)\Prob(C_2)$. Because the population mixes
classes with different rates, a claim in year one raises the probability of the high-risk
class and therefore raises the probability of a claim in year two:
$\Prob(C_2 \mid C_1) > \Prob(C_2)$. A candidate who treats the years as unconditionally
independent answers with the marginal one-year rate, which is always among the choices.
Multiply the two years' probabilities inside each row of the table, never across the total
row.
\end{trap}

\begin{trap}[Conditioning on $\cc{A}$ requires $1 - \text{likelihood}$ in every row]
When the observed event is ``no accident'', ``tested negative'', or ``survived'', the
likelihood in every row is $1 - \Prob(A \mid B_i)$, and the denominator is
$\Prob(\cc{A}) = 1 - \Prob(A)$. Two wrong shortcuts appear in the choices:
using the accident rates unchanged (which gives $\Prob(B_i \mid A)$), and computing
$1 - \Prob(B_i \mid A)$ (which is $\Prob(\cc{B_i} \mid A)$, the complement on the wrong
side of the bar). Neither equals $\Prob(B_i \mid \cc{A})$. For a negative medical test the
likelihoods are $1 - \text{sensitivity}$ in the disease row and the specificity in the
healthy row.
\end{trap}

\begin{trap}[Normalising by the wrong total]
The posterior is the joint probability divided by the sum of the joint column, not by the
prior, not by the likelihood, and not by 1. Dividing the joint by the class prior returns the
likelihood; dividing by the likelihood returns the prior; leaving it undivided gives
$\Prob(A \cap B_i)$, which the choices usually include (0.001 in the driver example, 0.018 in
a test example). The posterior column must sum to 1; if it does not, the divisor is wrong.
\end{trap}

\subsection{Practice problems}

\begin{practice}
An insurer classifies its policyholders as low risk (45\%), medium risk (35\%), and high
risk (20\%). The probability of a claim in a year is 0.02 for low risk, 0.05 for medium
risk, and 0.10 for high risk. A policyholder chosen at random files a claim. Calculate the
probability that the policyholder is high risk.
\par\smallskip\noindent (A) 0.200\qquad (B) 0.387\qquad (C) 0.430\qquad (D) 0.465\qquad (E) 0.500
\end{practice}

\begin{practice}
A screening test for a condition has sensitivity 0.90 and specificity 0.95. The condition
is present in 2\% of the population screened. A person screened at random tests positive.
Calculate the probability that the person has the condition.
\par\smallskip\noindent (A) 0.018\qquad (B) 0.067\qquad (C) 0.269\qquad (D) 0.474\qquad (E) 0.900
\end{practice}

\begin{practice}
An insurer's policyholders are 60\% class A, with annual claim probability 0.10, and 40\%
class B, with annual claim probability 0.30. Given the class, claims in different years are
independent. A policyholder chosen at random has a claim in year one. Calculate the
probability that the policyholder has a claim in year two.
\par\smallskip\noindent (A) 0.180\qquad (B) 0.200\qquad (C) 0.233\qquad (D) 0.250\qquad (E) 0.300
\end{practice}

\begin{practice}
A factory produces a part on two machines. Parts from machine 1 are defective with
probability 0.02 and parts from machine 2 are defective with probability 0.05. The
proportion of all parts that are defective is 0.032. A part chosen at random is defective.
Calculate the probability that it was produced on machine 2.
\par\smallskip\noindent (A) 0.400\qquad (B) 0.500\qquad (C) 0.600\qquad (D) 0.625\qquad (E) 0.750
\end{practice}

\begin{practice}
An auto insurer has three times as many preferred drivers as standard drivers. In a year,
a standard driver has an accident with probability 0.12 and a preferred driver with
probability 0.04. A driver chosen at random has an accident. Calculate the probability that
the driver is standard.
\par\smallskip\noindent (A) 0.250\qquad (B) 0.333\qquad (C) 0.500\qquad (D) 0.667\qquad (E) 0.750
\end{practice}

\begin{practice}
Of a life insurer's policyholders, 25\% are smokers and 75\% are non-smokers. The
probability of dying within the year is 0.05 for a smoker and 0.01 for a non-smoker. A
policyholder chosen at random survives the year. Calculate the probability that the
policyholder is a smoker.
\par\smallskip\noindent (A) 0.238\qquad (B) 0.242\qquad (C) 0.250\qquad (D) 0.375\qquad (E) 0.625
\end{practice}

\begin{practice}
A distributor stocks a component from four suppliers. Supplier A provides 40\% of the
stock with a 1\% defect rate, supplier B provides 30\% with a 2\% defect rate, supplier C
provides 20\% with a 3\% defect rate, and supplier D provides 10\% with a 5\% defect rate.
A component chosen at random is defective. Calculate the probability that it came from
supplier A.
\par\smallskip\noindent (A) 0.100\qquad (B) 0.190\qquad (C) 0.286\qquad (D) 0.400\qquad (E) 0.476
\end{practice}

\begin{practice}
An insurer's policyholders are 30\% young and 70\% not young. The probability that a young
policyholder files a claim in a year is 0.15. The overall probability that a policyholder
files a claim in a year is 0.09. A policyholder chosen at random files no claim during the
year. Calculate the probability that the policyholder is young.
\par\smallskip\noindent (A) 0.064\qquad (B) 0.255\qquad (C) 0.280\qquad (D) 0.300\qquad (E) 0.500
\end{practice}

\subsection{Answers to practice problems}

\begin{enumerate}[nosep,leftmargin=2em]
  \item \textbf{(C)}. Joint probabilities $0.45(0.02) = 0.0090$, $0.35(0.05) = 0.0175$,
        $0.20(0.10) = 0.0200$; total $0.0465$. $\Prob(\text{high} \mid \text{claim}) = 0.0200/0.0465 = 0.430$.
  \item \textbf{(C)}. Joint: $0.02(0.90) = 0.018$ and $0.98(0.05) = 0.049$; total $0.067$.
        $\Prob(D \mid T) = 0.018/0.067 = 0.269$. Choice (A) is the joint, (B) is $\Prob(T)$.
  \item \textbf{(C)}. $\Prob(C_1) = 0.6(0.10) + 0.4(0.30) = 0.18$;
        $\Prob(C_1 \cap C_2) = 0.6(0.10)^2 + 0.4(0.30)^2 = 0.006 + 0.036 = 0.042$;
        $\Prob(C_2 \mid C_1) = 0.042/0.18 = 0.233$. Choice (A) is the unconditional rate.
  \item \textbf{(D)}. With $x = \Prob(\text{machine 1})$: $0.02x + 0.05(1 - x) = 0.032$ gives
        $x = 0.60$. $\Prob(\text{machine 2} \mid \text{def}) = 0.40(0.05)/0.032 = 0.020/0.032 = 0.625$.
  \item \textbf{(C)}. Priors $\Prob(\text{std}) = 1/4$, $\Prob(\text{pref}) = 3/4$. Joint:
        $0.25(0.12) = 0.03$ and $0.75(0.04) = 0.03$; total $0.06$. $\Prob(\text{std} \mid A) = 0.03/0.06 = 0.500$.
  \item \textbf{(B)}. Survival likelihoods $0.95$ and $0.99$. Joint: $0.25(0.95) = 0.2375$ and
        $0.75(0.99) = 0.7425$; total $0.98$. $\Prob(\text{smoker} \mid \text{survive}) = 0.2375/0.98 = 0.242$.
        Choice (E) is $\Prob(\text{smoker} \mid \text{die}) = 0.0125/0.02$; choice (D) is its complement.
  \item \textbf{(B)}. Joint: $0.004$, $0.006$, $0.006$, $0.005$; total $0.021$.
        $\Prob(\text{A} \mid \text{def}) = 0.004/0.021 = 0.190$. Choice (C) is the posterior for B or C.
  \item \textbf{(C)}. First solve for the not-young rate $y$: $0.3(0.15) + 0.7y = 0.09$ gives
        $y = 0.045/0.7 = 0.0643$. Then $\Prob(\text{no claim}) = 1 - 0.09 = 0.91$ and
        $\Prob(\text{young} \mid \text{no claim}) = 0.3(0.85)/0.91 = 0.255/0.91 = 0.280$.
        Choice (A) is $y$; choice (B) is the joint probability.
\end{enumerate}

\subsection{One-page summary}

\begin{itemize}[nosep,leftmargin=2em]
  \item A partition $B_1, \dots, B_n$: mutually exclusive, exhaustive, priors sum to 1.
        The two-class case is $\{B, \cc{B}\}$.
  \item Law of total probability: $\Prob(A) = \sum_i \Prob(A \mid B_i)\Prob(B_i)$.
        Two classes: $\Prob(A) = \Prob(A \mid B)\Prob(B) + \Prob(A \mid \cc{B})\Prob(\cc{B})$.
        The overall rate is the prior-weighted average of the class rates.
  \item Bayes' theorem: $\Prob(B_k \mid A) = \dfrac{\Prob(A \mid B_k)\Prob(B_k)}{\sum_i \Prob(A \mid B_i)\Prob(B_i)}$.
        Posterior $=$ prior $\times$ likelihood, divided by the sum of prior $\times$ likelihood
        over all classes. Posteriors sum to 1.
  \item Table method: (1) one row per class; (2) prior column, sums to 1; (3) likelihood
        column $\Prob(\text{observed event} \mid \text{class})$; (4) joint $=$ prior $\times$
        likelihood; (5) sum the joint column to get $\Prob(\text{observed event})$;
        (6) posterior $=$ joint $/$ column sum.
  \item Natural frequencies: take 1000 or 10\,000 people, count each class, count the
        outcome within each class, and divide the count in the target class by the total
        count with the outcome.
  \item Observed event on the right of the bar, asked-for event on the left. The data give
        $\Prob(A \mid B_i)$; the question asks $\Prob(B_i \mid A)$.
  \item Conditioning on $\cc{A}$ (no claim, negative test, survived): likelihood
        $1 - \Prob(A \mid B_i)$ in every row; denominator $1 - \Prob(A)$.
        Negative test: disease row $1 - \text{sensitivity}$, healthy row $=$ specificity.
  \item Sensitivity $= \Prob(T \mid D)$; specificity $= \Prob(\cc{T} \mid \cc{D})$;
        false-positive rate $= 1 - \text{specificity}$. Low prevalence makes $\Prob(D \mid T)$
        small even for an accurate test.
  \item Two years: independence holds within a class only. Multiply year probabilities inside
        each row: $\Prob(C_1 \cap C_2 \mid B_i) = p_i^2$, $\Prob(N_1 \cap C_2 \mid B_i) = (1 - p_i)p_i$,
        $\Prob(\text{at least one claim} \mid B_i) = 1 - (1 - p_i)^2$. Then
        $\Prob(C_2 \mid C_1) = \Prob(C_1 \cap C_2)/\Prob(C_1)$, each side by total probability.
        Equivalent: update the class posteriors after year one, then average the year-two rates.
  \item Unknown proportion: set $\Prob(B_1) = x$, write the given overall rate as
        $\sum_i \Prob(A \mid B_i)\Prob(B_i)$, solve for $x$, then use the given overall rate as
        the Bayes denominator.
  \item Ratios and odds: ``$k$ times as many in class 1 as class 2'' or odds $k : 1$ means
        priors $k/(k+1)$ and $1/(k+1)$.
  \item Distractors to recognise: the likelihood, the joint probability, the overall rate,
        the prior, the posterior of a neighbouring class, and the complement on the wrong
        side of the bar.
\end{itemize}


\end{document}
